\newif\ifanon\anonfalse

\documentclass[11pt,letterpaper]{article}
\usepackage[margin=1in]{geometry}
\usepackage[utf8]{inputenc}
\usepackage[T1]{fontenc} %
\usepackage[dvipsnames]{xcolor}
\usepackage{color} 
\usepackage{times}
\usepackage{cancel}
\usepackage{mathtools}
\usepackage{mathpazo} 


\usepackage{manfnt} %
\linespread{1.09} %

\usepackage{figchild}
\usepackage{fontawesome5}

\usepackage{accents}
\newcommand{\comb}[1]{\ensuremath{\accentset{\frown}{\text{#1}}}}


\definecolor{niceRed}{RGB}{190,38,38}
\definecolor{niceYellow}{HTML}{f5b400}
\definecolor{blueGrotto}{HTML}{059DC0}
\definecolor{royalBlue}{HTML}{057DCD}
\definecolor{navyBlue}{HTML}{0B579C}
\definecolor{yaleBlue}{HTML}{00356b}
\definecolor{limeGreen}{HTML}{81B622}
\definecolor{nicePurple}{HTML}{9c27b0}
\definecolor{lightRoyalBlue}{HTML}{def2ff}  
\definecolor{gold}{HTML}{ffa300}

\newcommand*{\white}[1]{\textcolor{white}{#1}}
\newcommand*{\red}[1]{\textcolor{red}{#1}}
\newcommand*{\blue}[1]{\textcolor{blue}{#1}}
\newcommand*{\gray}[1]{\textcolor{lightgray}{#1}}


% --- begin macros.tex ---
\DeclarePairedDelimiter{\set}{\{}{\}}
\DeclarePairedDelimiter{\iprod}{\langle}{\rangle}
\DeclarePairedDelimiter{\card}{|}{|}

\DeclareMathOperator{\conv}{conv}
\DeclareMathOperator{\Span}{span}
\DeclareMathOperator{\proj}{proj}
\DeclareMathOperator{\round}{round}
\DeclareMathOperator{\diag}{diag}

\newcommand{\up}{\uparrow}
\newcommand{\down}{\downarrow}
\newcommand{\Lattice}[2]{\mathcal{L}_{#1}^{#2}}
\newcommand{\bmat}[1]{\begin{bmatrix}#1\end{bmatrix}}

\newcommand{\sset}{\subseteq}
\newcommand{\ones}{\mathbb{1}}
\newcommand{\Cube}{\text{cube}} %
% --- end macros.tex ---

\usepackage{hyperref}
\hypersetup{
  colorlinks = true, 
  urlcolor = {blueGrotto},
  linkcolor = {royalBlue},
  citecolor = {navyBlue}
}
\usepackage[
  backref=true,
  backend=biber,
  natbib=true,
  style=alphabetic,
  sorting=alphabeticlabel,
  sortcites=true,
  minbibnames=3,
  maxbibnames=999,
  mincitenames=4,
  maxcitenames=4,
  minalphanames=4,
  maxalphanames=4,
  doi=false, url=false, eprint=false, isbn=false, 
]{biblatex}

\DeclareSortingTemplate{alphabeticlabel}{
  \sort[final]{%
    \field{labelalpha} %
  }
  \sort{%
    \field{year}   %
  }
  \sort{%
    \field{title}  %
  }
}

\AtBeginRefsection{\GenRefcontextData{sorting=ynt}}
\AtEveryCite{\localrefcontext[sorting=ynt]}
\addbibresource{refs.bib}


\usepackage[most]{tcolorbox}

\tcbset{
    mybox/.style={
        colframe=black,        %
        colback=gray!0,       %
        boxrule=0.5mm,         %
        width=\linewidth,      %
        arc=3mm,               %
        left=5mm,              %
        right=5mm,             %
        top=5mm,               %
        bottom=5mm,            %
        enhanced,              %
        halign=center        %
    }
}

\usepackage{booktabs} 
\usepackage{multicol} 
\usepackage{multirow} 
\usepackage{makecell} 
\usepackage{longtable}

\usepackage{graphicx}
\usepackage{float} 
\usepackage{tikz}
\usepackage{pgfplots}
\pgfplotsset{compat=1.17}
\usepgfplotslibrary{fillbetween}
\usetikzlibrary{arrows.meta}
\usepackage{setspace}



\tikzset{
  myNodeFlex/.style={
    draw,
    rectangle,
    rounded corners,
    text centered,
    minimum height=1.5em,
  }
}

\tikzset{
  myNode/.style={
    draw,
    rectangle,
    rounded corners,
    text centered,
    minimum height=1.5em,
    minimum width=3cm,
    text width=5cm,    
  }
}

\tikzset{
  myNodeNarrow/.style={
    draw,
    rectangle,
    rounded corners,
    text centered,
    minimum height=1.5em,
    minimum width=1cm,
  }
}

\tikzset{
  myNodeWide/.style={
    draw,
    rectangle,
    rounded corners,
    text centered,
    minimum height=1.5em,
    minimum width=6cm,
  }
}

\usepackage{xinttools} %
\usetikzlibrary{calc}

\def\biglen{20cm} %
\tikzset{
  half plane/.style={ to path={
       ($(\tikztostart)!.5!(\tikztotarget)!#1!(\tikztotarget)!\biglen!90:(\tikztotarget)$)
    -- ($(\tikztostart)!.5!(\tikztotarget)!#1!(\tikztotarget)!\biglen!-90:(\tikztotarget)$)
    -- ([turn]0,2*\biglen) -- ([turn]0,2*\biglen) -- cycle}},
  half plane/.default={1pt}
}


\usepackage{physics} %
\usepackage{amsmath}
\usepackage{amssymb}
\usepackage{amsthm}
\usepackage{bbm}
\usepackage{blkarray}
\usepackage{mathtools}
\usepackage{nicefrac}
\usepackage{dsfont}
\usepackage{xspace}
\usepackage{pifont} %

\usepackage{thm-restate} %
\usepackage[capitalise,noabbrev,nameinlink,sort]{cleveref}

\usepackage{enumerate} 
\usepackage[inline,shortlabels]{enumitem} 
\usepackage{typed-checklist}
\usepackage[framemethod=TikZ]{mdframed}
\mdfsetup{%
backgroundcolor=white, %
roundcorner=4pt,
linewidth=1pt}
\usepackage{changepage} %

\usepackage{algorithm}
\usepackage{algpseudocode}[1]
\renewcommand{\algorithmicrequire}{\textbf{Input:}}
\renewcommand{\algorithmicensure}{\textbf{Oracles:}}

\usepackage{mathrsfs} %
\usepackage{soul} %

 

\theoremstyle{plain} 
\newtheorem{theorem}{Theorem}[section]
\newtheorem{corollary}[theorem]{Corollary}
\newtheorem{proposition}[theorem]{Proposition}
\newtheorem{lemma}[theorem]{Lemma}
\newtheorem{fact}[theorem]{Fact}
\newtheorem{conjecture}[theorem]{Conjecture}
\newtheorem{claim}[theorem]{Claim}

\newtheorem{assumption}{Assumption}
\newtheorem{infassumption}{Informal Assumption}
\newtheorem{inftheorem}{Informal Theorem}
\newtheorem{minftheorem}{Main Informal Theorem}
\newtheorem{maintheorem}{Main Theorem}
\newtheorem{definition}{Definition}
\newtheorem*{definition*}{Definition}

\theoremstyle{definition} 
\newtheorem{example}[theorem]{Example}
\newtheorem{remark}[theorem]{Remark}

\theoremstyle{remark}
\newtheorem{observation}[theorem]{Observation}
\newtheorem{problem}{Problem}


\NewDocumentEnvironment{pf}{o}
  {\IfNoValueTF{#1}{\begin{proof}}{\begin{proof}[Proof of #1.]}}
  {\IfNoValueTF{#1}{\end{proof}}{\end{proof}}}

\AfterEndEnvironment{definition}{\noindent\ignorespaces}
\AfterEndEnvironment{assumption}{\noindent\ignorespaces}
\AfterEndEnvironment{lemma}{\noindent\ignorespaces}
\AfterEndEnvironment{theorem}{\noindent\ignorespaces}
\AfterEndEnvironment{proposition}{\noindent\ignorespaces}
\AfterEndEnvironment{fact}{\noindent\ignorespaces}
\AfterEndEnvironment{question}{\noindent\ignorespaces}
\AfterEndEnvironment{corollary}{\noindent\ignorespaces}
\AfterEndEnvironment{model}{\noindent\ignorespaces}
\AfterEndEnvironment{remark}{\noindent\ignorespaces}
\AfterEndEnvironment{proof}{\noindent\ignorespaces}
\AfterEndEnvironment{fact}{\noindent\ignorespaces}
\AfterEndEnvironment{minftheorem}{\noindent\ignorespaces}
\AfterEndEnvironment{inftheorem}{\noindent\ignorespaces}
\AfterEndEnvironment{maintheorem}{\noindent\ignorespaces}
\AfterEndEnvironment{restatable}{\noindent\ignorespaces}
\AfterEndEnvironment{infassumption}{\noindent\ignorespaces}
\AfterEndEnvironment{conjecture}{\noindent\ignorespaces}
\AfterEndEnvironment{claim}{\noindent\ignorespaces}

\crefname{section}{Section}{Sections}
\crefname{theorem}{Theorem}{Theorems}
\crefname{lemma}{Lemma}{Lemmas}
\crefname{definition}{Definition}{Definitions}
\crefname{conjecture}{Conjecture}{Conjectures}
\crefname{corollary}{Corollary}{Corollaries}
\crefname{construction}{Construction}{Constructions}
\crefname{conjecture}{Conjecture}{Conjectures}
\crefname{claim}{Claim}{Claims}
\crefname{observation}{Observation}{Observations}
\crefname{proposition}{Proposition}{Propositions}
\crefname{fact}{Fact}{Facts}
\crefname{question}{Question}{Questions}
\crefname{problem}{Problem}{Problems}
\crefname{remark}{Remark}{Remarks}
\crefname{example}{Example}{Examples}
\crefname{equation}{Equation}{Equations}
\crefname{appendix}{Appendix}{Appendices}
\crefname{algorithm}{Algorithm}{Algorithms}
\crefname{model}{Model}{Models}
\crefname{figure}{Figure}{Figures}
\crefname{inftheorem}{Informal Theorem}{Informal Theorems}
\crefname{infassumption}{Informal Assumption}{Informal Assumptions}
\crefname{minftheorem}{Main Informal Theorem}{Main Informal Theorems}
\crefname{maintheorem}{Main Theorem}{Main Theorems}
\crefname{assumption}{Assumption}{Assumptions}
\crefname{case}{Case}{Cases}
\crefname{program}{Program}{Programs}
\crefname{inequality}{Inequality}{Inequalities}

\newlist{asmpenum}{enumerate}{1} %
\setlist[asmpenum]{label={\arabic*.},ref=\theassumption.{\arabic*}}
\crefname{asmpenumi}{Assumption}{Assumptions}



\usepackage{etoolbox}



\newcommand{\yesnum}{\addtocounter{equation}{1}\tag{\theequation}}  

\makeatletter
\renewcommand{\eqref}[1]{\textup{\eqrefform@{\ref{#1}}}}
\let\eqrefform@\tagform@

\newcommand{\changetag}[1]{%
  \renewcommand\tagform@[1]{\maketag@@@{(\ignorespaces#1\unskip\@@italiccorr)}}%
}
\makeatother

\makeatletter
\newcommand{\tagnum}[2]{%
    \refstepcounter{equation}%
    \tag{#1) \ (\theequation}%
    \protected@write \@auxout {}{%
        \string \newlabel {#2}{{\theequation}{\thepage}{}{equation.\theequation}{}}%
    }%
}
\makeatother

\newcommand{\Stackrel}[2]{\stackrel{\mathmakebox[\widthof{\ensuremath{#2}}]{#1}}{#2}}

\newcommand{\quadtext}[1]{\quad\text{#1}\quad}
\newcommand{\qquadtext}[1]{\qquad\text{#1}\qquad}
 
\newcommand{\quadand}{\quadtext{and}}
\newcommand{\qquadand}{\qquadtext{and}}

\newcommand{\quadwhere}{\quadtext{where}}
\newcommand{\qquadwhere}{\qquadtext{where}}


\newcommand{\mat}[1]{\begin{bmatrix}#1\end{bmatrix}}
\def\abs#1{\left| #1 \right|}
\def\biggabs#1{\bigg| #1 \bigg|}
\def\sabs#1{| #1 |}
\newcommand{\given}{\;\middle|\;}
\newcommand{\sinparen}[1]{\ensuremath{(#1)}}
\newcommand{\sinbrace}[1]{\ensuremath{\{#1\}}}
\newcommand{\sinsquare}[1]{\ensuremath{[#1]}}
\newcommand{\inbrace}[1]{\ensuremath{\left\{#1\right\}}}
\newcommand{\sinangle}[1]{\ensuremath{\langle#1\rangle}}
\newcommand{\inparen}[1]{\ensuremath{\left(#1\right)}}
\newcommand{\insquare}[1]{\ensuremath{\left[#1\right]}}
\newcommand{\inangle}[1]{\left\langle#1\right\rangle}
\newcommand{\floor}[1]{\ensuremath{\left\lfloor#1\right\rfloor}}
\newcommand{\sfloor}[1]{\ensuremath{\lfloor#1\rfloor}}
\newcommand{\ceil}[1]{\ensuremath{\left\lceil#1\right\rceil}}
\newcommand{\sceil}[1]{\ensuremath{\lceil#1\rceil}}
\newcommand{\snorm}[1]{\ensuremath{\ensuremath{\| #1 \|}}}
\let\norm\relax
\newcommand{\norm}[1]{\ensuremath{\left\lVert #1 \right\rVert}}

\newcommand{\midsepremove}{\aboverulesep = 0mm \belowrulesep = 0mm}
\newcommand{\midsepadd}{\aboverulesep = 1mm \belowrulesep = 1mm}
\newcommand{\midsepdefault}{\aboverulesep = 0.605mm \belowrulesep = 0.984mm}


\newcommand{\N}{\mathbb{N}}
\newcommand{\R}{\mathbb{R}}
\newcommand{\Q}{\mathbb{Q}}
\newcommand{\Z}{\mathbb{Z}}

\newcommand{\evA}{\ensuremath{\mathscr{A}}}
\newcommand{\evB}{\ensuremath{\mathscr{B}}}
\newcommand{\evC}{\ensuremath{\mathscr{C}}}
\newcommand{\evD}{\ensuremath{\mathscr{D}}}
\newcommand{\evE}{\ensuremath{\mathscr{E}}}
\newcommand{\evF}{\ensuremath{\mathscr{F}}}
\newcommand{\evG}{\ensuremath{\mathscr{G}}}
\newcommand{\evH}{\ensuremath{\mathscr{H}}}
\newcommand{\evI}{\ensuremath{\mathscr{I}}}
\newcommand{\evJ}{\ensuremath{\mathscr{J}}}
\newcommand{\evK}{\ensuremath{\mathscr{K}}}
\newcommand{\evL}{\ensuremath{\mathscr{L}}}
\newcommand{\evM}{\ensuremath{\mathscr{M}}}
\newcommand{\evN}{\ensuremath{\mathscr{N}}}

\newcommand{\E}{\operatornamewithlimits{\mathbb{E}}} 
\newcommand{\Ex}{\E}
\newcommand{\Exp}{\Ex}
\newcommand{\Cov}{\ensuremath{\operatornamewithlimits{\rm Cov}}}
\newcommand{\cov}{\ensuremath{\operatornamewithlimits{\rm Cov}}}
\let\var\relax
\newcommand{\var}{\ensuremath{\operatornamewithlimits{\rm Var}}}
\newcommand{\Var}{\var}
\newcommand\ind{\mathds{1}}
\newcommand{\sgn}{\operatornamewithlimits{sgn}}
\newcommand{\argmin}{\operatornamewithlimits{arg\,min}}
\newcommand{\argmax}{\operatornamewithlimits{arg\,max}}

\newcommand{\tv}[2]{\operatorname{d}_{\mathsf{TV}}{\inparen{#1,#2}}}
\newcommand{\kl}[2]{\operatornamewithlimits{\mathsf{KL}}{\inparen{#1\|#2}}} 
\newcommand{\chidiv}[2]{\chi^2\inparen{#1\|#2}}
\newcommand{\srenyi}[3]{\cR_{#1}\sinparen{#2\|#3}}
\newcommand{\renyi}[3]{\cR_{#1}\inparen{#2\|#3}}

\newcommand{\zo}{\ensuremath{\inbrace{0, 1}}}
\newcommand{\oo}{\ensuremath{\inbrace{-1, 1}}}
\newcommand{\IntsTo}[1]{\setb{1, 2, \dotsc, #1}}
\newcommand{\IntsTon}{\IntsTo{n}}
\newcommand{\sfrac}[2]{{#1/#2}} 
\newcommand{\nfrac}[2]{\nicefrac{#1}{#2}}
\newcommand{\OPT}{\mathrm{OPT}} 
\newcommand{\opt}{\textrm{OPT}}
\newcommand{\Err}{\mathrm{Err}}
\newcommand{\err}{\textrm{\rm Err}}
\newcommand{\diam}{\mathrm{diam}}
\newcommand{\poly}{\mathrm{poly}}
\newcommand{\polylog}{\mathrm{polylog}}
\newcommand{\vc}{\textrm{\rm VC}}
\newcommand{\vol}{\textrm{\rm vol}}
\newcommand{\supp}{\operatorname{supp}}
\newcommand{\cdf}{\ensuremath{{\rm cdf}}}
\newcommand{\iid}{i.i.d.}
\newcommand{\unif}{{\rm Unif}}
\newcommand{\nll}{\mathrm{Null}}

\newcommand{\cupdot}{\charfusion[\mathbin]{\cup}{\cdot}}
\newcommand{\bigcupdot}{\charfusion[\mathop]{\bigcup}{\cdot}}
\newcommand*\circled[1]{\tikz[baseline=(char.base)]{
            \node[shape=circle,draw,inner sep=1pt] (char) {#1};}}

\newcommand{\tick}{\ding{51}}
\let\cross\relax
\newcommand{\cross}{\ding{55}}

\newcommand{\eps}{\varepsilon}
\renewcommand{\epsilon}{\varepsilon}
\makeatletter
\newcommand*{\tran}{{\mathpalette\@tran{}}}
\newcommand*{\@tran}[2]{\raisebox{\depth}{$\m@th#1\intercal$}}
\makeatother

\mathchardef\NABLA"272
\newcommand*{\Nabla}{\boldsymbol\NABLA}
\let\nabla\Nabla

\renewcommand{\hat}{\widehat}
\newcommand{\wh}[1]{\widehat{#1}}
\newcommand{\hf}{\wh{f}}

\newcommand{\hD}{\wh{D}}
\newcommand{\hF}{\wh{F}}
\newcommand{\hR}{\wh{R}}
\newcommand{\hS}{\wh{S}}
\newcommand{\hW}{\wh{W}}

\newcommand{\sS}{S^\star}

\renewcommand{\bar}{\overline}
\newcommand{\oE}{\overline{E}}
\newcommand{\oS}{\overline{S}}

\renewcommand{\tilde}{\widetilde}
\newcommand{\wt}[1]{\widetilde{#1}}
\newcommand{\tS}{\wt{S}}
\newcommand{\tT}{\wt{T}}

\newcommand{\customcal}[1]{\euscr{#1}}
\newcommand{\cA}{\customcal{A}}
\newcommand{\cB}{\customcal{B}}
\newcommand{\cC}{\customcal{C}}
\newcommand{\cD}{\customcal{D}}
\newcommand{\cE}{\customcal{E}}
\newcommand{\cF}{\customcal{F}} 
\newcommand{\cG}{\customcal{G}} 
\newcommand{\cH}{\customcal{H}} 
\newcommand{\cI}{\customcal{I}} 
\newcommand{\cJ}{\customcal{J}} 
\newcommand{\cK}{\customcal{K}}
\newcommand{\cL}{\customcal{L}}
\newcommand{\cM}{\customcal{M}}
\newcommand{\cN}{\customcal{N}}
\newcommand{\cO}{\customcal{O}}
\newcommand{\cP}{\customcal{P}} 
\newcommand{\cQ}{\customcal{Q}} 
\newcommand{\cR}{\customcal{R}} 
\newcommand{\cS}{\customcal{S}}
\newcommand{\cT}{\customcal{T}} 
\newcommand{\cU}{\customcal{U}} 
\newcommand{\cV}{\customcal{V}} 
\newcommand{\cW}{\customcal{W}}
\newcommand{\cX}{\customcal{X}}
\newcommand{\cY}{\customcal{Y}}
\newcommand{\cZ}{\customcal{Z}}


\DeclareMathAlphabet{\mathdutchcal}{U}{dutchcal}{m}{n}
\SetMathAlphabet{\mathdutchcal}{bold}{U}{dutchcal}{b}{n}
\DeclareMathAlphabet{\mathdutchbcal}{U}{dutchcal}{b}{n}


\DeclareMathAlphabet\urwscr{U}{urwchancal}{b}{n}%
\DeclareMathAlphabet\rsfscr{U}{rsfso}{m}{n}
\DeclareMathAlphabet\euscr{U}{eus}{m}{n}
\DeclareFontEncoding{LS2}{}{}
\DeclareFontSubstitution{LS2}{stix}{m}{n}
\DeclareMathAlphabet\stixcal{LS2}{stixcal}{m} {n}

\newcommand{\toa}[1]{\textit{#1}} %
\renewcommand{\paragraph}[1]{\medskip \noindent\textbf{#1}~}
\newcommand{\itparagraph}[1]{\medskip \noindent\textit{#1}~}
\newcommand{\paragraphit}[1]{\itparagraph{#1}}

\newcommand{\eg}{\emph{e.g.}}
\newcommand{\ie}{\emph{i.e.}}

\usepackage{upgreek}
\renewcommand{\gamma}{\upgamma}
\renewcommand{\pi}{\uppi}


\newcommand\D{\mathcal{D}}
\newcommand\C{\mathcal{C}}

\renewcommand{\S}{\mathcal{S}} 
\newcommand{\T}{\mathbb{T}} 
\newcommand{\W}{\mathcal{W}}
\newcommand{\X}{\mathcal{X}} 
\newcommand{\Y}{\mathcal{Y}}  
\newcommand{\pred}{\mathrm{pred}}
\newcommand{\eat}[1]{}
\newcommand{\negLL}{\ensuremath{\mathscr{L}}}



\newcommand{\hypo}[1]{\mathdutchcal{#1}}
\newcommand{\hyA}{\hypo{A}}
\newcommand{\hyB}{\hypo{B}}
\newcommand{\hyC}{\hypo{C}}
\newcommand{\hyD}{\hypo{D}}
\newcommand{\hyE}{\hypo{E}}
\newcommand{\hyF}{\hypo{F}}
\newcommand{\hyG}{\hypo{G}}
\newcommand{\hyH}{\hypo{H}}
\newcommand{\hyI}{\hypo{I}}
\newcommand{\hyJ}{\hypo{J}}
\newcommand{\hyK}{\hypo{K}}
\newcommand{\hyL}{\hypo{L}}
\newcommand{\hyM}{\hypo{M}}
\newcommand{\hyN}{\hypo{N}}
\newcommand{\hyO}{\hypo{O}}
\newcommand{\hyP}{\hypo{P}}
\newcommand{\hyQ}{\hypo{Q}}
\newcommand{\hyR}{\hypo{R}}
\newcommand{\hyS}{\hypo{S}}
\newcommand{\hyT}{\hypo{T}}
\newcommand{\hyU}{\hypo{U}}
\newcommand{\hyV}{\hypo{V}}
\newcommand{\hyW}{\hypo{W}}
\newcommand{\hyX}{\hypo{X}}
\newcommand{\hyY}{\hypo{Y}}
\newcommand{\hyZ}{\hypo{Z}} 

\renewcommand{\d}{{\rm d}}

\newcommand{\normal}[2]{\cN{\inparen{#1,#2}}}
\newcommand{\coarseNormal}[3]{\cN_{#3}\inparen{#1,#2}}
\newcommand{\truncatedNormal}[3]{\cN\inparen{#1,#2,#3}}
\newcommand{\normalMass}[3]{\cN\inparen{#1, #2; #3}}

\newcommand{\wstar}{{w^\star}}
\newcommand{\Wstar}{{W^\star}}
\newcommand{\fstar}{{f^\star}}
\newcommand{\ystar}{{y^\star}}
\newcommand{\xstar}{{x^\star}}
\newcommand{\muStar}{{\mu^\star}}
\newcommand{\SigmaStar}{{\Sigma^\star}}

\newcommand{\maxcut}{\textsf{MAXCUT}\xspace}
\newcommand{\promisemaxcut}{\textsf{PromiseMAXCUT}\xspace}
\newcommand{\uniqueSAT}{\textsf{Unique3SAT}\xspace}
\newcommand{\SAT}{\textsf{3SAT}\xspace}
\newcommand{\localLL}{\negLL_{\rm Local}}                 

\usepackage{caption}
\usepackage{subcaption}
\usepackage{graphicx}



\def\n{23} %
\def\maxxy{4} %

\newcommand{\Pclass}{\ensuremath{\mathsf{P}}\xspace}
\newcommand{\np}{\ensuremath{\mathsf{NP}}\xspace}
\newcommand{\NP}{\ensuremath{\mathsf{NP}}\xspace}
\newcommand{\RP}{\ensuremath{\mathsf{RP}}\xspace}


\newcommand{\gmruntime}{\ensuremath{\poly(d,{\eps^{-k}}) + \inparen{\nfrac{\log\inparen{k}}{\eps}}^{O(k)}}}

\newcommand{\ymax}{y_{\max}}
\newcommand{\ynoisymax}{y_{\text{noisy-max}}}


\newcommand{\customcdf}[2]{\Phi\sinparen{#1; \sigma_{#2}^2}}
\newcommand{\proofconstant}{{3}}
\renewcommand{\cL}{\negLL}
\renewcommand{\hyL}{\negLL}


\newcommand{\smax}{\max,2} %
\newcommand{\imax}{i_{\max}}
\newcommand{\ysmax}{y_{\smax}}

\newcommand\blfootnote[1]{%
  \begingroup
  \renewcommand\thefootnote{}\footnote{#1}%
  \addtocounter{footnote}{-1}%
  \endgroup
}

\usepackage{array}
\newcolumntype{L}[1]{>{\raggedright\let\newline\\\arraybackslash\hspace{0pt}}m{#1}}
\newcolumntype{C}[1]{>{\centering\let\newline\\\arraybackslash\hspace{0pt}}m{#1}}
\newcolumntype{R}[1]{>{\raggedleft\let\newline\\\arraybackslash\hspace{0pt}}m{#1}}

\usepackage{fancyhdr}

\fancypagestyle{custom}{
  \fancyhf{}                %
  \fancyfoot[C]{\textcolor{black}{\fcFishD{0.025}{black}{0.5}}} %
  \renewcommand{\headrulewidth}{0pt}
  \renewcommand{\footrulewidth}{0pt}
}

\newcommand{\negsep}{\vspace{0mm}}

\title{
    Can SGD Select Good Fishermen?\\
    Local Convergence under Self-Selection Biases and Beyond
} 
%
%
%
%
%                   ███                                                           ██
%                ██    ██                                                       ██  ██
%            █████      ████                                                ████      █████
%           ██              █                                              █              ████
%           ██             ██                                             ██                 █
%             █████████████                                                 ██████████████████
%
%
%      █            ███████                  █████████          ██        ████████           ██
%       ██         ██     ██                 █      ██         ███       ██     ██          ███
%        ██      █████████████            ██████████████      ████     █████████████       ████
%         ██      ██       ██               ██       ██      ██ ██       █       ██       ██ ██
%           ██    █   █  █ ██               █   ██  █ █    ██   ██      ██   █ ██ █     ██   ██
%            ██   ██  █  ████               ██       ██    █    ██      ██       ██    ██    ██
%             ██   ██     ██                 ██     ██   ██     ██       ███    ██    ██     ██
%              ██    █████                     █████    ██      ██         █████     ██      ██
%               ██    ███                       ███    ██       ██          ██      ██       ██
%                ██  █████                    ██████  ██        ██         ████    ██        ██
%                 ████ ████                 ██  ██ ████         ██         █████ ██          ██
%                  ██  ██  ██             ██    ██  ██          ██        ████ ███           ██
%                 ████ ██   ██            ████  ██ ████         ██       ██ ██  ██           ██
%                 ████████████                ██████████        ██       █████████           ██
%                      ██                       ███             ██          ████             ██
%                      ███                      ██              ██          ███              ██
%                      ██                       ███             ██          ██               ██
%                     ███                       ███             ██         ████              ██
%                     █ ██                     ██ █             ██         █  █              ██
%         █          ██  ██                   ██   █            ██        ██  ██             ██
%      ██████       ██    █      ███ ██       █    ██      ███  ██  ███   █    ██      ███   ██
%   ██████████████  ██    ██████████████████ ██     █ ████ ███████ █████ ██     ███████ ███████ ██
%                  ██      ██                ██     ██          ██       █      ██           ██
%                  █       ██                █       █          ██      ██       ██          ██
%               ████        ████          ████        ██        ██      ███                  ██
%                                                               ██                           ██
%  ████████████████████████████████████████████████████████████████████████████████████████████████
%                                                               ██                           ██
%              ██     ████████                                  ██                           ██
%               ██████      █ ███                               ██        ██████████  ████   ██
%               ██      █████ █ ██    ██    ██████████                  ████ █       █ ██    ██
%              ████ ████   █████      ██████    █  █████                ███  █     ███████   ██
%                       █████         ██ ██     ████   ██                  ████████     █    ██
%                                     ███  ████    ████                                      ██
%      ██      █████████                        ███            ████████    ██                ██
%      ████████      █ ███                                  ███ █     ███████
%       ██        ████    ██                               ██ █ ██       █ ██
%      █████████     █████                                  █████  █████  ████
%               ███████                                         ████
%
%
%
%
%
%
%
\ifanon
\author{Anonymous Author(s)}
\else
\author{
        \begin{tabular}{C{5.1cm}C{5.1cm}C{5.1cm}}
        {\bf Alkis Kalavasis} 
            & {\bf Anay Mehrotra} 
                & {\bf Felix Zhou}\\[2mm]
        {Yale University} 
            & {Yale University} 
                & {Yale University}\\[1mm]
        \mbox{\small\texttt{\href{mailto:alkis.kalavasis@yale.edu}{alkis.kalavasis@yale.edu}}} 
            & \mbox{\small\texttt{\href{mailto:anaymehrotra1@gmail.com}{anaymehrotra1@gmail.com}}}
                & \mbox{\small\texttt{\href{felix.zhou@yale.edu}{felix.zhou@yale.edu}}}
        \\
        \end{tabular}
}
\fi



\date{}

\begin{document}

\maketitle
\thispagestyle{empty}

\begin{abstract}
    We revisit the problem of estimating $k$ linear regressors with self-selection bias in $d$ dimensions with the maximum selection criterion, as introduced by \citet[STOC'23]{cherapanamjeri2023selfselection}. 
    Our main result  
    is a $\poly(d,k,\nfrac{1}{\eps}) + {k}^{O(k)}$ time algorithm for this problem, 
    which yields an improvement in the running time of the algorithms of \cite{cherapanamjeri2023selfselection} and \citet[arXiv]{gaitonde2024selfselection}. We achieve this by providing the first local convergence algorithm for self-selection, thus resolving the main open question of \cite{cherapanamjeri2023selfselection}.

    To obtain this algorithm, we 
    reduce self-selection to a seemingly unrelated statistical problem called coarsening. Coarsening occurs when one does not observe the exact value of the sample but only some
    set (a subset of the sample space) that contains the exact value. 
    Inference from coarse samples arises in various real-world applications due to rounding by humans and algorithms, limited precision of instruments, and lag in multi-agent systems.  

    Our reduction to coarsening is intuitive and relies on the geometry of the self-selection problem, which enables us to bypass the limitations of previous analytic approaches. To demonstrate its applicability, we provide a local convergence algorithm for linear regression under another self-selection criterion, which is related to second-price auction data. 
    Further, we give the first polynomial time local convergence algorithm for coarse Gaussian mean estimation given samples generated from a convex partition. Previously, only a sample-efficient algorithm was known 
    due to \citet[COLT'21]{fotakis2021coarse}.
\end{abstract}

\newpage 
\thispagestyle{empty}
\setcounter{tocdepth}{3}
{   
    \setstretch{0.9}
    \tableofcontents
}
\thispagestyle{empty}
\newpage

\pagenumbering{arabic}

\negsep{}

\section{Introduction}
    Self-selection bias occurs when data is systematically selected rather than randomly sampled.
    Inference under selection biases has a rich history in Statistics and Econometrics, starting with the foundational works of \citet{roy1951earnings,willis1979education,heckman1979sample}.
    {This framework has since been applied in} various scientific fields, from causal inference and imitation learning \cite{heckman1990varieties}, to learning from strategically reported data \cite{hardt2016strategic,dong2018strategic,kleinberg2020classifiers}, and to learning from auction data \cite{athey2002identification,athey2007nonparametric,cherapanamjeri2022auctionEstimation}.
    
    A concrete application of self-selection bias appears in the work of \citet{fair1972methods}, who studied estimation in markets at disequilibrium, where supply does not match demand.
    Fair and Jaffee modeled the housing market, with $y_S(x)$ representing supply and $y_D(x)$ {denoting} demand, as {functions} of {features} $x$ encoding, \eg{}, location, size, and amenities. 
    For a given $x$, {only the transaction price} $(x,\min\{y_S(x), y_D(x)\})$ {is observed} instead of the complete sample $(x,y_S(x),y_D(x))$, 
    leaving it unclear whether the market imbalance stems from excess supply or demand.

    Another classic model due to \citet{roy1951earnings} examines workers' occupational choices based on potential earnings.
    Suppose there are $k$ occupations (\eg{}, hunting, fishing, and woodcutting) and, for each occupation $i$, the expected income for workers with feature vector $x$ is given by $y_i(x)=x^\top w_i^\star + \xi_i$, 
    where $w_i^\star$ is a parameter vector and $\xi_i \sim \normal{0}{1}$ {captures} the sum of independent market events related to the $i$-th occupation.
    {Workers} select their occupations \textit{strategically} to maximize revenue, considering all potential incomes $y_1(x), y_2(x),\dots, y_k(x)$, and selecting the occupation that offers the highest income.
    From the analyst's perspective, however,
    only the feature vector $x$ and the maximum income achieved $\max \{y_1(x), y_2(x), \dots, y_k(x)\}$ are observed.
    The goal is {then} to estimate the unknown parameters $w^\star_1,w^\star_2,\dots,w^\star_k$ {which determine the expected income from becoming a hunter, fisherman, woodcutter, and so on.} 
    
    
    These examples, as many others, are cases of self-selection bias with the \textit{maximum} selection criteria.\footnote{{The minimum in Fair-Jaffee's model can be converted to a maximum by negating the parameter vectors and noise.}}  
    Since the seminal works of \citet{roy1951earnings,fair1972methods}, 
    estimation under self-selection bias has received significant attention from studies of participation in the labor force \cite{heckman1974shadow,hanoch1976multivariate,NELSON1977309,heckman1979sample,Hanoch1980,Cogan1980}, 
    to migration and income \cite{nakosteen1980migration,borjas1987selfSelection}, 
    to effects of unions on wages \cite{lee1978union,abowd1982job}, 
    to returns to education \cite{willis1979education,kenny1979college,griliches1978missing}, 
    and to choices between tenure choice and demand for housing \cite{lee1978estimation,king1980tenure,rosen1979housing}; see \cite{maddala1986limited} {for more in-depth discussions and references}. 
    {Despite this extensive history, efficient algorithms for estimation under self-selection -- even for natural selectors like \( \min \) and \( \max \) -- were not known until \citet{cherapanamjeri2023selfselection} initiated the design of sample- and computationally-efficient linear regression under self-selection biases.}
    Let us first give a more formal exposition of the model they studied: linear regression with maximum self-selection bias. 
    \begin{definition}
    [Max Self-Selection \cite{cherapanamjeri2023selfselection}]
    In linear regression with self-selection bias under the maximum selection criterion, a sample $(x,\ymax) \in \R^d \times \R$ is generated as follows: 
    \begin{enumerate}
        \item A covariate $x$ is drawn from $\normal{0}{I}$, and 
        \item $\ymax$ is the maximum of $k$ different unknown linear functions of $x$ that are independently perturbed by noise, \ie{}, $\ymax = \max_{i \in [k]} \set{x^\top w_i^\star + \xi_i}$, where $w_1^\star,\dots,w_k^\star$ are the unknown target parameters and $\xi_1,\dots,\xi_k$ are independent $\normal{0}{1}$ random variables.
    \end{enumerate}
    The unknown parameters $w_1^\star,\dots,w_k^\star$ satisfy the following conditions: %
    \begin{enumerate}
        \item[(a)] (Separability)~~  $ \norm{w^\star_i}_2^2 
            \geq \Omega(1) + \max_{j\neq i}~\sabs{\sinangle{w^\star_j, w^\star_i}}.$
        \item[(b)] (Boundedness) ~~$\max_{1\leq i\leq k}\norm{w^\star_i}_2\leq O(1)\,.$
    \end{enumerate}
    \label{def:ssb}
    \end{definition}
    \noindent Both conditions on the parameters are standard and also appear in prior works  \cite{cherapanamjeri2023selfselection,gaitonde2024selfselection}.
    Without condition (a), it is information-theoretically impossible to separate the weight vectors (see Remark 6.1 in \cite{gaitonde2024selfselection}), and condition (b) is a standard boundedness assumption.
    
    This model captures the aforementioned classical works of \citet{roy1951earnings} and \citet{fair1972methods}.
    Remarkably, \cite{cherapanamjeri2023selfselection} presented an algorithm to estimate regression parameters $w^\star_1,\dots,w^\star_k$ to within $\poly(\nfrac{1}{k})$-error in $\poly(d)\cdot e^{\poly(k)}$ time under \Cref{def:ssb}.
    This result is surprising since even the \emph{identifiability} of the parameters under \Cref{def:ssb} is highly non-trivial.
    Subsequently, the follow-up work by \citet{gaitonde2024selfselection} significantly improved the sample complexity {to the near optimal:} they recover the regression parameters to error $\eps$ with $O(d) \cdot \poly(k, \nfrac{1}{\eps})$ samples in {$\poly(d,k,\nfrac{1}{\eps}) + O(\nfrac{(\log{k})}{\eps})^{O(k)}$} time.
    {Both algorithms first identify a \(k\)-dimensional subspace containing \(w_1^\star, \dots, w_k^\star\) by computing the span of the top eigenvectors of some weighted covariance matrix. They then perform a brute-force search within this subspace to identify a regressor \(w_i\) (and then iterate to find the remaining regressors). Unfortunately, this brute-force step necessarily introduces an \(\eps^{-k}\) dependence in the running time. 
        }
    
    Our main result is a faster algorithm for the above model {that avoids the $\eps^{-k}$ dependence}.  
    
    \begin{inftheorem}[Efficient Estimation under Self-Selection; see \cref{thm:SS}]
        \label{infthm:Intro1}
        \label{infthm:Intro:SS}
        Under \Cref{def:ssb}, there is an algorithm for linear regression with self-selection bias under the maximum selection criterion that recovers the weights $w^\star_1,w^\star_2,\dots,w^\star_k$ up to $\eps$-error in $\poly(d,k,\nfrac{1}{\eps}) + {k}^{O(k)}$ time using $O(d) \cdot \poly(k,\nfrac{1}{\eps})$ samples.\footnote{{As in previous work, we recover the weights up to permutation, which is information-theoretically optimal.}}
    \end{inftheorem}
    This is the first $\poly(d,\nfrac{1}{\eps}) \cdot k^{O(k)}$ time algorithm for the self-selection problem under maximum selection criteria (henceforth, the \textit{self-selection problem}).
    {It} is based on the first local convergence algorithm for self-selection, which, in particular, resolves an open problem posed by \cite{cherapanamjeri2023selfselection}.
    \begin{inftheorem}[Local Convergence under Self-Selection; see \Cref{thm:SS:localConvergence}]
\label{infthm:Intro:SS:localConvergence}
        Under \Cref{def:ssb}, 
        there is an algorithm for linear regression with self-selection bias that, given a $\poly\inparen{\nfrac{1}{k}}$-warm start, 
        recovers weights $w_1^\star,w_2^\star,\dots,w_k^\star$ up to $\eps$-error with $O\inparen{\nfrac{d^2}{\eps^2}}\cdot \poly(k)$ time and $O\inparen{\nfrac{d}{\eps^2}}\cdot \poly(k)$ samples.
    \end{inftheorem} 
        Our approach to prove \cref{infthm:Intro:SS,infthm:Intro:SS:localConvergence},  
        is geometric, rooted in a reduction to a classical statistical problem known as \textit{coarsening} \cite{heitjan1990inference} and complements the important analytic techniques of prior works \cite{cherapanamjeri2023selfselection,gaitonde2024selfselection} 
        by offering a new perspective to the self-selection problem.
        Our idea is extensively described in \Cref{section:tool}.
        We believe that our framework will find applications in other problems with exogenous and endogenous selection biases. 
        In this direction, we show how to apply our tools to other self-selection criteria (see \Cref{sec:apps}). 

\negsep{}

\subsection{Coarsening}
\label{section:tool}
Our main conceptual contribution is a reduction of the self-selection problem to a seemingly unrelated statistical problem known as \emph{coarsening.}
In this section,  we introduce learning from coarse data in its simplest form and the tools we need for the reduction. 
This groundwork is important for \Cref{sec:reduction}, where we show how to {cast the self-selection problem as inference from coarse data and then} employ this framework to obtain faster algorithms for the self-selection problem.

Coarsening occurs when the exact value of a sample is not observed; instead, only a subset of the sample space containing the exact value is known.
Coarse data naturally arises in diverse fields, including Economics, Engineering, Medical and Biological Sciences, and all areas of the Physical Sciences 
(\eg{}, \cite{heitjan1990inference,heitjan1991ignorability,heitjan1993ignorability,gostic2020practical,leung1997censoring}). 
One of the simplest forms of coarsening is rounding, where data values are mapped to the nearest point on a specified lattice.  

The problem of estimation from coarse data, in its simplest form, can be described as follows. Consider the family of normal distributions $\{\normal{\mu}{I}\}_{\mu}$  over the $d$-dimensional Euclidean space $\R^d$ with identity covariance matrix. Suppose {that} a  (potentially unknown) partition $\hyP$ {divides} $\R^d$ into sets {and our goal is to estimate the unknown parameter $\muStar$ of the target distribution $\normal{\muStar}{I}$.}
In the coarse learning problem, the  algorithm has access to an oracle that operates as follows:
\begin{center}
    {On each query, the oracle samples $x\sim \normal{\mu^\star}{I}$ and returns the unique set $P$ {in} $\hyP$ containing $x.$}
\end{center}
We denote the distribution on sets $P\in \hyP$ as $\coarseNormal{\mu^\star}{I}{\hyP}.$
Crucially, the {actual} point $x$ itself is not observed.
{The challenge is to design an algorithm that, given \iid{} (set) samples \(P_1, \dots, P_n \sim \coarseNormal{\mu^\star}{I}{\hyP}\), accurately estimates \(\mu^\star\).}
{We focus on algorithms with time- and sample-complexity polynomial in the dimension and the desired accuracy, achieving an \(L_2\)-norm approximation of the true mean.}
On one extreme, when $\hyP$ consists of singletons, 
{the problem reduces to} ``fine'' Gaussian mean estimation problem where simply computing the empirical average {suffices}.
On the other extreme, when $\hyP = \{\R^d\}$, \ie{}, the only set of the partition is the whole space, parameter recovery {becomes} information-theoretically impossible. 
For further examples, see \Cref{fig:convex-partitions}.
\begin{figure}[htbp]
    \centering
    \begin{subfigure}[t]{0.3\linewidth}
        \centering
        \includegraphics[width=\linewidth]{fig/convex_partition_nonidentifiable_shaded.pdf}
    \end{subfigure}
    \hfill
    \begin{subfigure}[t]{0.3\linewidth}
        \centering
        \includegraphics[width=\linewidth]{fig/convex_partition_polytope_shaded.pdf}
    \end{subfigure}
    \hfill
    \begin{subfigure}[t]{0.3\linewidth}
        \centering
        \includegraphics[width=\linewidth]{fig/convex_partition_known_index_shaded.pdf}
    \end{subfigure}
    \caption{
        {The figure illustrates different partitions $\hyP$ of $\R^2$.}
        The figure on the left corresponds to a partition that is not identifiable: Given any $\mu^\star \in \R^2$, the vector $\mu_t = \mu^\star + t e_1$ {(for any $t\in \R$)} induces the same coarse Gaussian distribution and so $\coarseNormal{\mu^\star}{I}{\hyP}$ and $\coarseNormal{\mu^\star_t}{I}{\hyP}$ are identical.
        The middle and right figures are identifiable and correspond to convex partitions of the space.
    }
    \label{fig:convex-partitions}
\end{figure}


    
    

\negsep{}

\subsubsection{Coarse Negative Log-Likelihood Objective}
A classical approach to estimating parameters from coarse samples is to use the negative log-likelihood objective.
This idea is used in the foundational works on learning from coarsened data \cite{heitjan1990inference,heitjan1991ignorability} and has also appeared in recent work in Theoretical Computer Science \cite{fotakis2021coarse}. 
Specifically, to estimate $\mu^\star$ from coarse samples, one can consider the population negative log-likelihood objective (NLL)
\begin{equation}
      \cL(\mu) = -\E_{P \sim \coarseNormal{\mu^\star}{I}{\hyP}} \log \normalMass{\mu}{I}{P}\,. 
      \label{eq:NLL-intro}
\end{equation}
  Namely, we will see that identifiability, sample complexity, and computational efficiency all depend on the structure of the sets in the partition $\hyP$.
  This dependence on $\hyP$ was also observed by \citet{fotakis2021coarse},
  who discovered that (identifiable) convex partitions are statistically ``easy,''
  while even simple non-convex sets can induce computationally intractable instances.
  Our exploration further refines the boundary of tractability by showing that some non-convex partitions, such as those resulting from self-selection (\Cref{sec:reduction,sec:apps}), can still be computationally tractable.

\negsep{}

\subsubsection{Ingredient I: Information Preservation}
\label{sec:tool1}

As we have seen, there are simple cases of coarse partitions where inferring the true mean is information-theoretically impossible (left partition in \cref{fig:convex-partitions}). 
To determine when consistent estimation under coarsening is possible, we must assess the extent to which the coarsening mechanism (which we identify with the partition \(\hyP\)) \emph{distorts} the fine space.
To this end, we revisit the notion of \emph{information preservation} of a partition $\hyP$, 
introduced by \cite{fotakis2021coarse},
and provide a more general formulation.

\begin{definition}
[\textbf{Information Preservation}]
\label{def:informationPreservation}
            Given constants $\alpha\in [0,1]$, 
            true parameters $(\muStar,\SigmaStar)$ of a $d$-dimensional Gaussian distribution $\normal{\muStar}{\SigmaStar}$, 
            and a partition $\hyP$ of the domain $\R^d$, 
            $\hyP$ is said to be \emph{$\alpha$-information preserving 
            at radius $R>0$
            with respect to $\normal{\muStar}{\SigmaStar}$} if, 
            for any parameters $(\mu,\Sigma)$ 
            {that are $R$-close to $(\muStar,\SigmaStar)$ in $\ell_2$-norm},\footnote{
            For a matrix $A$, $\norm{A}_F$ denotes its Frobenious norm and is defined as $\norm{A}_F^2=\sum_i\sum_j A_{ij}^2$.}
            it holds that
            \[
                \tv{
                    \coarseNormal{\mu}{\Sigma}{\hyP}
                }{  
                    \coarseNormal{\muStar}{\SigmaStar}{\hyP}
                }
                \geq 
                \min\inbrace{
                    1,
                    \alpha \inparen{
                        \norm{\mu - \muStar}_2^2 
                        +
                        \norm{\Sigma - \SigmaStar}_F^2
                    }^{1/2}
                }
                \,.\footnote{ Taking a minimum on the RHS is somewhat redundant since the TV distance is upper bounded by $1$. However, this serves as a useful notational reminder of this restriction when carrying out calculations.}
            \]
        \end{definition}
        Information preservation is parameterized by $\alpha \in [0,1]$ (and $R$).
        Intuitively, large values of $\alpha$ correspond to partitions that preserve a lot of information about the original (non-coarsened) probability measures; 
        {while as} $\alpha \to 0$, {the coarse measures become nearly indistinguishable in total variation, making the problem statistically impossible}.
        Note that information preservation is defined with respect to the true parameters $(\mu^\star,\SigmaStar)$, and {it} can vary significantly with different true parameters.
        {Moreover, the radius parameter \(R\) imposes a local structure by allowing information to be preserved only in a ball of radius \(R\) around the true parameters.}
      {This locality is essential; for example, in problems with permutation invariance} on the true parameters {(\eg{}, the self-selection problem of \Cref{def:ssb} or mixture models)},
      as any permutation of the true parameters 
      would make the LHS of \Cref{def:informationPreservation} zero, 
      but the RHS would be positive.
        
        Our starting observation is that information preservation {implies ``local growth''} for the negative log-likelihood ({due to} Pinsker's inequality and standard connections between the KL divergence and MLE).
       
    \begin{lemma}
    [Information Preservation Implies Quadratic Growth]
    \label{thm:quadGrowth}
    Consider a partition $\hyP$ of $\R^d$ that is $\alpha$-information preserving at radius $R$ with respect to $\normal{\mu^\star}{\SigmaStar}$. Then,
            \begin{equation}
                 \hyL(\mu, \Sigma)
            -
            \hyL(\mu^\star, \Sigma^\star) 
            \geq
            \min\inbrace{
                2,
                2\alpha^2\inparen{\|\mu - \mu^\star\|_2^2 + \|\Sigma - \Sigma^\star\|_{F}^2}
            }\,.
            \label{eq:QuadGrowth}
            \end{equation}
        {for $(\mu, \Sigma)$ that are $R$-close to $(\mu^\star, \SigmaStar)$ in $\ell_2$-norm.}
    \end{lemma}
    Thus, if $\alpha$-information preservation holds, then $\sinparen{\muStar,\SigmaStar}$ is the \textit{unique} minimizer of the negative log-likelihood {\textit{locally}}.
    {However, minimizing the negative log-likelihood can be challenging:}
    Since the distribution of coarse samples does not belong to an exponential family, the negative log-likelihood {is, in general,} non-convex.
    Hence, while the quadratic growth {property} is reminiscent of strong convexity, it does not imply global {or even local} convexity of the negative log-likelihood; in fact, quadratic growth {may hold} even when the NLL is highly non-convex
    (see \Cref{fig:quadratic-growth}).


    \begin{figure}
        \centering
        \includegraphics[width=0.3\linewidth]{fig/quadratic_growth_bounded.pdf}
        \caption{This figure illustrates a function $f$ {that} is {not convex}
        but {does} satisfy a quadratic local growth condition.
        {Indeed, $f$ is lower bounded by a quadratic function (shown by the dotted line).}
        }
        \vspace{-4mm}
        \label{fig:quadratic-growth}
    \end{figure}
    

\subsubsection{Ingredient II: Local Convexity of NLL}
\label{sec:tool2}

{For the remainder of this section, assume} the partition $\hyP$ is $\alpha$-information preserving with respect to $(\muStar, \SigmaStar)$. The key idea to designing efficient algorithms lies in understanding the \emph{optimization landscape} induced by the coarse likelihood objective of \eqref{eq:NLL-intro}. As shown in \Cref{fig:quadratic-growth}, this objective is, in general, {non-convex} and contains undesirable stationary points.

In such settings, one can hope for a \emph{local convergence} algorithm -- one that converges to the optimum given a sufficiently good warm start.
A natural algorithm of this kind is stochastic gradient descent (SGD) performed on \eqref{eq:NLL-intro}.  
However, proving theoretical guarantees for {SGD} is highly non-trivial because it relies on showing that the negative log-likelihood objective of \eqref{eq:NLL-intro} is convex {in a sufficiently large neighborhood (of radius \(R\)) around the optimum}, {which is generally not guaranteed in the presence of undesirable stationary points}. 
{Further, while one might be able to establish convexity for very small radii (\eg{}, on the order of \(\poly(\nfrac{1}{d})\)),} to get useful and interesting algorithmic guarantees (such as in \Cref{infthm:Intro:SS}), one needs $R$ to be ``non-trivially'' large (\eg{}, have no dependence on $\nfrac{1}{d}$).
 
\negsep{}

\subsubsection{A Recipe for Local Convergence}\label{sec:localConvReceipe}
    Assuming the ingredients in the previous two sections, we get the following local convergence algorithm for inference from a coarse partition.
   In particular, if we can demonstrate that the partition is information preserving and that the negative log-likelihood (NLL) is locally convex, then SGD (with a warm start for which both properties hold) will converge to the true parameters, given access to stochastic gradients of the NLL objective.
    \begin{mdframed}
        \textbf{A Recipe for Local Convergence.}\quad Suppose the following hold.
        \vspace{-2mm}
        \begin{enumerate}[itemsep=-1.5pt,leftmargin=12pt]
            \item Partition $\hyP$ is $\alpha$-information preserving with respect to $\normal{\mu^\star}{\SigmaStar}$ in an $\ell_2$-ball of radius $R_{1}.$
            \item The NLL is convex in an $\ell_2$-ball of radius $R_2$ centered at $(\mu^\star, \SigmaStar)$.
            \item There is an oracle that outputs stochastic gradients $g$ of NLL; with $\Ex[\norm{g}_2^2]\leq G^2$. 
        \end{enumerate}
        Then stochastic gradient descent on the NLL initialized within radius $R  = \min\{R_1,R_2\}$ of $\inparen{\muStar,\SigmaStar}$
        and run for $O\inparen{\nfrac{G^2}{(\alpha^4\eps^2)}}$ steps
        computes an estimate $\sinparen{\wh{\mu},\wh{\Sigma}}$ such that 
        \[
            \norm{\wh{\mu} - \muStar}_2 + \snorm{\wh{\Sigma} - \SigmaStar}_F \leq \eps\,.
        \]
    \end{mdframed}

\negsep{}

\subsection{Speeding-Up Self Selection Algorithms}
\label{sec:reduction}   
{Having discussed the ingredients for local convergence in the coarsening problem}, we are ready to explain how these ingredients can be realized for the self-selection problem. The ingredients of \Cref{sec:tool1,sec:tool2} suggest the following program to obtain an algorithm for self-selection. 
Later in this section, we explain how to implement each one of these steps.

\begin{tcolorbox}[mybox,breakable]
    \vspace{-2mm}
\begin{description}[itemsep=1pt]
    \item[Step 1 {(From Self-Selection to Coarse-Inference):}]~
        We show that linear regression under self-selection with the maximum criterion (\Cref{def:ssb}) can be encoded as a coarse-inference task.
        However, even when seen through the lens of coarsening, the NLL (see \eqref{eq:NLL-intro}) is \emph{non-convex} under the induced partition.
    

    \item[Step 2 {(Information Preservation for Self-Selection):}]
        We lower bound the \emph{information preservation} of the self-selection coarsening mechanism (\Cref{thm:InfoPresSelfSelection}). 
        This holds in an $O(\nfrac{1}{\log k})$ radius around the optimal
        {and} directly gives us a quadratic growth condition for the NLL, thanks to \Cref{thm:quadGrowth}. 
    \item[Step 3 {(Local Convexity for Self-Selection):}] 
        We show that the resulting NLL satisfies \emph{local convexity} around the optimal parameters (\Cref{thm:LocalConvSelfSelection}). 
        The radius of the local convexity is $\poly(\nfrac{1}{k})$ -- which, perhaps surprisingly, is independent of $d$ and is the key for the runtime in \cref{infthm:Intro:SS:localConvergence}. 
       
    \item[Step 4 {(Local Convergence Algorithm):}]
        Steps 2 and 3 give us the local convergence theorem (\Cref{infthm:Intro:SS:localConvergence}), which resolves an open problem of \cite{cherapanamjeri2023selfselection}.
    \item[Step 5 {(End-to-End Algorithm):}]
        Finally, using the algorithm of \cite{gaitonde2024selfselection} as a warm start, we obtain an end-to-end algorithm that improves on \cite{cherapanamjeri2023selfselection,gaitonde2024selfselection} and gives us \Cref{infthm:Intro:SS}.
\end{description}
    \vspace{-4mm}
\end{tcolorbox}

\smallskip

\noindent Next, we show how to implement each one of these steps.

\paragraph{Step 1 (From Self-Selection to Coarse-Inference).}
Recall that in self-selection, given the feature vector $x \in \R^d,$ we observe the maximum of $k$ variables $y_1,\dots,y_k$ determined by linear regression models $y_i = x^\top w_i^\star + \xi_i$. 
We view the observation $(x,\ymax)$ with $\ymax = \max\{y_1,\cdots,y_k\}$ as a \emph{coarse label} on the $k$-dimensional Euclidean space $\R^k$, where the partition consists of the following sets 
\begin{equation}
   P_{y} = \inbrace{ 
    (y_1,\dots,y_k)\colon \exists i \in [k], ~~~~\text{such that},~~~~ y_i = y 
        \quadand\max_{j \neq i} y_{j} \leq y
    }\,,\quad y \in \R \,.
\label{eq:L-shapeSets} 
\end{equation}
Hence, the learner observes the coarse example $(x,P_{\ymax})$, where $P_{\ymax}$ replaces the fine label $(y_1,\dots,y_k)$. 
Observe that $\{P_y\}_{y \in \R}$ partitions the space $\R^k$ and the information obtained by the tuple $(x,\ymax)$ is exactly the same as the information of $(x, P_{\ymax}).$ For $k=2$, the partition is depicted in \Cref{fig:self-selection-partitions}.

\begin{figure}[htbp]
    \centering
    \hspace{5em}
    \begin{subfigure}[t]{0.3\linewidth}
        \centering
        \includegraphics[width=\linewidth]{fig/self_selection_partitions_thicc.pdf}
    \end{subfigure}
    \hfill
    \begin{subfigure}[t]{0.3\linewidth}
        \centering
        \includegraphics[width=\linewidth]{fig/self_selection_one_observation_thicc.pdf}
    \end{subfigure}
    \hspace{5em}
    \caption{The left figure is an approximate illustration of the partition over $\R^2$ that the self-selection mechanism (\cref{def:ssb}) induces
    over the dependent variable space for $k=2$. Each set $P_{y_{\max}}$ corresponds to some green $L$-shape set.
    The true partition covers the entire space with the $2$-dimensional $L$-shapes.
    The right figure is an example of an observation from the self-selection model in the dependent variable.
    {See \cref{fig:self-selection-3D} for an illustration of the partitions with $k=3$.}
    }
    \label{fig:self-selection-partitions}
\end{figure}

\paragraph{Step 2 (Information Preservation for Self-Selection).} The above reduction gives rise to a partition of the $k$-dimensional Euclidean space. A natural first question is whether one can compute the information preservation of the self-selection coarsening partition with respect to the true parameters $W^\star = [w_1^\star, \dots, w_k^\star] \in \R^{d \times k}.$ 

\begin{inftheorem}
[Information Preservation for Self-Selection; see \Cref{thm:SS:localConvexity}]
Consider the model of \Cref{def:ssb}  with true parameters $W^\star = [w_1^\star,w_2^\star,\dots,w_k^\star]$. For any matrix $W\in \R^{d\times k}$ that is $O\inparen{\nfrac{1}{\log{k}}}$-close to $\Wstar$ in Frobenius norm, it holds that
            \[
                \tv{\cM(W^\star)}{~\cM(W)} \geq 
                    \poly\inparen{\nfrac{1}{k}}
                    \cdot 
                        {\norm{W^\star - W}_F}
                        \,,
            \]
where $\cM(W)$ (respectively $\cM(W^\star))$ is the distribution of {the observations} ${(x,\max_i y_i)}$ induced by $W$ (respectively $W^\star)$. 
\label{thm:InfoPresSelfSelection}
\end{inftheorem}
\noindent Here, obtaining a $\poly(\nfrac{1}{k})$ lower bound on the information preservation parameter
-- which is independent of $d$ -- is crucial in avoiding exponential-in-$d$ dependence in the running time of the local convergence algorithm {and enables us to obtain the optimal 
$O(\nfrac{d}{\eps^2})$ sample complexity and {a corresponding} $O(\nfrac{d^2}{\eps^2})$ running time} (\Cref{infthm:Intro:SS:localConvergence}).
We postpone the technical details of this result to the upcoming Technical Overview section (\cref{sec:Overview}) and present the full proof in \cref{sec:SS:proof:informationPreservation}.

\paragraph{Step 3 (Local Convexity for Self-Selection).} Having shown information preservation with respect to the true parameters, we obtain that the negative log-likelihood enjoys quadratic growth around $w_1^\star,\dots,w_k^\star.$ However, this says nothing about whether the optimization landscape is favorable for a gradient-based method (\eg{}, it could be highly non-convex). Our second contribution is a local convexity guarantee for the NLL objective. 


\begin{inftheorem}
[Local Convexity for Self-Selection; see \Cref{thm:SS:localConvexity}]
Consider the model of \Cref{def:ssb} with true parameters $W^\star = [w_1^\star,w_2^\star,\dots,w_k^\star]$.
For any matrix $W\in \R^{d\times k}$ that is $\poly(\nfrac{1}{k})$-close to $\Wstar$ in Frobenius norm, it holds that
                    \[
                        \grad^2 \cL(W) \succeq 0\,,
                    \]
where $\cL(\cdot)$ is the population negative log-likelihood objective of the self-selection problem evaluated at $W.$
\label{thm:LocalConvSelfSelection}
\end{inftheorem}
This result shows that the radius where the landscape is convex is of order $\poly(\nfrac{1}{k})$ {which} has no dependence on $d$.
As for information preservation, we postpone the technical details of this result to the upcoming Technical Overview section (\cref{sec:Overview}) and present the full proof in \cref{sec:SS:proof:localConvexity}.

\paragraph{Step 4 (Local Convergence Algorithm).}
Having shown
\Cref{thm:InfoPresSelfSelection,thm:LocalConvSelfSelection}, we are ready to obtain \Cref{infthm:Intro:SS:localConvergence}. To do that, we have to employ the recipe in \cref{sec:localConvReceipe}, which requires some careful analysis of the SGD algorithm, which we discuss further in the {Technical Overview} section (\Cref{sec:Overview}).

\paragraph{Step 5 (End-to-End Algorithm).} To obtain the improved end-to-end self-selection algorithm of \Cref{infthm:Intro:SS}, we use the Gaitonde--Mossel algorithm \cite{gaitonde2024selfselection} as a warm start, setting $\eps = \poly(\nfrac{1}{k}).$ Then we apply our local convergence method of \Cref{infthm:Intro:SS:localConvergence}. In total, we get an algorithm that uses $O(d)\cdot  \poly(k,\nfrac{1}{\eps})$ samples and $\poly(d,k,\nfrac{1}{\eps}) + k^{O(k)}$ runtime.

\subsection{Technical Overview for the Self-Selection Algorithm}
\label{sec:Overview}
    In this section, we present the main ideas and challenges for obtaining the local convergence algorithm in \Cref{infthm:Intro:SS:localConvergence} for the self-selection problem (\cref{def:ssb}). To this end, we will explain how to ensure information preservation for self-selection (\cref{thm:InfoPresSelfSelection}) and local convexity (\cref{thm:LocalConvSelfSelection}).

\subsubsection*{{Conversion to a Coarse Labels Problem}} 
    Our main conceptual idea is to convert the self-selection problem (\cref{def:ssb}) to a coarse labels problem, where the learner observes samples $(x,P_{\ymax})$.
    Here, $P_{\ymax} \subseteq\R^k$ corresponds to the ``$L$-shape sets'' of \Cref{fig:self-selection-partitions} and $\ymax \in \R$ is the observed maximum of the self-selection model given feature $x$, (recall \eqref{eq:L-shapeSets}). 
    For this conversion to be useful, the resulting partition $\{P_{\ymax}\}_{\ymax \in \R}$ of $\R^k$ must be $\alpha$-information preserving with respect to the true parameters $W^\star$ for some non-trivially large $\alpha$ (and within {a} sufficiently large radius {$R$}). Additionally, {(local) convexity must hold within this radius as well}.
    In particular, to get a $\poly(d,\nfrac{1}{\eps})$-time algorithm, the parameter $\alpha$ must be independent of $d$ and to obtain a $\poly(d,k,\nfrac{1}{\eps})$-time local convergence algorithm, {it} must be at least $\poly(\nfrac{1}{k})$.

    \subsubsection*{Establishing Information Preservation} 
    \smallskip 

         \noindent First, we discuss several challenges in lower-bounding $\alpha$, and then we present our approach. %
    
    
    \paragraph{{Challenge I: ``Variance Reduction'' Fails.}} 
            If the sets of the partition were convex, then the situation would be much easier. 
            For instance, \Cref{fig:convex-partitions} (right) corresponds to a convex partition where each 
            $P_{\ymax}$ is replaced by two sets $P_{\ymax}^1$ and $P_{\ymax}^2$ for $\ymax\in \R$. 
            {With this partition, the estimation} problem reduces to the \emph{known-index} self-selection task, where the learning algorithm observes the maximum value but {crucially} also the \textit{index} of the model realizing {the maximum}. 
            In this setting, \cite{cherapanamjeri2023selfselection} showed that negative log-likelihood is {{(globally)} \textit{strongly}} convex and the convexity of the sets was crucial in this proof. 
            At the core, this uses the important result that conditioning a Gaussian to a convex set always reduces its variance.
            Since the sets in our task are non-convex, we cannot use this variance reduction property to deduce strong convexity (which would have implied information preservation). 
        
        
    \paragraph{{Challenge II: No Notion of Distance to $P_{\ymax}$.}} %
            {A natural approach to lower bounding $\alpha$ is to use \cref{def:informationPreservation}. This requires showing that}
            if a parameter $W$ is $\eps$-far from $\Wstar$,
            then the distributions $\cM(W)$ and $\cM(\Wstar)$ of the observed maximum $\ymax$ with parameters $W$ and $\Wstar$ respectively are also $\eps$-far from each other (in total variation).
            To show this, {we need} to show that $W$ assigns a higher (or lower) mass to sufficiently many sets $\inbrace{P_{m}}_m$ compared to the mass $\Wstar$ assigns.
            Fix a particular feature $x$ and a value $m\in \R$.
            If the distance to the set $P_{m}$ was well-defined and $W^\top x$ was further from $P_m$ than $( W^\star)^\top x$, then it would be easy to show that $W$ places a lower mass on $P_m$ than $\Wstar$ (and vice versa).
            This follows because the values $y_1,y_2,\dots,y_k$ in \cref{def:ssb} are concentrated around $(W')^\top x$ for each parameter $W'$.
            However, since {the set} $P_m$ is non-convex, we cannot use this argument for a general parameter $W$ close to $\Wstar$.
            The only situation in which this change-in-mass is clear is when $W^\top x$ is a translation of $(W^\star)^\top x$ along the direction of the axis-of-symmetry of the $L$-shapes, \ie{}, along the vector $1_k = (1,1,\dots,1)$
            (\Cref{fig:information-preservation-self-selection} (left)).
            

        
            \begin{figure}
                \centering
                \hspace{5em}
                \begin{subfigure}[t]{0.3\linewidth}
                    \centering
                    \includegraphics[width=\linewidth]{fig/information_preservation_self_selection_italy.pdf}
                \end{subfigure}
                \hfill
                \begin{subfigure}[t]{0.3\linewidth}
                    \centering
                    \includegraphics[width=\linewidth]{fig/information_preservation_far_italy.pdf}
                \end{subfigure}
                \hspace{5em}
                \caption{
                    The left figure is an illustration of moving $W^\star$ to $W \in \R^{d \times k}$. In general, it is unclear how the assigned mass on the set $P_{y_{\max}}$ changes. The depicted direction of change is along an 'easy' direction where the new point $W^\top x \in \R^{k \times 1}$ assigns less mass than $( W^\star)^\top x.$ This is because any point $p$ on the $L$-shape, is further from $(W^\star)^\top x$ than from $W^\top x$. 
                    The right figure gives an example where the $L$-shape behaves like a convex set since the 
                    random variable $x^\top W^\star$ is ''deep inside'' the $L$-shape, conditioned on the good event $\evE$ and this effectively enables us to ignore all but one part of the $L$-shape.
                }
                \label{fig:information-preservation-self-selection}
                \label{fig:information-preservation-far}
            \end{figure} 


        \paragraph{{Challenge III: Existing Techniques are Insufficient.}}
        While we cannot directly use a notion of distance to the sets $P_m$ (as a notion of distance does not exist), we may still be able to show information preservation via a more complex argument. 
            One idea is to use the analysis in prior works \cite{cherapanamjeri2023selfselection,gaitonde2024selfselection}.
            \cite{cherapanamjeri2023selfselection}, roughly speaking, study the behavior of high-degree moments of $\ymax$ conditioned on certain $e^{-\poly(1/\eps)}$-probability events.
            However, conditioned on an event $\evE$, the largest lower bound on $\alpha$ one can deduce is $\Pr[\evE]$ (\cref{lem:SS:IP:reductionToConditionalIP}) and, hence, \cite{cherapanamjeri2023selfselection}'s approach can only provide a weak lower bound of $\alpha \geq e^{-\poly(1/\eps)}$.
            This is insufficient to obtain a $\poly(d,\nfrac{1}{\eps})$-time algorithm.
            \cite{gaitonde2024selfselection}'s techniques are more promising: they are able to analyze the low-degree moments of $\ymax$ conditioned on an event $\evE$, which happens with a nontrivially large probability, $\Pr[\evE]\geq \poly(\nfrac{1}{k})$.
            Hence, in principle, one might hope to prove a lower bound of $\poly(\nfrac{1}{k})$ conditioned on $\evE$.
            To lower bound $\alpha$, one has to show that the total variation distance between distributions $\cM(W)$ and $\cM(V)$ induced by candidate regressors $W$ and $V$ scales \textit{linearly} with $\norm{W-V}_F$.
            While the moments of $\cM(W)$ and $\cM(V)$ (conditioned on $\evE$), along with other properties of the model, might be sufficient to give some lower bound on $\tv{\cM(W)}{~\cM(V)}$,
            it does not lead to the linear dependence on $\norm{W-V}_F$ necessary for information-preservation. 
            We need significantly tighter analysis to establish this linear dependence.
            
    \paragraph{Our Approach to Show Information Preservation.}
    To obtain the $\poly(\nfrac{1}{k})$-information preservation, 
    we draw inspiration from the events considered in \cite{gaitonde2024selfselection} to define certain, slightly different, ''good'' events $\evE_i$ for $i\in [k]$
    over the randomness in $x$ that happens with constant probability (\cref{def:SS:event}).
    Subsequently, we follow a fundamentally different analysis to obtain a sharp linear dependence
    \[
        \tv{\cM(\Wstar)}{~\cM(W)}=\Omega(\norm{\Wstar-W}_F)\,.
    \]
    {The key to obtaining the above linear lower bound is the novel observation that conditioned} on the event $\evE_i$, 
    the $L$-shape $P_y$ {(for certain values $y$)} ``behaves like the convex set'' of the known-index variant of the self-selection problem, which makes the analysis much more tractable.
  {The core intuition} is the following:
    First, $\evE_i$ implies that, for any matrix $A \in \{W^\star, W\}$, 
    the value of the $i$-th model $y_{i, A}$ is far from any other value $y_{j, A}$, $j \neq i$ (\cref{lem:SS:propertiesOfRho}). 
    This means that, conditioned on $\evE_i$, the mass of the measure of $(y_{1,A},\dots,y_{k,A})$ is concentrated ``far'' from the main diagonal of $\R^k$ (where $y_{i,A} \approx y_{j,A}$ for any $j$), thus making the $L$-shape {$P_y$ (for $y\approx y_{i, A}$)}  look like a convex set (see \cref{fig:information-preservation-far}).
    
    {If this was true for all sets $P_y$, then one could hope to reduce our analysis to the analysis of the known index of selection problem studied in \cite{cherapanamjeri2023selfselection}.
    However, because this only holds for sets $P_y$ with  $y\approx y_{i, A}$ we need to do a more careful analysis.}
    {Toward this,} we use the fact that total variation corresponds to the supremum over all events in the $\sigma$-algebra. 
    Hence, {after some calculations,} we get that it suffices to prove that (see \cref{eq:SS:probabilityDiff})
    \[
       \abs{
                \Pr\insquare{y_{\max, W^\star} \leq 0\mid \evE_i}
                -
                \Pr\insquare{y_{\max, W} \leq 0\mid \evE_i}
            }
            \geq 
            \frac{1}{\Pr[\evE_i]}
            \cdot \poly{(\nfrac{1}{k})}
            \cdot {\norm{\Wstar-W}_F}\,.
    \] 
    {Here $y_{\max, W^\star}$ and $y_{\max, W}$ denote the random variable $\ymax$ when the parameter in the self-selection model is $W^\star$ and $W$ respectively.}
    Now, {we use the} structure imposed by the event $\evE_i$ {along with careful} approximations of the {ratios of the} Gaussian cumulative distribution function {(\cref{lem:SS:cdfratios}) to establish the linear dependence in the above inequality}.
        

       
    
\subsubsection*{{Establishing} Local Convexity}  Next, we turn to establishing local convexity.
A calculation of the Hessian of the negative log-likelihood for the self-selection problem gives that
\[
\nabla^2 \cL(W) = I_{dk} - \Ex_{x,\ymax}\inparen{
            \cov_{z\sim \normal{W^\top x}{I}}\insquare{z\mid z\in P_{\ymax}}
            \otimes xx^\top}\,,
            \yesnum\label{eq:convexity:intro:1}
\]
where $(x,y_{\max})$ has {distribution} induced by \Cref{def:ssb} and the distribution of the covariance matrix is the normal distribution conditioned on the $L$-shape $P_{y_{\max}}$. 

\paragraph{{Challenges in Establishing Local Convexity.}}
To argue about the (local) convexity of $\cL(\cdot)$, we would like to show that certain covariance {matrices} are spectrally upper bounded by the identity (see \eqref{eq:convexity:intro:1}). 
A natural idea is to use classical variance reduction inequalities:
The Brascamp-Lieb inequality states that the variance of the Gaussian distribution reduces when conditioning on a convex set \cite{harge2004convex}. Unfortunately,
the set $P_{y_{\max}}$ is non-convex,
{making} such inequalities {inapplicable}. 
{Another} idea is to use our ''good'' event $\evE_i$ used before for information preservation in order to reduce the non-convex set $P_{y_{\max}}$ to some convex set. However, this idea {also} falls short: the event $\evE_i$ has only constant mass and hence, with the remaining constant probability,  variance reduction will fail to hold {(this was not an issue for establishing information preservation as the contributions from the part of the event-space where $\evE_i$ does not hold is non-negative)}. %

\paragraph{Our Approach to Show Local Convexity.}
Instead, we follow a more algebraic route. First, thanks to information preservation, we know that $\cL(W)$ is strongly convex at the true optimal parameters $W^\star;$ also we know that $\cL$ has quadratic growth (\Cref{thm:quadGrowth}).
We hence directly bound the 
change in Hessian value, \ie{}, upper bound the difference $\nabla^2 \cL(W) - \nabla^2 \cL(W^\star)$ for parameters $W$ close to $W^\star$. 
It is indeed not difficult to show that this change is at most of order $\poly(d).$ 
However, this bound gives no algorithmic guarantees. A much more delicate analysis is required to obtain the desired $\poly(k)$ bound, which allows us to get the $\poly(\nfrac{1}{k})$-warm-start of \Cref{thm:LocalConvSelfSelection}. 

To be more concrete, first, we show that there is an event $\evF$ that happens with probability $1-e^{-\poly(k)}$ (\cref{def:SS:LC:event}), conditioned on which we can perform a change of measure between the covariances in $\cL(W)$ and $\cL(\Wstar)$ to bound the spectral norm of $\nabla^2 \cL(W) - \nabla^2 \cL(W^\star)$ by $\poly\inparen{\nfrac{1}{k}}$ (\cref{lem:SS:LC:changeOfMeasure}).
Next, it remains to bound $\snorm{\nabla^2 \cL(W) - \nabla^2 \cL(W^\star)}_2$ when $\evF$ does not occur.
This is tricky as conditioning on the complement of $\evF$ skews the distributions of all involved random variables.
To overcome this, we use the fact that $\Pr[\lnot\evF]=e^{-\poly(k)}$, to show that it is sufficient to bound certain moments of $zz^\top \otimes xx^\top$ {where $z\sim \normal{W^\star x}{I_k}$} (\textit{without} conditioning on any event) by $\poly(k)$.
This bound follows by, roughly, observing that only $k$ ``directions'' of $x$ are correlated with $z$ (those along the vectors $w^\star_1,w^\star_2,\dots,w^\star_k$) 
and that $\Ex[xx^\top]=I$ which has a constant spectral norm (\cref{lem:SS:LC:step2}).


\subsubsection*{{Completing the Analysis of the Local-Convergence} Algorithm of \Cref{infthm:Intro:SS:localConvergence}} The above two results are almost all we need to obtain the local convergence algorithm of \Cref{infthm:Intro:SS:localConvergence}. The algorithm is stochastic gradient descent on the negative log-likelihood objective given a $\poly(\nfrac{1}{k})$-warm-start. To complete the analysis of the algorithm, we need to get a bound on the second moment of the norm of the gradient (see \Cref{sec:SS:secondMoment}) and unbiased estimates of the gradients (which reduces to sampling from Gaussian distribution conditional on the sets $P_{y_{\max}}$). Both of these steps are standard and appear in \Cref{sec:SS}.





    


\subsection{Applications to Other Self-Selection Mechanisms}
\label{sec:apps}
    We believe that the connection we draw between models with self-selection biases and coarse-inference tasks can have further applications to other problems with limited-dependent variables \cite{maddala1986limited}. 
    To {illustrate} the applicability of our tools (namely, information preservation and local convexity of {the log-}likelihood), we {demonstrate} an application to estimation {from second-price auction data}.

Estimating bid distributions from observed {auction sequences} is a fundamental {primitive} in Econometrics with many practical applications. 
{The main challenge in this task is that observed information 
(e.g., the winner and the price) 
is strategically selected} 
(see \cite{cherapanamjeri2022auctionEstimation} for further discussion 
{of the challenges and applications}). 
\begin{restatable}[]{definition}{defSecondPrice}
    \label{def:secondPrice}
    In linear regression from second-price auction data, a sample $(x,\imax,\ysmax) \in \R^d \times [k] \times \R$ is generated as follows: 
    \begin{enumerate}
        \item $x$ is drawn from $\normal{0}{I}$,
        \item $\imax$ is the index of the \emph{winner} when bids are generated by unknown linear functions of $x$ that are independently perturbed by noise, \ie{}, $\imax = \argmax_{i \in [k]} x^\top w_i^\star + \xi_i$, where $w_1^\star,\dots,w_k^\star$ are the unknown target parameters and $\xi_1,\dots,\xi_k$ are independent $\normal{0}{1}$ random variables, and 
        \item ${\ysmax}$ is the \emph{second-highest bid}, \ie{}, ${\ysmax} = \max_{j \neq i} x^\top w_j^\star + \xi_j$.
    \end{enumerate}
    The unknown parameters $w_1^\star,\dots,w_k^\star$ satisfy the same conditions as in \cref{def:ssb}.
\end{restatable} 
 \cref{fig:2nd-price-observation} illustrates the coarsening arising from this model.
{This illustration shows that the partition arising from the second-price auction model is, in general, non-convex. When there are at least $k\geq 3$ participants in the auction, the partition is always non-convex. In the special case with $k=2$ participants,
the partition is convex, which makes the estimation task simpler.}
We give the following local convergence guarantee for this model {for all $k$}.
\begin{inftheorem}[see \cref{thm:2ndPrice:localConvergence}]
Under \Cref{def:secondPrice}, there is an algorithm for linear regression from second-price auction data that given {a} $\poly(\nfrac{1}{k})$ warm-start, recovers the weights $w_1^\star, \dots , w_k^\star$ up to $\eps$-error 
in $O\inparen{\nfrac{d^2}{\eps^2}}\cdot \poly(k)$ time and using $O\inparen{\nfrac{d}{\eps^2}}\cdot \poly(k)$ samples.
\end{inftheorem}
The formal statement of this result appears in \cref{sec:2ndPrice},
and its full proof appears in \cref{sec:2ndPrice:proof:informationPreservation}.
In fact, this result can be straightforwardly extended to $\ell$-th price auctions,
which capture realistic self-selection biases arising in Reliability Theory.
Concretely, it arises in cases where some property of a complex system (\eg{}, the functionality of a machine) depends on properties of individual components (\eg{}, the lifetime of each component of the machine). The statistical task of interest is to infer properties of each component of a complex system, 
such as its lifetime \cite{meilijson1981autopsy}, 
from observations of the corresponding properties of the complex system (also see \cref{rem:reliability}). The selection function is controlled by the structure of the components in the system; see, \eg{}, 
the works of \citet{gertsbakh1988statistical,huang2019reliability}.








 








 \subsection{Towards Global Convergence Algorithms}
\label{sec:apps-2}

Our approach to get the improved self-selection algorithm was to re-interpret the problem through the lens of coarsening 
and show a local convergence guarantee. 
A natural question is whether and when it is possible to obtain a \emph{global convergence} guarantee. 
This question again reduces to understanding the optimization landscape of \eqref{eq:NLL-intro}. 
In this section, we address this question by focusing on the {fundamental} problem of Gaussian mean estimation from coarse examples. 
This problem has been recently explored by \citet{fotakis2021coarse}. 
They studied the following version of the problem 
(which is slightly more general than the one we introduced before).
\begin{definition}
[Coarse Mean Estimation \cite{fotakis2021coarse}]
Consider the Gaussian distribution $\normal{\mu^\star}{I}$, with mean $\mu^\star \in \R^d$ and identity covariance matrix.
Consider a distribution $\pi$ over partitions of $\R^d$.
We generate a sample as follows:
\begin{enumerate}[itemsep=1pt]
    \item Draw $z$ from $\normal{\mu^\star}{I}$ and, independently, draw a partition $\hyP$ of $\R^d$ from $\pi$.
    \item Observe the unique set $P \in \hyP$ that contains $z$.
\end{enumerate}
We denote the distribution of $S$ as $\cN_\pi(\mu^\star, I)$.
\end{definition}
As we have already discussed, the structure of the sets in the partition $\hyP$ {determines} the optimization landscape of \eqref{eq:NLL-intro}. In fact, in order to obtain a computationally efficient algorithm for estimating $\mu^\star$, we require that all the sets of the partitions are \emph{convex}. One of the results of \cite{fotakis2021coarse} shows that if the convexity condition is dropped, then the problem becomes computationally hard in general.
They prove their hardness result using a reduction from \maxcut{}. Given a \maxcut{} instance, they design
partitions of $\R^d$ consisting of intersections of two halfspaces, ellipsoids, and their complements and (roughly speaking) show that if there were an efficient algorithm for coarse Gaussian mean estimation for general non-convex partitions, this would imply a polynomial time $\sinparen{\nfrac{16}{17}+c}$-approximation algorithm for \maxcut{} with $c > 0$ which is precluded by classical results of \citet{hastad1997inapproximability} {unless $\RP=\NP$}. The simplicity of the sets in the partition indicates that the computational hardness is inherent and not due to overly complicated sets in the partition. 

We complement the negative result of \cite{fotakis2021coarse} with a computationally efficient algorithm when all the sets of the partition(s) are convex and the algorithm is given a warm start.

\begin{inftheorem}[see \cref{thm:convexPartitionMeanEstimation}]
Consider an $\alpha$-information preserving distribution $\pi$ over partitions of $\R^d$ with convex sets.
There exists an algorithm that, given an $R$-warm-start $\mu_0$ for $\mu^\star$ and $n = \tilde O(\nfrac{d}{(\alpha^4 \eps^2)})$ samples from $\cN_\pi(\mu^\star, I)$, computes an estimate $\mu $ in $\poly(d,\nfrac{1}{\alpha},\nfrac{1}{\eps})$ time that satisfies $\|\mu - \mu^\star\|_2 \leq \eps$ with high probability.
\label{infthm:convex}
\end{inftheorem}
We underline that \cite{fotakis2021coarse} only studied the statistical efficiency of the NLL objective but did not deal explicitly with computation. 
{To obtain computational efficiency, our algorithm} has to deal with various additional challenges compared to \cite{fotakis2021coarse}, which we review below.

\paragraph{Techniques for Convex Coarsening Algorithm.}
Our algorithmic approach is to, roughly speaking, run stochastic gradient descent on the negative likelihood objective of \eqref{eq:NLL-intro}. The starting observation, going back to \cite{fotakis2021coarse}, is that the NLL population objective is \emph{convex} with respect to the mean when each set of the partition is convex. This step requires the use of the Brascamp-Lieb inequality.
Observe that, even in this case,
information-preservation is needed because otherwise the (convex) NLL will be ``flat.'' In the convex partition case, information preservation can be used to show \emph{local strong convexity} around the true solution $\mu^\star$. However, outside of this small region, strong convexity no longer holds. In order to obtain an efficient algorithm, we have to (i) bound the second moment of the norm of the gradient and (ii) be able to sample from the Gaussian $\cN(\mu, I)$ conditional on some set $S$ of the convex partition $\hyP$. 

To argue about (i) and (ii), let us compute the gradient of NLL:
\[ \grad\cL \inparen{\mu}~~
    =~~ 
        \mu
        -
        \Ex_{P\sim \coarseNormal{\muStar}{I}{\hyP}}
        {
            \Ex_{\truncatedNormal{\mu}{I}{P}}{
                {[x]}
            }
        }\,.
        \]
Bounding the second moment of the stochastic gradients is not straightforward since the inner and outer expectations in the gradient expressions are not over the same means, and the sets $P$ can be of arbitrarily large diameter.
To overcome this issue, we introduce an idealized class of partitions, which we call local partitions (see \Cref{sec:localPartitions}), and derive an algorithm that recovers the mean
under this ideal class of partitions.
We then show that we can implement this algorithm using samples from the actual class of observed partitions (see \Cref{sec:SGD-local}). Finally, to be able to run PSGD, we also
need to obtain stochastic gradients for $\negLL$. Since the partition is convex, it must consist of polyhedra. Hence, we can employ classical tools from the polyhedral sampling literature \cite{lovasz2006simulated,lovasz2006fast} to implement this step (see \Cref{apx:sampling-polytopes}).


\subsection{Takeaways for Self-Selection and Open Questions}

The tractability for convex partitions and intractability for some non-convex partitions positions the problem of self-selection \emph{in between} {the algorithm in \cref{infthm:convex} and the hardness result of \cite{fotakis2021coarse}}: on the one hand, the self-selection problem of \Cref{def:ssb} gives rise to \emph{non-convex} partitions, and hence the algorithm of \Cref{infthm:convex} is not useful. On the other hand, these non-convex sets are \emph{structured} and \emph{well-behaved}, making existing hardness results not applicable and leaving open the possibility of an SGD-based global convergence algorithm.

\begin{description}
    \item[Open Problem 1.] 
    {
    \emph{Can an SGD-based algorithm solve the self-selection problem in $\poly(d,k,\nfrac{1}{\eps})$ time?}}
\end{description}
More broadly, we pose the following question regarding inference from coarse data.

\begin{description}
    \item[Open Problem 2.] \emph{Are there polynomial-time algorithms for coarse Gaussian parameter estimation that do not require a warm start for convex partitions and for structured non-convex partitions?}
\end{description}


\subsection{Related Work}
Our paper falls under the umbrella of algorithmic statistics for 
\emph{limited-dependent variables} and high-dimensional
\emph{censored statistics.}

\paragraph{Learning with Self-Selection Biases.}  
Bias due to outcome self-selection is a
well-documented phenomenon across statistics, econometrics, and
the social sciences \cite{roy1951earnings,fair1972methods,lee1978estimation,willis1979education,olsen1982distributional,maddala1986limited,heckman1979sample}. We refer to the works of \citet{cherapanamjeri2023selfselection,cherapanamjeri2022auctionEstimation} for an exhaustive overview of applications and of related Econometrics works. Compared to our results, the closest works are those of \cite{cherapanamjeri2023selfselection} and \citet{gaitonde2024selfselection}. \cite{cherapanamjeri2023selfselection} provide two sets of results. 
First, they work under \Cref{def:ssb} and give an algorithm estimates the regression parameters to within $\poly(\nfrac{1}{k})$-error in $\poly(d)\cdot e^{\poly(k)}$. One of the most challenging parts of this result is the proof of \emph{identifiability}, which was the inspiration of their moment-based algorithm. Their second set of results concerns the easier problem, which they call the known-index setting. In this case, apart from observing only the maximum of $k$ models, they also observe the \emph{index} of the selected model. This makes the problem more tractable. In the known index setting, \cite{cherapanamjeri2023selfselection} manage to give an efficient algorithm that goes beyond the maximum criterion and works for any \emph{convex-inducing} selection rule. At a technical level, this algorithmic idea is very similar to the idea of \citet{fotakis2021coarse}, who gave a sample-efficient algorithm for coarse Gaussian mean estimation under convex partitions. In fact, under the viewpoint of coarsening, the known index self-selection problem with the maximum criterion (or more generally convex-inducing selection rules) corresponds to a convex partition of the $k$-dimensional Euclidean space.
The follow-up work by \cite{gaitonde2024selfselection}  improved the sample complexity for the maximum self-selection problem to the near-optimal $O(d) \cdot \poly(k, \nfrac{1}{\eps})$ and the running time to $\poly(d,k,\nfrac{1}{\eps}) + (\nfrac{(\log{k})}{\eps})^{k}$. Interestingly, they also provide a moment-based algorithmic approach that works for sub-Gaussian noise. They further provide an algorithm for the related problem of max-affine regression (where the noise is added after taking the maximum). Their result for max-affine regression is an algorithm with sample complexity $O(d)\cdot \poly(k,\nfrac{1}{\eps})$ and runtime $\poly(d,k,\nfrac{1}{\eps}) + O(k)^{O(k)},$ improving, in some parameter regime, the result of \citet{ghosh2022maxaffine}. We mention that our techniques could be used to obtain local convergence guarantees for max-affine regression; however, such a result already exists by \citet{kim2024max} (handling sub-Gaussian noise).

            {Both \cite{cherapanamjeri2023selfselection} and \cite{gaitonde2024selfselection}'s algorithms for the model in \cref{def:ssb}
    {proceed in two steps. First, they construct a $k$-dimensional subspace containing $w_1^\star,\dots,w_k^\star$ by computing the span of the top eigenvectors of a carefully-constructed weighted covariance matrix. Then, they recover $w_1^\star,\dots,w_k^\star$ from this subspace by} relying} on conditional moments of the observations.
            These works differ in the specific event upon which they perform the conditioning:
            \cite{cherapanamjeri2023selfselection} condition on the event that the covariate is highly correlated with a guess vector,
            while \cite{gaitonde2024selfselection} choose an event with significantly higher probability
            by only requiring the correlation to be somewhat non-trivial.
            In either case, both approaches explicitly compute estimates of their respective conditional moments and use those estimates to locate $\eps$-close guess vectors to some regressor $w^\star_i$
            {by searching over an $\eps$-net over a $k$-dimensional ball}.
            {Our approach of optimizing a suitably chosen likelihood function is quite different {from the brute force methods (over carefully chosen $k$-dimensional subspaces) in these works, and is crucial to avoid the $\eps^{-k}$ dependence suffered by them}.

\paragraph{Inference with Coarse Data.}
Estimation under coarse examples has a long history in Statistics,
going back to the works of \citet{heitjan1990inference,heitjan1991ignorability,heitjan1993ignorability}, who introduce the coarsening at random framework (CAR), generalizing the missing at random (MAR) model of \citet{rubin1976inference,little1989analysis}. These works and follow-ups \cite{little2019statistical,gill1997coarsening,robins1992recovery,van1996efficient,jacobsen1995coarsening,grunwald2003updating}  focus on statistical aspects of coarsening. In these works, coarsening is based on very simple coarsening mechanisms, such as intervals, motivated by survival analysis problems, and the standard estimation method is based on maximum likelihood estimation, whose computational complexity was not carefully treated in this Statistics line of work. 

In terms of algorithms in high dimensions, the work of
\cite{fotakis2021coarse} is the closest to our work regarding coarsening. In terms of distribution learning, their work provides a sample-efficient algorithm for coarse Gaussian mean estimation when the sets of the partition are convex. We improve on that by providing rigorous guarantees for the runtime of this estimation task, given a warm start.
In terms of computational hardness, \cite{fotakis2021coarse} show that if there exists an efficient algorithm for coarse Gaussian mean estimation under general non-convex partitions, then one can approximate \maxcut{} better than $\nfrac{16}{17}$, which is a well-known $\NP$-hard problem. 

    \paragraph{Truncated and Censored Statistics.}
            Our work is also closely related to the literature on learning from censored-truncated data.  
            Truncation and censoring have a long history of work in Statistics \cite{Cohen91} (and \cite{Galton1897,pearson1902v,pearson1908generalised,Lee1915,fisher31,Cohen91,hannon1999estimation,raschke2012inference,maddala1986limited}), Econometrics \cite{maddala1986limited} {(and references therein)}, and a growing recent literature in Computer Science \cite{daskalakis2018efficient,daskalakis2019computationally,daskalakis2020truncated,ilyas2020theoretical,fotakis2020efficient,lee2023learning,lee2024unknown,Kontonis2019EfficientTS,daskalakis2021statistical,plevrakis2021censored,diakonikolas2024statistical,trunc_regression_unknown_var,truncatedDiscrete,truncated_sm,nagarajan2020truncatedMixtureEM,nagarajan2023EMMixtureTruncation,tai2023mixtureCensored}.
            Truncation occurs when samples falling outside of some subset $S^\star$ of the support of the distribution, called \textit{survival set}, are not observed. 
            Truncation arises in a variety of fields from Econometrics \cite{maddala1986limited}, to Astronomy and other physical sciences \cite{woodroofe1985estimating}, to Causal Inference \cite{imbens2015causal,hernan2023causal}.
            Another recent line of work tackles the problem of \textit{testing} whether a given source of data is truncated or not \cite{canonne2020learning, de2023testing, de2024detecting}. 
            The main connection between our algorithms and this line of research is the use of the negative log-likelihood as the key optimization objective. 
            Interestingly, while in the works of \cite{daskalakis2018efficient}, the structure of the truncation set is not related to the computational efficiency of the algorithm (the set should only have mass but can be highly non-convex), in the coarse setting of \Cref{infthm:convex}, the convexity of the sets in the partition is crucial for a favorable optimization landscape.





\paragraph{Learning Bid Distributions in First- and Second-Price Auctions.} 
Our application for self-selection auctions is related to the work of
\citet{cherapanamjeri2022auctionEstimation}. They provide efficient estimation methods for first- and second-price auctions when 
each sample contains the identity of the winner and the price they paid in a sequence of identical
auctions.
Compared to our application, they work in a non-parametric setting and focus on density estimation.
For data coming from first-price auctions and bid distributions with support in $[0,1]$, they provide finite-sample estimation guarantees under
Lévy, Kolmogorov and total variation distance. For the more relevant second-price auction case, they assume that the bid distributions are supported on $[0,1]$ and give an algorithm with running time $\eps^{-O(k)}$ to learn the $k$ CDFs up to a uniform $\eps$-approximation. We shortly mention that, in computer science literature, research has also explored non-parametric methods for estimating bid distributions in first-price auctions, when the econometrician can place bids without affecting the bidders' behavior \cite{blum2015learning}.

There is also an extensive line of work in the Econometrics literature for identification and estimation from auction data. In terms of identification of bid and value distributions from complete or partial observations of bids, the work of \citet{athey2002identification}
shows that with infinite samples, bid distributions are identifiable
for first-price, second-price, ascending (English), and descending (Dutch) auctions. For a survey on non-parametric identification, we refer to \citet{athey2007nonparametric}. Regarding the estimation of the bid distributions, the work of \cite{cherapanamjeri2022auctionEstimation} focuses on the asymmetric case, where the distribution of each bidder can be different. In the simpler symmetric case, there exist estimation guarantees for first- and second-price auctions (\cite{morganti2011estimating,menzel2013large}).

In terms of identification and estimation under parametric or semi-parametric assumptions, there has been important work from the Econometrics line of research (see \eg{}, \cite{athey2006empirical,donald1996identification,athey2011comparing,hendricks2003empirical}). However, we are not familiar with any local convergence method in the high-dimensional setting we are considering in \Cref{sec:2ndPrice}.

 

\begin{remark}[{$k$-th Price Auctions and Machine Autopsy}]
\label{rem:reliability}
Our results for second price auction data can be naturally extended to $k$-th price auctions. 
As mentioned by \cite{cherapanamjeri2022auctionEstimation}, {apart from their importance in auction theory,} these models are {also relevant for} reliability theory and machine autopsy -- which have received significant attention \cite{meilijson1981autopsy,antoine1993identifiability,nowik1990identifiability,gertsbakh1988statistical}; 
{while many of these works focus on the nonparametric setting,
they have natural parametric counterparts where our results are applicable}.
\end{remark}



    
                 
     
\section{Preliminaries {and Notation}}
    {In this section, we introduce preliminaries and (standard) notations used throughout this work.}
% --- begin preliminaries.tex ---

\paragraph{Vectors and Matrices.}
We use standard notations for matrix and vector norms:
For a vector $v \in \R^d$, we use $\|v\|_2$ to denote its $\ell_2$-norm and $\|v\|_\infty$ its $\ell_\infty$-norm.
For a matrix $A = (A_{i,j})_{i \in [m], j \in [n]}$, we use $\|A\|_F \triangleq \sqrt{\sum_{i} \sum_{j} |A_{i,j}|^2}$ to denote its Frobenius norm and $\|A\|_2 \triangleq \max_{v \neq 0} \nfrac{\|A v\|_2}{\|v\|_2}$ for its spectral norm. Recall that $\|A\|_2 \leq \|A\|_F.$

\paragraph{Unit Vectors, Balls, and Miscellaneous Notation.}
Moreover, we set $\wh{v}$ to be the unit vector associated with the direction of $v$, \ie{}, $\wh{v} = \nfrac{v}{\|v\|_2}$ for $v \neq 0.$
Given a vector $v$, we use $B(v,R)$ to denote the $\ell_2$-ball centered at $v$ with radius $R$ and $B_{\infty}(v,R)$ for the associated $\ell_\infty$-ball.
Also, for $v\in \R^d$ and $i\in [d]$,
$v_{-i}\in \R^{d-1}$ denotes the vector obtained from $v$ by removing the $i$-th coordinate.


\paragraph{Span, Projection, and Kronecker Product.}
For two vectors $v,w$, we let $\mathrm{span}(v,w)$ denote their span, \ie{}, the set of
all linear combinations of those vectors. We also denote their inner product as $v^\top w.$
A projection of $v$ onto $w$ (formally to the span of $w$) is denoted by $\proj_w(v) \triangleq (w^\top v) \cdot w/\|w\|_2 = (w^\top v) \cdot \wh{w}.$ This can be extended to projections of $v$ onto a subspace $W$ by writing $W$ as the span of $\{w_1,\dots,w_k\}$ and letting $\proj_W(v) = \sum_{i \in [k]} (w_i^\top v) \cdot \wh{w}_i.$
Given two matrices $A \in \R^{a} \times \R^b, B \in \R^{c} \times \R^d$, we use $A \otimes B \in \R^{ac} \times \R^{bd}$ to denote the Kronecker product between them. When the identity matrix's dimension is unclear from the context, we use $I_k$ and $I_{k\times k}$ to denote the $k\times k$ identity matrix.
For a matrix $W \in \R^{d \times k}$ and $x \in \R^d$, we use $W^\top x$ to denote the result of the multiplication (which is a column vector).

\paragraph{Partitions.}
We usually denote by $\hyP$ a partition of $\R^d$ and by $P$ an element of this partition. We will call a partition convex if each $P \in \hyP$ is convex. We may also define a distribution $\pi$ over different partitions of $\R^d.$

\paragraph{Gaussians, Coarse Gaussians, and Distances Between Distributions.}
We use $\normal{\mu}{\Sigma}$ to denote the Gaussian distribution with parameters $\mu \in \R^d$ and $\Sigma \in \R^{d \times d}. $
The mass of some set $S \subseteq \R^d$ under this distribution is denoted by $\normalMass{\mu}{\Sigma}{S}.$ More broadly, for a distribution $p$  and set $S$, $p(S)$ is the mass of $S$ under $p$. For some partition $\hyP$, we use the notation $\cN_{\hyP}(\mu, \Sigma)$ to denote the coarse Gaussian distribution, defined as
$\cN_{\hyP}(\mu, \Sigma)(P) = \int_P \normal{\mu}{\Sigma}(x) \d x$ for each $P \in \hyP.$ For two distributions $p,q$ over $\R^d$, we denote their total variation distance as $\tv{p}{q} = \inparen{\sfrac{1}{2}} \int_{\R^d} |p(x) - q(x)| \d x$ and their KL divergence as 
$\kl{p}{q}= \E_{x \sim p}[\log ( \nfrac{p(x)}{q(x)} )].$


% --- end preliminaries.tex ---
    


\section{Models and Main Results}
    In this section, we present the main results of the paper. In \Cref{sec:info-pres-noisyChannel}, we give a general abstract formulation of information preservation and present its implication for quadratic growth of the negative log-likelihood objective. In \Cref{sec:info-pres-coarse}, we specialize this observation for coarsening.
    \Cref{sec:self-selection} presents the formal self-selection model with maximum selection criterion and our main results.
    \Cref{sec:second-price} presents the formal self-selection model with second-price data and our local convergence guarantees. 
    Finally,
    in \Cref{sec:coarseConvex}, 
    we present our algorithmic results for coarse Gaussian mean estimation. 
    
    \subsection{Information Preservation for General Noisy Channels}
    \label{sec:info-pres-noisyChannel}

    In this section, we present an abstract framework for information preservation.
    Consider a domain $\cX$ and a family of distributions $\{ P(\theta)\colon \theta \in \Theta\}$ over it.
    Let $f\colon \cX \to \cY$ be a  (deterministic) distortion mechanism that transforms each element $x \in \cX$ to some element of $\cY$.
    This mapping induces, for any distribution $P(\theta)$, a new distribution $P_f(\theta)$ over $\cY$, where $P_f(\theta; y) = \int_{x\colon f(x) = y} P(\theta; x) \d x $ for any $y \in \cY$.
    The statistical task of interest is the following: given \iid{} draws $f(x_1),\dots,f(x_n)$ from the distorted distribution $P_f(\theta^\star)$ for some unknown $\theta^\star \in \Theta$ (without observing $x_1,\dots,x_n)$, the goal is to estimate $\theta^\star.$

    \begin{definition}[Information Preserving Distortion Mechanism]\label{def:model:informationPreservingDistortion}
    We say that $f$ is \emph{$\alpha$-information preserving at radius $R>0$ with respect to $P(\theta^\star)$} if, 
    for any $\theta \in \Theta$ satisfying $\norm{\theta-\theta^\star}_2\leq R$,
    \[
        \tv{P_f(\theta)}{P_f(\theta^\star)}
        \geq \min\inbrace{1, \alpha \norm{\theta - \theta^\star}_2}\,.
    \]
    If the radius $R$ is not explicitly mentioned,
    we understand it to be $R=\infty$.
    \end{definition}
    A simple application of Pinsker's inequality shows the following important implication for the growth of the log-likelihood. Let us set
    \[
    \hyL_f(\theta) 
    =
    -\E_{y \sim P_f(\theta^\star)}
    \insquare{\log P_f(\theta; y)}
    =
    -\E_{y \sim P_f(\theta^\star)}
    \insquare{\log \int_{x\colon f(x) = y} P(\theta; x) d x}\,.
    \]
    
    \begin{lemma}\label{lem:model:informationPreservingImpliesGrowth}
    Consider a distortion mechanism $f$ that is {$\alpha$}-information preserving at radius $R$ with respect to $P(\theta^\star)$. Then, the negative log-likelihood $\hyL_f$ at any $\theta \in \Theta$ where $\norm{\theta-\theta^\star}_2\leq R$ satisfies
    \[
    \hyL_f(\theta)-
    \hyL_f(\theta^\star)
    \geq 
    \min\inbrace{2, 2\alpha^2 \norm{\theta - \theta^\star}_2^2}\,.
    \]
    \end{lemma}
    \begin{proof}
    We have that
    \[
    \hyL_f(\theta)-
    \hyL_f(\theta^\star)
    =
    \E_{y \sim P_f(\theta^\star)}
    \insquare{
    \log(P_f(\theta^\star) / P_f(\theta))
    }
    = \kl{P_f(\theta^\star)}{P_f(\theta)}
    \geq 
    2 \tv{P_f(\theta)}{P_f(\theta^\star)}^2 
    \,.
    \]
    Hence, this implies that
    \[
    \hyL_f(\theta)-
    \hyL_f(\theta^\star)
    \geq  
    \min\inbrace{2, 2\alpha^2 \norm{\theta - \theta^\star}_2^2} 
    \,. \qedhere
    \]
    \end{proof}
    \Cref{lem:model:informationPreservingImpliesGrowth} shows that it suffices to prove information preservation
    in order to show that the true parameter $\theta^\star$ is a local minima of the negative log-likelihood function.
    
    \subsection{Information Preservation for Coarsening}
    \label{sec:info-pres-coarse}
    In this section, we specialize the previous abstract setting to the coarsening mechanism.
        Let $\hyP$ be a partition of $\R^d$.
        Let $\mu\in\R^d$ and $\Sigma\in \R^{d\times d}$ be the parameters of Gaussian distribution $\normal{\mu}{\Sigma}$.
        Define the coarse Gaussian distribution $\coarseNormal{\mu}{\Sigma}{\hyP}$ as the discrete distribution on $\hyP$ satisfying that, for each set $S \in \hyP$, 
        \[
            \Pr_{P \sim \coarseNormal{\mu}{\Sigma}{\hyP}}\insquare{P= S}
            =
            \cN(S ;\mu,\Sigma)\,.
        \]
        With certain partitions (\eg{}, the singleton partition covering the entire domain, \ie{}, $\hyP = \{\R^d\}$), it is information-theoretically impossible to identify the parameters of the underlying Gaussian distribution.
        To ensure that identification is possible, we assume that the partition preserves information as introduced in \Cref{def:informationPreservation}. Our first observation is that one can employ information preservation of an arbitrary partition $\hyP$ to get quadratic growth guarantees for the negative log-likelihood objective
        \begin{equation}
        \label{eq:NLL}
        \hyL(\mu, \Sigma) = \E_{P \sim \coarseNormal{\mu^\star}{\Sigma^\star}{\hyP}}[- \log \cN(P; \mu, \Sigma)] \,.
        \end{equation}
        In the above, $(\mu, \Sigma)$ correspond to guesses of the true uknown parameters $(\mu^\star, \Sigma^\star).$

       The next result is an immediate consequence of \Cref{sec:info-pres-noisyChannel}.
        \begin{theorem}
        \label{thm:info-pres-implies-strong-conv}
        Consider a partition $\hyP$ of $\R^d$ which is $\alpha$-information preserving at radius $R$ with respect to $\normal{\mu^\star}{\Sigma^\star}$. Then the negative log-likelihood satisfies
        \[
        \hyL(\mu, \Sigma)
        -
        \hyL(\mu^\star, \Sigma^\star) 
        \geq  
        \min \inbrace{
            2, 
            2\alpha^2 \inparen{\|\mu - \mu^\star\|_2^2 + \|\Sigma - \Sigma^\star\|_{F}^2}} 
        \]
        for every $\mu, \Sigma$ such that $\inparen{\norm{\mu-\muStar}_2^2 + \norm{\Sigma-\SigmaStar}_F^2}^{\sfrac12}\leq R$. 
        \end{theorem}
        
    \subsection{Linear Regression with Self-Selection Bias}
    \label{sec:self-selection}
    In this section, we formally present our main result for linear regression under self-selection bias \cite{cherapanamjeri2023selfselection,gaitonde2024selfselection}.
    As is usual in multiple linear regression, the goal is to estimate $d\times k$ unknown parameters $\inbrace{w^\star_i\in \R^d\colon 1\leq i\leq k}$ that determine the relation between the independent variable $x\in \R^d$ and the dependent variables $y_1,y_2,\dots,y_k$ as follows 
    \[
        \text{for each $1\leq i\leq k$}\,,\quad 
            y_i = \inangle{x, w_i^\star} + \xi_i 
            \quadwhere \xi_i \sim \cN(0,I_{k})\,. 
    \]
    If we observe samples of the form $\inparen{x,y_1,\dots,y_k}$, then this problem reduces to $k$ independent linear regression problems and can be solved using standard methods, \eg{}, least squares regression.
    Under self-selection, instead of observing $\inparen{x,y_1,\dots,y_k}$, one observes $\inparen{x,f(y_1,\dots,y_k)}$ for some known function $f\colon\R^k\to\R$, which prevents us from using learning each regressor $w^\star_i$ separately. %
    We focus on the max-self-selection bias, where $f(\cdot)=\max(\cdot)$ is the maximum function.
    This setting was studied by \citet{cherapanamjeri2023selfselection,gaitonde2024selfselection}.
    In this model, we observe samples of the form (see \cref{def:ssb})
    \[
        \inparen{x\,,~~ y_{\max}\coloneqq \max_{1\leq i\leq k} y_i}\,.
        \yesnum\label{eq:SS:observations}
    \]
    {See \cref{fig:self-selection-partitions,fig:self-selection-3D} for an illustration fo the resulting coarse partitions with $k=2$ and $k=3$ respectively.}
    \begin{figure}[tbh!]
        \centering
        \includegraphics[width=0.5\linewidth]{fig/self_selection_3D.pdf}
        \caption{
            \label{fig:self-selection-3D}
            {This figure illustrates one set $P$ in the partition $\hyP$ arising in the self-selection problem with $k=3$. 
            The figure illustrates the set $P$ corresponding to the observation $\ymax=1$.
            $P$ is defined as the set of points $(y_1,y_2,y_3)$ with $\max\inbrace{y_1,y_2,y_3}=1$ and $y_1,y_2,y_3\leq 1$
            In other words, $P = \set{y_1= 1, y_2\leq 1, y_3\leq 1}\cup \set{y_1\leq 1, y_2= 1, y_3\leq 1}\cup \set{y_1\leq 1, y_2\leq 1, y_3= 1}$.
            See \cref{fig:self-selection-partitions} for an illustration of the entire partition $\hyP$ with $k=2$.}
        }
    \end{figure}  
    Since the distribution of these observations is invariant to permutations of  $\inbrace{w^\star_i\colon 1\leq i\leq k}$, the goal is to recover the parameters up to permutation.
    For an arbitrary ordering $w^\star_1,w^\star_2,\dots,w^\star_k$ of the parameters, we define $\Wstar$ to be the following matrix:
    \[
        \Wstar \coloneqq \begin{bmatrix}
            w^\star_1 & w^\star_2 & \dots & w^\star_k
        \end{bmatrix}\,.
    \]
    We study this problem under the following assumptions that match the ones in prior work \cite{cherapanamjeri2023selfselection,gaitonde2024selfselection}. First, we assume that the feature vectors $x$ are independently drawn from $\cN(0,I_{d}).$
    \begin{restatable}[Gaussianity]{assumption}{selfSelectionGauss}\label{asmp:SS:gaussianity}
        The feature $x$ follows the $d$-dimensional standard Gaussian distribution. 
    \end{restatable} 
    Next, we impose some separability  and boundedness conditions for the true parameter vectors.
    \begin{restatable}[Separability and Boundedness]{assumption}{selfSelectionSeparation}\label{asmp:SS:separationBoundedness}
        There are known $c\in (0,1]$ and $C\geq 1$ such that 
        \[
            \norm{w^\star_i}_2^2 
            \geq c + \max_{j\neq i}~\sabs{\sinangle{w^\star_j, w^\star_i}} 
            \qquadand
            \max_{1\leq i\leq k}\norm{w^\star_i}_2 \leq C\,.
        \]
    \end{restatable} 

    \noindent Even with this pair of assumptions, it is highly non-trivial to show that $\Wstar$ is identifiable from the data in \cref{eq:SS:observations} (see \cite{cherapanamjeri2023selfselection}for details).
    The Gaussian prior on $x$ (\cref{asmp:SS:gaussianity}) is a classic assumption that is common in similar problems, such as mixtures of linear regressions (\eg{}, \cite{silvia2020EMmixtures}) and is crucially used to show identifiability. 
    Regarding \cref{asmp:SS:separationBoundedness}, 
    its necessity is shown in \cite[Remark 6.1]{gaitonde2024selfselection},
    which gives a simple example demonstrating that at least $k^{\Omega\inparen{\sfrac{C^4}{c^4}}}$ samples are necessary to estimate $W^\star$ to any non-trivial accuracy.

        Our main result is a $\poly(d)\cdot \inparen{\nfrac{1}{\eps}}^2 \cdot 2^{\poly(k)}$ time algorithm to estimate the parameters $W^\star$ (up to permutation) under the above self-selection model. This algorithm improves upon the ${\poly(d)\cdot \inparen{\nfrac{1}{\eps}}^{O(k)} \cdot 2^{\poly(k)}}$ time of the algorithm by \cite{gaitonde2024selfselection} and the $\poly(d)\cdot 2^{\poly(k, {1}/{\eps})}$ algorithm of \cite{cherapanamjeri2023selfselection}.
        \begin{restatable}[Efficient Regression Under Self-Selection]{theorem}{thmSelfSelectionMain}\label{thm:SS}
            Suppose \cref{asmp:SS:gaussianity,asmp:SS:separationBoundedness} hold.
            There is an algorithm that, 
                given any $\eps,\delta\in (0,1)$ 
                and 
                given {$n=O(d) \cdot \poly(k,\nfrac{1}{\eps}, \log(\nicefrac1\delta))$} samples generated by the max-self-selection model with parameters $w^\star_1,w^\star_2,\dots,w^\star_k$,
            outputs a set of estimates $\inbrace{w_i\in \R^d\colon 1\leq i\leq k}$,
            such that, 
            with probability $1-\delta,$
            there is an ordering of these parameters $w_1,w_2,\dots,w_k$ satisfying 
            \[
                \max_{1\leq i\leq k}~\norm{w_i - w^\star_i}_2 \leq \eps\,,
            \]
            The algorithm runs in time \mbox{$\poly(d,k,\nfrac{1}{\eps},{\log(\nicefrac1\delta)}) + {2^{\wt{O}(k)}}$}.
        \end{restatable}
        The main new ingredient required to establish \cref{thm:SS} is a proof that the negative log-likelihood arising in self-selection is strongly convex in a $\poly\inparen{\nfrac{1}{k}}$-sized neighborhood of $\Wstar$ (\cref{thm:SS:localConvexity}).
        Note that globally the negative log-likelihood is highly non-convex: it has at least $k!$ distinct stationary points (corresponding to the permutations of $\Wstar$).
        This local strong convexity is sufficient to deduce \cref{thm:SS} as a point within this $\poly\inparen{\nfrac{1}{k}}$-sized neighborhood can be found using existing algorithms for self-selection and, subsequently, one can obtain an estimate $\eps$-close to $\Wstar$ in $\poly(d,k,\nfrac{1}{\eps})$ time using PSGD.
        To formally state this local convergence result, we need a definition of a warm-start:
        A matrix $W$ is said to be an $R$-warm start for $\Wstar$ (up to permutation of columns) if 
        \[
            \norm{W-W^\star}_F\leq R\,.
            \tag{$R$-Warm Start}
        \]
        \begin{restatable}[Polynomial Time Local Convergence for Self-Selection]{theorem}{thmSelfSelectionLocalConvergence}\label{thm:SS:localConvergence}
            Fix $\eps,\delta\in (0,1)$ and 
            suppose \cref{asmp:SS:gaussianity,asmp:SS:separationBoundedness} hold.
            There is an algorithm that, given 
                a $\poly\inparen{\nfrac{1}{k}}$-warm start
                and
                \[n=\tilde O{\inparen{\nfrac{d}{\eps^2}\cdot {\log(\nicefrac1\delta)}}}\cdot \poly(k)\] 
                samples generated by the max-self-selection model with parameters $w^\star_1,w^\star_2,\dots,w^\star_k$,
            outputs a set of estimates $\inbrace{w_i\in \R^d\colon 1\leq i\leq k}$,
            such that,  with probability $1-\delta,$ there is an ordering of these parameters $w_1,w_2,\dots,w_k$ satisfying 
            \[
                \norm{W - \Wstar}_F\leq \eps\,.
            \]
            The algorithm runs in time \mbox{$\tilde O{\inparen{\nfrac{d^2}{\eps^2}\cdot {\log(\nicefrac1\delta)}}}\cdot \poly(k)$}.
        \end{restatable}
        As mentioned before, \cref{thm:SS:localConvergence} follows by a local strong-convexity property of the negative log-likelihood. 
        Given parameters $W$, define:
        \[
            \text{
                $\cM(W)\coloneqq$ the distribution of {$(x,y_{\max})$} in the self-selection instance corresponding to $W$\,.
            }
        \]
        Further, define $\cL(\cdot)$ as the negative log-likelihood at point $W$.
        We divide the proof of local strong convexity into two parts (quadratic growth and local convexity), as stated in the next result.
        \begin{restatable}[Information Preservation and Local Convexity for Self-Selection]{theorem}{thmSelfSelectionLocalConvexity} 
            \label{thm:SS:localConvexity}
            Suppose \cref{asmp:SS:gaussianity,asmp:SS:separationBoundedness} hold with constants $c,C>0$.
            Fix a permutation $w_1^\star,w_2^\star,\dots,w_k^\star$ of the vectors $\inbrace{w_i^\star:1\leq i\leq k}$.
            The following hold. 
            \begin{enumerate}[leftmargin=15pt]
                \item \textit{(Information Preservation)}
                    For any matrices $V,W\in \R^{d\times k}$ that are $O\inparen{\nfrac{1}{\log{k}}}$-warm starts for $\Wstar$, 
                    \[
                        \tv{\cM(V)}{~\cM(W)} \geq 
                        \inparen{\frac{c}{ek}}^{{O({C^2/c^2}){}}}
                        \cdot \norm{V - W}_F\,.
                        \yesnum\label{eq:SS:informationPreservation}
                    \]
                \item \textit{(Local Convexity)} For any matrix $W\in \R^{d\times k}$ that is a $\poly\inparen{\nfrac{1}{k}}$-warm start for $\Wstar$, 
                    \[
                        \grad^2 \cL(W) \succeq 0\,.
                    \]
            \end{enumerate}
            
        \end{restatable}  
        For the proof of \cref{thm:SS:localConvexity}, we refer the reader to \Cref{sec:SS}.


    \subsection{Second Price Auctions}
    \label{sec:second-price}
    In this section, we study (multiple) linear regression under the self-selection bias arising in a second-price auction \cite{cherapanamjeri2022auctionEstimation}.
    Recall that in a second price auction, there are $k$-bidders and each bidder places a (hidden) bid.
    At the end of the auction, the winner's identity and the price they pay (the second highest bid) are revealed.
    This leads to the following formal model.
    \defSecondPrice*
    \noindent 
    See \cref{fig:2nd-price-observation} for an illustration of the coarse partition created by this model.
    We study this problem under the same assumptions as we considered for the max-self-selection problem.
    \begin{figure}[tbh!]
        \centering
        \includegraphics[width=0.5\linewidth]{fig/2nd_price_one_observation.pdf}
        \caption{
            This figure illustrates one possible observation 
            from the second-price auction model {(\cref{def:secondPrice})}.
            {The figure illustrates the set $P$ corresponding to the observation $\ysmax=1$ and $\imax=1$.}
            $P$ is defined as 
            {the set of points $(y_1,y_2,y_3)$ with $y_1\geq 1$ (as $\imax=1$ and $y_{\imax}\geq \ysmax$ and $\max\inbrace{y_2,y_3} = 1$ (as $\max_{i\neq \imax} y_i=\ysmax$). 
            In other words,} 
            $P = \set{y_1\geq 1, y_2=1, y_3\leq 1}\cup \set{y_1\geq 1, y_2\leq 1, y_3=1}$.
            {Observe that the figure has two slabs: a ``vertical'' one and a ``horizontal'' one.}
            The vertical slab corresponds to $\set{y_1\geq 1, y_2=1, y_3\leq 1}$ and the horizontal slab is $\set{y_1\geq 1, y_2\leq 1, y_3=1}$.
            }
        \label{fig:2nd-price-observation}
    \end{figure}  

    \selfSelectionGauss*
    \selfSelectionSeparation*
    \noindent {The identifiability of this problem follows from similar techniques as \cite{cherapanamjeri2023selfselection}.}
    Our main result is a local convergence algorithm for this problem.
    \begin{restatable}[]{theorem}{thmSecondPriceLocalConvergence}\label{thm:2ndPrice:localConvergence}
        Consider the self-selection model arising from Second Price Auction Data (\cref{def:secondPrice}).
        Suppose \cref{asmp:SS:gaussianity,asmp:SS:separationBoundedness} hold.
        There is an algorithm that, given any $\eps,\delta\in (0,1)$, any $\poly\inparen{\nfrac{1}{k}}$-warm-start for $W^\star$, and \mbox{$n=\wt{O}(\nfrac{d}{\eps^2}\cdot {\log(\nicefrac1\delta)})\cdot \poly(k)$} samples generated by the self-selection model arising from Second Price Auction Data (\cref{def:secondPrice}), outputs an estimate $W$, such that with probability $1-\delta$, 
        \[
            \norm{W-W^\star}_F\leq \eps\,.
        \]
        The algorithm runs in time \mbox{$\wt{O}(\nfrac{d^2}{\eps^2}\cdot {\log(\nicefrac1\delta)} )\cdot \poly(k)$}.
    \end{restatable} 
    As with our main result, the main new ingredients used to prove \cref{thm:2ndPrice:localConvergence} is (1) a lower bound on the information preservation of the problem if provided a warm start and (2) a proof that the negative-log-likelihood is convex close to the optimum $W^\star$.
    We say that a matrix $W$ be an $\alpha$-warm start for $W^\star$ if
    \[
        \norm{W - W^\star}_F\leq \alpha\,.
    \]
    Further, let $\cM(W)$ be the distribution of the observations in the self-selection problem arising in second-price auctions, {\ie{}, the distribution of $(x,i_{\max}, \ysmax)$}.
    Our information preservation and local convexity results are as follows.
    \begin{restatable}[Information Preservation and Local Convexity for Second-Price]{theorem}{thmsecondPriceLocalConvexity} 
        \label{thm:2ndPrice:localConvexity}
        Suppose \cref{asmp:SS:gaussianity,asmp:SS:separationBoundedness} hold with constants $c,C>0$.
        Then, %
        \begin{enumerate}[leftmargin=15pt]
            \item \textit{(Information Preservation)}
                For any matrices $V,W\in \R^{d\times k}$ that are $O\inparen{\nfrac{1}{\log{k}}}$-warm starts for $\Wstar$, 
                \[
                    \tv{\cM(V)}{~\cM(W)} \geq 
                    {\inparen{\frac{c}{ek}}^{{O({C^2/c^2}){}}}}
                    \cdot \norm{V - W}_F\,.
                    \yesnum\label{eq:2ndPrice:informationPreservation}
                \]
            \item \textit{(Local Convexity)} For any matrix $W\in \R^{d\times k}$ that is {a} $\poly\inparen{\nfrac{1}{k}}$-warm start for $\Wstar$, 
                \[
                    \grad^2 \cL(W) \succeq 0\,.
                \]
        \end{enumerate}
        
    \end{restatable}  
    For the proof of \cref{thm:2ndPrice:localConvexity}, we refer the reader to \Cref{sec:2ndPrice}.

    \subsection{Gaussian Mean Estimation From Convex Partitions}
    \label{sec:coarseConvex}
    In this section, we will provide our result for efficiently learning the mean $\mu^\star \in \R^d$ of an isotropic Gaussian  $\normal{\mu^\star}{I}$ from coarse examples when the partition $\hyP$ is $\alpha$-information preserving with respect to $\cN(\mu^\star, I)$ and consists of convex sets.  
    Our result provides a local convergence algorithm; to this end, we will assume that we are provided a warm start by letting the true mean satisfy that $\| \mu^\star \|_2 \leq D.$
    The goal is to design an efficient algorithm that, 
    for any input $\eps, D > 0$, 
    after seeing sufficiently many coarse samples from $\cN_{\hyP}(\mu^\star, I)$, outputs parameter $\mu$ such that with high probability 
$\norm{\mu-\muStar}_{2} \leq \eps$.
     Our main result for this problem is as follows.
    
    \begin{restatable}{theorem}{thmConvexPartitionMeanEstimation}\label{thm:convexPartitionMeanEstimation}
        Let $\eps\in (0, 1]$.
        Suppose $\hyP$ is a convex $\alpha$-information preserving partition of $\R^d$ 
        with respect to $\normal{\muStar}{I}$
        and $\norm{\muStar}_2\leq D$.
        There is an algorithm that outputs an estimate $\tilde \mu$ satisfying
        \[
            \norm{\tilde \mu - \muStar}_2 \leq \eps
        \]
        with probability $1-\delta$. Moreover, the algorithm requires
        \[
            m = \tilde O\inparen{\frac{dD^2\log(\nfrac1\delta)}{\alpha^4} + \frac{d\log(\nfrac1\delta)}{\alpha^4\eps^2}}
        \]
        \iid{} samples from $\coarseNormal{\muStar}{I}{\hyP}$ and $\poly(m, T_s)$ time,
        where $T_s$ is the time complexity of sampling from a Gaussian distribution truncated to a set $P\in \hyP$.
        {Moreover,
        if the facet-complexity\footnote{The \emph{facet-complexity} \cite[Definition 6.2.2]{grotschel1988geometric} of a polytope is a natural measure of the complexity of the polytope.
        See \Cref{apx:sampling-polytopes:removing-oracle-assumption} for the formal definition and more details.}
        of every observed set $P$ is bounded above by $\varphi$,
        we may assume $T_s = \poly(d, \varphi)$.}
    \end{restatable}
    {We remark that similar to \cite{fotakis2021coarse},
    our result can straightforwardly be extended to mixtures of convex partitions. 
    In this setting,
    there are convex partitions $\hyP_1, \dots, \hyP_\ell$ and the observations are sets $P$ where $P\in \hyP_i$ for some $1\leq i\leq \ell$, where $i$ is chosen independently by some prior distribution over $[\ell]$.
    The underlying  data generating distribution $\normal{\muStar}{I}$ remains the same
    but the observations are created by first drawing a partition index $1\leq i\leq \ell$ according to some mixture probability,
    independently drawing $x\sim \normal{\muStar}{I}$,
    and then outputting $P\in \hyP_i$ containing $x$.
     For details about the proof, we refer to \Cref{sec:coarseGaussianEstimation}.
    


\addtocontents{toc}{\protect\setcounter{tocdepth}{2}}
\section{Organization of the Rest of the Paper} 
We organize the rest of the paper as follows:
\begin{itemize}
    \item \cref{sec:SS} provides proofs for linear regression with self-selection bias under the maximum selection criterion; specifically,  \cref{sec:SS:proof:informationPreservation} establishes information preservation and \cref{sec:SS:proof:localConvexity} proves local convexity.
    \item \cref{sec:2ndPrice} presents proofs for linear regression with self-selection bias under the second maximum selection criterion and its implications for second-price auction data; specifically, details on information preservation are in \cref{sec:2ndPrice:proof:informationPreservation} and on local convexity are in \cref{sec:2ndPrice:localConvexity}.
    \item \cref{sec:coarseGaussianEstimation} has proofs for coarse Gaussian mean estimation under convex partitions.
    \item \cref{sec:sgd-analysis} offers a general analysis of Projected Stochastic Gradient Descent (PSGD) for functions that are locally convex and satisfy local quadratic growth, which is key to all our local convergence methods.
\end{itemize}

\section{Linear Regression with Self-Selection Bias}
\label{sec:SS}
% --- begin self_selection.tex ---
In this section, we prove \cref{thm:SS:localConvexity} (restated below) which implies our main result -- an efficient local convergence algorithm for the self-selection with max selection rule. 
\thmSelfSelectionLocalConvexity*
\noindent We divide the proof of \cref{thm:SS:localConvexity} into two parts corresponding to the two claims.
\subsection{Information Preservation}\label{sec:SS:proof:informationPreservation}
     In this section, we prove the information preservation promised in \cref{thm:SS:localConvexity} (namely \cref{eq:SS:informationPreservation}); {we follow the outline in \cref{fig:outline:SS:IP}.}
     
    \begin{figure}[!ht]
          \centering
    \vspace{4mm}
    \resizebox{\textwidth}{!}{%
    \begin{tikzpicture}[node distance=2.5cm, auto]
    \node[myNodeNarrow, line width=1pt] (thm) at (0,0cm) 
        {{Information Preservation (Claim 1 of \cref{thm:SS:localConvexity})}};
    \node[myNodeNarrow, text width=8cm, line width=1pt] (subres1) at (-5cm, -3cm) 
        {
        {Reduction to ``Conditional''\\ Information Preservation (\cref{lem:SS:IP:reductionToConditionalIP})}\\[1mm]
        \begin{tabular}{p{0.01cm}@{}p{7cm}@{}}
          &
          \textcolor{gray}{\textit{\small $\beta$-Information Preservation Conditioned on $\evE$}}\\
          &\textcolor{gray}{\textit{\small $\implies\beta\cdot \Pr[\evE]$-Information Preservation}}
        \end{tabular}
        };
    \node[myNodeNarrow, text width=7.25cm, line width=1pt] (subres3) at (5cm, -2.69cm) 
        {
            {Lower bound on Information Preservation Conditioned on Event $\evE$ (\cref{def:SS:event})}
        };

    \node[myNodeNarrow, draw, line width=1pt] (step1) at (-1.5cm, -4.775cm) {Properties of Event $\evE$};
    \node[myNodeNarrow, draw, text width=7cm, line width=1pt] (step2) at (5cm, -5cm) {Elementary Set Witnesses Information Perservation Conditional on $\evE$ {(Eq.~\eqref{eq:SS:probabilityDiff})}};
    
    \node[myNodeNarrow, text width=3cm, dashed] (probOfE) at (-7.5cm, -7.22cm) {$\Pr[\evE] \geq \poly(\nfrac{1}{k})$ (\cref{lem:SS:probability})};
    \node[myNodeNarrow, text width=3.75cm, dashed] (separation) at (-1.5cm, -7.22cm) {
        $\evE$ Implies Separation (\cref{lem:SS:propertiesOfRho})
    }; 

    \node[myNodeNarrow, dashed, text width=2.5cm] (decomposition) at (2.1cm, -7.45cm) {Decomposition Into Terms A, B, and C \eqref{eq:SS:differenceAsProduct}}; 
    \node[myNodeNarrow, dashed, text width=2.5cm] (bounds) at (5cm, -7.45cm) {Bounds on Terms A, B, and C}; 
    \node[myNodeNarrow, dashed, text width=2.5cm] (lowerBound) at (7.9cm, -7.225cm) {Lower Bound on {Eq.~\eqref{eq:SS:probabilityDiff}}}; 
    
    \node[myNodeNarrow, dashed, text width=2cm] (termC) at (2.5cm, -9.5cm) {Bound on Term C}; 
    \node[myNodeNarrow, dashed, text width=2cm] (termB) at (5cm, -9.5cm) {Bound on Term B}; 
    \node[myNodeNarrow, dashed, text width=2cm] (termA) at (7.5cm, -9.5cm) {Bound on Term A}; 

    \node[myNodeNarrow, dashed, text width=2.75cm] (technicalLemma) at (3.75cm, -11cm) {\cref{lem:SS:cdfratios}}; 

    \draw[-stealth, line width=0.5mm] (thm) -- (subres1.north); 
    \draw[-stealth, line width=0.5mm] (thm) -- (subres3.north);

    \draw[-stealth, line width=0.5mm] (subres3) -- (step1);
    \draw[-stealth, line width=0.5mm] (subres3) -- (step2);
    
    \draw[-stealth, line width=0.5mm] (step1) -- (probOfE);
    \draw[-stealth, line width=0.5mm] (step1) -- (separation);  

    \draw[-stealth, line width=0.5mm] (step2) -- (decomposition);  
    \draw[-stealth, line width=0.5mm] (step2) -- (bounds);   
    \draw[-stealth, line width=0.5mm] (step2) -- (lowerBound);  

    \draw[-stealth, line width=0.5mm] (bounds) -- (termA);  
    \draw[-stealth, line width=0.5mm] (bounds) -- (termB);  
    \draw[-stealth, line width=0.5mm] (bounds) -- (termC);  

    \draw[-stealth, line width=0.5mm] (termB) -- (technicalLemma);  
    \draw[-stealth, line width=0.5mm] (termC) -- (technicalLemma);  
    
    \end{tikzpicture}
    }
        \caption{Outline of Proof of Information Preservation for Self-Selection.}
        \label{fig:outline:SS:IP}
    \end{figure}

    
   
    \noindent Toward proving \cref{{eq:SS:informationPreservation}}, fix any parameters $V=[v_1,\dots,v_k]$ and $W=[w_1,\dots,w_k]$ {close to each other and $\Wstar$ in the following sense} 
    \[
        \norm{V - W^\star}_F\,,~\norm{W - W^\star}_F
        \leq \frac{c^3}{400C^2\log{\nfrac{k}{c}}}\,.
        \yesnum\label{eq:SS:distance}
    \]
    Fix $i$ to be any index satisfying 
    \[
        i
        \in 
        \argmax_{1\leq j\leq k} \norm{v_j - w_j}_2\,.
        \yesnum\label{eq:SS:i}
    \]
    {Observe that due to the Pigeonhole Principle,
    \[
        \norm{v_i - w_i}_2 \geq \frac{1}{\sqrt{k}}\norm{V-W}_F\,.
        \yesnum\label{eq:SS:lowerboundonI}
    \]
    Consider any event $\evE$ over the draw of the covariate $x$ from a self-selection instance; we will specify $\evE$ later and its choice will depend on $V$ and $W$.
        We will sometimes abuse the notation and also treat $\evE$ as a subset of $\R^d$ which contains all covariates $x$ for which the event $\evE$ holds.
    
        \subsubsection{Reduction to ``Conditional'' Information Preservation}
        First, we show that it is sufficient to prove information preservation conditioned on event $\evE$ {(left child of the root in \Cref{fig:outline:SS:IP})}.
        {Recall that $\cM(W)$ denotes the distribution of $(x,\ymax)$ observed by the self-selection model with parameters $W \in \R^{d \times k}$.}
        \begin{lemma}[Reduction to ``Conditional'' Information Preservation]
            \label{lem:SS:IP:reductionToConditionalIP}
            For a non-empty set $\evE\subseteq \R^d$, 
            \[
                \tv{\cM(V)}{~\cM(W)} 
                \geq 
                \Pr_{x\sim \normal{0}{I}}[x\in\evE]\cdot \tv{\cM(V\mid \evE)}{~\cM(W\mid \evE)}\,,
            \]
            where $\cM(V\mid \evE)$ (respectively $\cM(W\mid \evE)$) is the distribution of $(x,\ymax)\sim \cM(V)$ (respectively $(x,\ymax)\sim \cM(W)$) conditioned on $x\in \evE$. 
        \end{lemma}
        \begin{proof}
            Observe that 
            {
            \begin{align*}
                \tv{\cM(V)}{~\cM(W)} 
                &= {\frac12 \int_x\int_{\ymax} \abs{\cM(V)(x, \ymax) - \cM(W)(x, \ymax)}} \\
                &= {\frac12 \int_x \int_{\ymax} \abs{\cM(V\mid x)(\ymax) - \cM(W\mid x)(\ymax)}\cdot \normalMass{0}{I}{~\d x}}\,. %
            \end{align*}
            }
            Here $\cM(V\mid x)$ is the distribution of $\ymax$ conditioned on $x$ for $(x,\ymax)\sim\cM(V)$ and, for a set $S$, $\cM(V\mid x)(S)$ is the mass $\cM(V\mid x)$ assigns to $S$ (and analogously for $\cM(W\mid x)$).
            Thus, flipping the order of the integrals and restricting the integral over $x$ to $\evE$, we get the following lower bound
            \begin{align*}
                \tv{\cM(V)}{~\cM(W)} 
                &\geq {\frac12 \int_{\ymax} \int_{x\in \evE} \abs{\cM(V\mid x)(\ymax) - \cM(W\mid x)(\ymax)}\cdot \normalMass{0}{I}{~\d x}} \\
                &\geq \frac12 \int_{\ymax} \abs{\int_{x\in \evE} \insquare{\cM(V\mid x)(\ymax) - \cM(W\mid x)(\ymax)}\cdot \normalMass{0}{I}{~\d x}} \\
                &= \frac12 \int_{\ymax} \abs{\cM(V\mid \evE)(\ymax) - \cM(W\mid \evE)(\ymax)}\cdot \normalMass{0}{I}{\evE} \yesnum\label{eq:SS:integrateConditionalProbability} \\
                &= \Pr_{x\sim \normal{0}{I}}[x\in\evE]\cdot \tv{\cM(V\mid \evE)}{~\cM(W\mid \evE)}\,.
                \yesnum\label{eq:SS:conditionalInfoPreservationStatement}
            \end{align*} 
            Here $\cM(V\mid \evE)$ (respectively $\cM(W\mid \evE)$) is the distribution of $(x,\ymax)\sim \cM(V)$ (respectively $(x,\ymax)\sim \cM(W)$) conditioned on event $\evE$.
            Further, for any set $S$, $\cM(V\mid \evE)(S)$ (respectively $\cM(W\mid \evE)(S)$) is the mass assigned by distribution $\cM(V\mid \evE)$ to set $S$.
            \Cref{eq:SS:integrateConditionalProbability} follows from the fact that
            \begin{align*}
                \int_{x\in \evE} \Pr[y\mid x]\cdot \Pr[x]
                = \int_{x\in \evE} \Pr[x\mid y]\cdot \Pr[y] %
                = \Pr[\evE\mid y]\cdot \Pr[y] %
                = \Pr[y\mid \evE]\cdot \Pr[\evE]\,.
            \end{align*}
        \end{proof}
        Thus, \cref{lem:SS:IP:reductionToConditionalIP} shows that to prove the information preservation claim in \cref{thm:SS:localConvexity} (\ie{}, \cref{eq:SS:informationPreservation}), it is sufficient to find event $\evE$ such that 
        \begin{align*}
            \tv{\cM(V\mid \evE)}{~\cM(W\mid \evE)} 
            &\geq~~ \frac{1}{\Pr[\evE]} \cdot 
            \inparen{\frac{c}{ek}}^{{O({C^2/c^2}){}}}
            \cdot \norm{V - W}_F \,.
        \end{align*}
        Further, due to \cref{eq:SS:lowerboundonI}, it is also sufficient to prove that
        \begin{align*}
            \tv{\cM(V\mid \evE)}{~\cM(W\mid \evE)} 
            &\geq~~ \frac{1}{\Pr[\evE]} \cdot 
            \inparen{\frac{c}{ek}}^{{O({C^2/c^2}){}}}
            \cdot \norm{v_i - w_i}_2
            \yesnum\label{eq:SS:conditionalInfoPreservation}\,.
        \end{align*}
    }
    \subsubsection{Definition of Event $\evE$ and its Properties}
    We now move on to the right child of the root in \Cref{fig:outline:SS:IP}.
    We will condition on the following type of event, which controls the length of projections of $x$ along certain directions {specified by} $V$ and $W$.
    For a matrix $W \in \R^{d \times k}$, we denote {its columns by $w_1,\dots,w_k$ with $w_i \in \R^d$.}
    \begin{definition}\label{def:SS:event}
        Given an index $1\leq i\leq k$, 
        parameters $W,V\in \R^{d\times k}$, and 
        constants $\gamma\in (0,\nfrac{1}{2})$ and $R\geq 2$, 
        let $\evE=\evE_{i,\gamma,R}$ be the following event 
        \[
           {\left\{ x \in \R^d :  1 ~\leq~ 
            \frac{{x{^\top} \wh{w}_i}}{
                R\sqrt{\log{\nfrac{1}{\gamma}}}
            }\,,
            ~
            \frac{{x{^\top} \wh{u}}}{
                2\sqrt{\log{\nfrac{1}{\gamma}}}
            }
            ~\leq~ 2 \right\}\,.}
        \]
        Here {$\wh w_i$ is the unit vector parallel to $w_i$,}
        $u\coloneqq v_i - {(v_i^\top \wh{w}_i}) \cdot \wh{w}_i$ is the component of $v_i$ orthogonal to $w_i${, and $\wh u$ is its corresponding unit vector}.
        Further, in the special case, where $v_i$ is parallel to $w_i$ (and, hence, $u=0$), the definition of $\evE$ omits the bound on ${x^\top \wh{u}}.$
    \end{definition}
    This event, {which is inspired by the work of \cite{gaitonde2024selfselection},} is relevant to us for two reasons {which are outlined in the left sub-branch of the right branch of \cref{fig:outline:SS:IP}.}
    
    \paragraph{Property 1 ($\evE$ occurs with {$\poly(\nfrac{1}{k})$} probability).} 
        First, $\evE$ occurs with a constant probability for any fixed $\gamma$ and $R$.
        \begin{restatable}[]{lemma}{lemSSprobability}
            \label{lem:SS:probability}
            For any $\gamma\in (0,\nfrac{1}{2}]$, $R\geq 2$, and the corresponding event $\evE=\evE_{i,\gamma,R}$ (\cref{def:SS:event}), 
            \[
                \Pr[\evE] = \gamma^{28+6R^2}\,.
            \]
        \end{restatable} 
        {Eventually, we will set $\gamma = \nfrac{1}{e}$ and $R$ to be of order $\sqrt{\log k}$, which will imply that the event $\evE$ occurs with $\poly(\nfrac{1}{k})$ probability.}
        \Cref{lem:SS:probability} is important as the largest information preservation constant we can deduce from \cref{eq:SS:conditionalInfoPreservation} is $\Pr[\evE]$: this is because when removing the conditioning on $\evE$ to deduce \cref{eq:SS:informationPreservation}, we lose a factor of $\Pr[\evE]$.
        \cref{lem:SS:probability} can be proved using standard tail bounds for Gaussian random variables and the fact that $x\sim \normal{0}{I}$; the proof appears in \cref{sec:proofof:lem:SS:probability}.

    \paragraph{Property 2 ($\evE$ implies ``separation'' between $i$ and other indices).}
            The second reason why $\evE$ is crucial is that, roughly speaking, conditioned on $\evE$, $y_{i,V}$ and $y_{i,W}$ are ``far'' from $y_{j,V}$ and $y_{j,W}$ for any $j\neq i$.
            To make this concrete, consider the following decomposition of $y_{j, V}$ and $y_{j, W}$ (for any $j$)
            \begin{equation}
                \begin{split}
                  \hspace{-2.5mm}y_{j, V} 
                    &= \rho_{j,V} + \zeta_{j,V}
                ~~\text{where}~~
                    ~\rho_{j,V}\coloneqq \proj_{{\rm span}(v_i,w_i)}(x)^\top v_j 
                    ~~\text{and}~~ 
                    \zeta_{j,V}\coloneqq {
                        \proj_{{\rm span}(v_i,w_i)^\bot}(x)^\top v_j
                        + \xi_j
                    }\,,\\
                \hspace{-2.5mm}y_{j, W} 
                    &= \rho_{j,W} + \zeta_{j,W}
                ~\text{where}~~
                    \rho_{j,W}\coloneqq \proj_{{\rm span}(v_i,w_i)}(x)^\top w_j 
                    ~\text{and}~~ 
                    \zeta_{j,W}\coloneqq {
                        \proj_{{\rm span}(v_i,w_i)^\bot}(x)^\top w_j
                        + \xi_{j}
                    }\,.
                \end{split}
                \label{eq:SS:decompositionOfY}
            \end{equation} 
            Since $x\sim \cN(0, I)$, the projections of $x$ to orthogonal subspaces are independent random variables.
            This combined with the fact that $\xi\sim \normal{0}{I}$ implies that $\rho_{j,V},\rho_{j,W},\zeta_{j,V},\zeta_{j,W}$ are Gaussian random variables satisfying the following:
            \begin{center}
                The random variables ~$\inbrace{\rho_{j,V},\rho_{j,W}}$~ are independent of the set of random variables ~$\inbrace{\zeta_{j,V},\zeta_{j,W}}$.
            \end{center}
            One consequence of this is that $\zeta_{j,V},\zeta_{j,W}$ \textit{remain} Gaussian random variables even after conditioning on the event $\evE=\evE_{i,\gamma,R}$.
            To see this, observe that $\evE_{i,\gamma,R}$ only depends on the projection of $x$ onto $\Span(v_i,w_i)$ which affects $\rho_{j,V}$ and $\rho_{j,W}$ but not $\zeta_{j,V}$ and $\zeta_{j,W}$.
            Further, $\zeta_{j,V}$ and $\zeta_{j,W}$ have the following distributions
            \begin{equation}
                \begin{split}
                    \zeta_{j,V}\sim 
                        \cN{\inparen{0, \sigma_{j,V}^2}}
                        &\quadtext{where}
                        \sigma_{j,V}^2 \coloneqq 1 + \norm{\proj_{{\rm span}(v_i,w_i)^\bot}(v_j)}_2^2
                        \,,\\
                    \zeta_{j,W}\sim 
                    \cN{\inparen{0, \sigma_{j,W}^2}}
                        &\quadtext{where}
                        \sigma_{j,W}^2 \coloneqq 1 + \norm{\proj_{{\rm span}(v_i,w_i)^\bot}(w_j)}_2^2\,.
                \end{split}
                \label{eq:SS:distributionOfZeta}
            \end{equation}
            It is informative to see bounds on the above variances:
            \[
                1\leq \sigma_{j,V}^2\,,~\sigma_{j,W}^2\leq 2+C^2\,.
                \label{eq:SS:covarianceUB}
            \]
            To deduce the upper bound, we used that  
            \begin{enumerate}[(i),noitemsep]
                \item for any vector $z$, 
                \mbox{$\snorm{\proj_{{\rm span}(v_i,w_i)^\bot}(z)}_2\leq \snorm{z}$},
                \item $\snorm{v_j},\snorm{w_j}\leq \snorm{w_j^\star}+\nfrac{c^3}{(6C)}$,
                which is implied by the stronger statement of \cref{eq:SS:distance}, and
                \item $\norm{w_i^\star}_2\leq C$ from \cref{asmp:SS:separationBoundedness}.
            \end{enumerate} 
            Now, we can explain the separation induced by $\evE$ more formally:
            At a high level, we will show that, conditioned on $\evE$, the following holds (where $i$ is the index in \cref{eq:SS:i} and $j$ is any other index) 
            \begin{align*}
                \begin{split}
                    \min_{A\in \inbrace{V,W}}~\rho_{i,A} - \rho_{j,A} &\geq \Omega(cR)\,,\quad \\
                    \max_{\ell} \abs{\rho_{\ell,V} - \rho_{\ell,W}} &\leq O({R})\cdot \norm{V-W}_F\,, \quadand\\
                    \abs{\rho_{i,V} - \rho_{i,W}} &\geq 
                    \Omega({R})\cdot \norm{V-W}_F\,.
                \end{split}
                \yesnum\label{eq:SS:informalPropertiesOfRho}
            \end{align*}   
            We pause here to make a subtle note; even after conditioning on event $\evE$, $\rho_{j, V}$ and $\rho_{j, W}$ are random variables (for all $1\leq j\leq k$).
            We will show that the random variables $\rho_{j, V}$ and $\rho_{j, W}$ \textit{always} satisfy the above properties.
            Next, let us discuss the usefulness of \cref{eq:SS:informalPropertiesOfRho}:
            \begin{itemize}
                \item 
                    For a large value of $R$, the first condition implies that $\rho_{i,V}$ and $\rho_{i,W}$ are significantly larger than $\rho_{j,V}$ and $\rho_{j,W}$ respectively.
                    This, combined with the upper bound on the variances of $\zeta_{j,V}$ and $\zeta_{j,W}$, implies that for $R\geq C/c$, $y_{i,V}$ (respectively $y_{i,W}$) is significantly larger than $y_{j,V}$ (respectively $y_{j,W}$) conditioned on $\evE$.
                    We will use this property to show that small changes in $\rho_{j,V}$ and $\rho_{j,W}$ do not affect the total variation distance in \cref{eq:SS:conditionalInfoPreservation} by much.
                    \item The second condition implies that when $V$ and $W$ are close, \ie{}, when $\norm{V-W}_F\ll c$, then $\rho_{j,V}$ and $\rho_{j,W}$ are close to each other for each $j$.
                    Since for any $j\neq i$, small changes between $\rho_{j,V}$ and $\rho_{j,W}$ do not affect the total variation distance in \cref{eq:SS:conditionalInfoPreservation} much, we can just focus on small changes in $\rho_{i,V}$ and $\rho_{i,W}$ (for the index $i$ in \cref{eq:SS:i}).
                    \item So far, we mentioned that due to the first two properties, to lower bound the total variation distance in \cref{eq:SS:conditionalInfoPreservation} it suffices to focus on the $i$-th coordinate.
                    The last condition shows that $\rho_{i,V}$ and $\rho_{i,W}$ are not too close to each other: their distance is at least $\Omega(\norm{V-W}_F)$.
                    This enables us to lower bound the total variation distance in \cref{eq:SS:conditionalInfoPreservation}.
                    It turns out that to obtain a lower bound that is linear in $\norm{V-W}_F$ (which is necessary to imply quadratic growth),
                    the lower bound on $\abs{\rho_{i,V}-\rho_{i,W}}$ must also be linear in $\Omega(\norm{V-W}_F)$.
            \end{itemize}
        The following result formalizes \cref{eq:SS:informalPropertiesOfRho}.
        Its proof relies on carefully applying \cref{asmp:SS:separationBoundedness} and using the properties implied by $\evE$; the complete proof appears in \cref{sec:proofof:lem:SS:propertiesOfRho}. 
        \begin{restatable}[Separation Properties of $\rho$]{lemma}{lemSSpropertiesOfRho}
            \label{lem:SS:propertiesOfRho} 
            Fix any constants $\gamma\in (0,\nfrac{1}{2}]$ and $R\geq 3 + (\sfrac{9C}{c})$.
            Let $i$ be the index in \cref{eq:SS:i}.
            The following guarantees hold {with probability 1}
            conditioned on the event $\evE= \evE_{i,\gamma,R}$ (\cref{def:SS:event}):
            \begin{enumerate}
                \item 
                    For any $j\neq i$,
                    $\rho_{i,V} - \rho_{j, V}\geq 3cR \cdot \sqrt{\log{\nfrac{1}{\gamma}}}$ and 
                    $ \rho_{i, W} - \rho_{j, W}\geq  3cR \cdot  \sqrt{\log{\nfrac{1}{\gamma}}}$;
                \item For each $j$,
                    $\abs{\rho_{j, V} - \rho_{j, W}}\leq 3R \cdot \sqrt{\log{\nfrac{1}{\gamma}}}\cdot \norm{v_j - w_j}_2$; 
                \item Further, for the $i$ in \cref{eq:SS:i}, $\abs{\rho_{i,V} - \rho_{i,W}}\geq \inparen{\sfrac{5R}{6}} \cdot \sqrt{\log{\nfrac{1}{\gamma}}}\cdot \norm{v_i - w_i}_2$.
            \end{enumerate}
        \end{restatable}
    Now, we are ready to explain how we prove \cref{eq:SS:conditionalInfoPreservation}. 
    At a high level, {we will show that} the above lemma implies that the event $\ymax\leq \Delta$ (for some constant $\Delta$) has a very different likelihood with respect to $V$ and $W$ -- enabling us to lower bound the total variation distance in \cref{eq:SS:conditionalInfoPreservation}.
    
\subsubsection{{Lower bound on Information Preservation Conditioned on Event $\evE$}}\label{sec:proofof:eq:SS:conditionalInfoPreservation}
    In this section, we prove \cref{eq:SS:conditionalInfoPreservation}, which we restate below:
    \[
        \tv{\cM(V\mid \evE)}{~\cM(W\mid \evE)} \geq 
        \frac{1}{\Pr[\evE]}
        \cdot 
        {\inparen{\frac{c}{ek}}^{{O({C^2/c^2}){}}}}
        \cdot \norm{v_i - w_i}_F\,,
        \tag{Eq. \eqref{eq:SS:conditionalInfoPreservation} restated}
    \]
    {where $\evE$ is the event from \cref{def:SS:event}.}
    Recall that $\rho_{j,V}$ and $\rho_{j,W}$ are the projections $\rho_{j,V}\coloneqq \proj_{{\rm span}(v_i,w_i)}(x)^\top v_j$ and $\rho_{j,V}\coloneqq \proj_{{\rm span}(v_i,w_i)}(x)^\top v_j$ respectively (where $i$ is the index in \cref{eq:SS:i}).
    Conditioning on $\evE$ bounds the norms of $\proj_{{\rm span}(v_i,w_i)}(x)$ to lie in a specific interval.
    However, it does not fix their values and, therefore, $\rho_V =  \inparen{\rho_{1,V}, \dots, \rho_{k,V}}$ and $\rho_W = \inparen{\rho_{1,W}, \dots, \rho_{k,W}}$ are random variables even after conditioning on $\evE$.
    To prove \cref{eq:SS:conditionalInfoPreservation}, we will show that for \textit{each} possible value $r_V,r_W$ of $\rho_V$ and $\rho_W$ conditioned on $\evE$, it holds that 
    \[
        \tv{\cM(V\mid \evE, \rho_V=r_V,\rho_W=r_W)}{~\cM(W\mid \evE, \rho_V=r_V,\rho_W=r_W)}
        \geq \frac{1}{\Pr[\evE]} \cdot 
        {\inparen{\frac{c}{ek}}^{{O({C^2/c^2}){}}}}
        \cdot \norm{V - W}_F\,. 
    \]
    This is sufficient to prove \cref{eq:SS:conditionalInfoPreservation} due to the following equality {(which can be verified using a similar calculation as in the proof of \Cref{lem:SS:IP:reductionToConditionalIP})}.
    \[
        \tv{\cM(V\mid \evE)}{~\cM(W\mid \evE)}
        =
        \Ex_{r_V,r_W\mid \evE}\insquare{
            \tv{\cM(V\mid \evE, \rho_V=r_V,\rho_W=r_W)}{~\cM(W\mid \evE, \rho_V=r_V,\rho_W=r_W)}
        }\,.
    \]
    In the remainder of this section, we fix arbitrary values of $\rho_V$ and $\rho_W$ that are possible conditioned on $\evE$.
    With some abuse of notation, we continue to use $\rho_V =  \inparen{\rho_{1,V}, \dots, \rho_{k,V}}$ and $\rho_W = \inparen{\rho_{1,W}, \dots, \rho_{k,W}}$ to denote these realizations.
    We also do not explicitly mention the conditioning on $\rho_V$ and $\rho_W$ and use $\cM(V\mid \evE)$ and $\cM(W\mid \evE)$ to denote $\cM(V\mid \evE,\rho_V)$ and $\cM(W\mid \evE,\rho_W)$ respectively.

    \paragraph{Pre-Processing to Simplify Analysis.}
    Now, observe that if we translate all of $\rho_{1,V},\dots,\rho_{k,V}$ and $\rho_{1,W},\dots,\rho_{k,W}$ by the same constant then the total variation distance between $\cM(V\mid \evE)$ and $\cM(W\mid \evE)$ remains unchanged.
    To simplify the arguments, we can translate all $\rho_{1,V},\dots,\rho_{k,V}$ and $\rho_{1,W},\dots,\rho_{k,W}$ to ensure that 
    \[
        \min\inbrace{\rho_{i,V}, \rho_{i,W}} = 0\,.
        \yesnum\label{eq:SS:translation}
    \]
    With some abuse of notation, we continue to use $y_V$, $y_W$, $\rho_V$, and $\rho_W$ to refer to the corresponding vectors in the translated space.
    
    {\paragraph{Elementary Set Witnesses Information Perservation Conditional on $\evE$.}
    Next, we move to the right sub-branch of the right branch of \cref{fig:outline:SS:IP}.
    }
    Since {$\tv{\cM(V\mid \evE)}{~\cM(W\mid \evE)}$ is defined as a supremum over all sets $S\subseteq\R$}
    \[
        \tv{\cM(V\mid \evE)}{~\cM(W\mid \evE)} 
        =
        \max_{S\subseteq \R}
            \abs{
                \Pr\insquare{y_{\max, V} \in S\mid \evE}
                -
                \Pr\insquare{y_{\max, W} \in S\mid \evE}
            }\,,
        \yesnum\label{eq:SS:conditionalTVExpression}
    \]
    it suffices to prove that {a specific set -- in our case $\R_{\leq 0}$ -- witnesses that $\tv{\cM(V\mid \evE)}{~\cM(W\mid \evE)} $ is large; \ie{}, it is sufficient to prove the following,}
    \[
        \abs{
            \Pr\insquare{y_{\max, V} \leq 0\mid \evE}
            -
            \Pr\insquare{y_{\max, W} \leq 0\mid \evE}
        }
        \geq 
        \frac{1}{\Pr[\evE]}
        \cdot {\inparen{\frac{c}{ek}}^{{O({C^2/c^2}){}}}}
        \cdot \norm{v_i-w_i}_2\,.
        \yesnum\label{eq:SS:probabilityDiff}
    \]
    To prove this, we compute expressions for the terms on the left-hand side:
    \begin{align*}
        \Pr\insquare{y_{\max, V} \leq 0\mid \evE}
        &= \prod_{j}~ 
            \Pr\insquare{\zeta_{j,V}\leq -\rho_{j,V}\mid \evE}
        ~\hspace{0.25mm}= 
            \prod_{j}~ \customcdf{-\rho_{j,V}}{j,V}
        \,,\\
        \Pr\insquare{y_{\max, W} \leq 0\mid \evE}
        &= \prod_{j}~ 
            \Pr\insquare{\zeta_{j,W}\leq -\rho_{j,W}\mid \evE}
        = \prod_{j}~ \customcdf{-\rho_{j,W}}{j,W}\,.
    \end{align*}
    Where, for any $\sigma^2\geq 0$, $\customcdf{\cdot}{}$ is the cumulative density function of the one-dimensional normal distribution $\cN(0,\sigma^2)$.
    {Hence, we get the following decomposition mentioned in \cref{fig:outline:SS:IP}}
    \begin{align*}
        \hspace{-2.5mm}
        &\abs{
                \Pr\insquare{y_{\max, V} \leq 0\mid \evE}
                - \Pr\insquare{y_{\max, W} \leq 0\mid \evE}
        } \\
        &= 
        \underbrace{\inparen{
            \prod_{j}~ \customcdf{-\rho_{j,V}}{j,V}
        }}_{\textsf{A}}
        \cdot  
        \biggabs{
            1 - 
            \underbrace{\frac{
                \customcdf{-\rho_{i,W}}{i,W}
            }{
                \customcdf{-\rho_{i,V}}{i,V}
            }}_{\textsf{B}}
            \cdot 
            \underbrace{\prod_{j\neq i}~ \frac{
                \customcdf{-\rho_{j,W}}{j,W}
            }{
                \customcdf{-\rho_{j,V}}{j,V}
            }}_{\textsf{C}}
        }\,.
        \yesnum\label{eq:SS:differenceAsProduct}
    \end{align*}
    {Next, following \cref{fig:outline:SS:IP}, we divide the remaining proof into two steps.}
    {The first step has three sub-parts, which} bound the Terms~\textsf{A}, \textsf{B}, and \textsf{C} in \cref{eq:SS:differenceAsProduct}.
    {The second step} uses these bounds to prove \cref{eq:SS:probabilityDiff} which implies \cref{eq:SS:conditionalInfoPreservation}.
    
    To bound Terms~\textsf{B} and \textsf{C}, we use the following result to control the ratios in the above equation.
    \begin{restatable}[Separating CDF Ratios]{lemma}{lemSScdfratios}\label{lem:SS:cdfratios}
        For any constants $\gamma\in (0,\nfrac{1}{2})$ and $R\geq 3 + (\sfrac{10C}{c})$, index $i$ in \cref{eq:SS:i}, and $j\neq i$
        \begin{align*}
                \frac{
                    \customcdf{-\rho_{i,V}}{i,V}
                }{
                    \customcdf{-\rho_{i,W}}{i,W}
                }
                &~~\notin~~ 
                \insquare{
                    1 \pm \frac{1}{100C}\cdot 
                    \min\inbrace{
                        1\,,
                        R \sqrt{\log{\nfrac{1}{\gamma}}}\cdot \norm{v_i-w_i}_2 
                    }
                }\,,
                \yesnum\label{eq:SS:cdfratios:lb}\\
                \frac{\customcdf{-\rho_{j,V}}{j,V}}{\customcdf{-\rho_{j,W}}{j,W}}
                &~~\in~~
                \insquare{
                    1
                    \pm 10R\sqrt{\log{\nfrac{1}{\gamma}}}\cdot 
                    \inparen{\norm{v_j-w_j}_2+C\norm{v_j-w_j}_2^2}
                    \cdot \gamma^{-\frac{c^2R^2}{72C^2}}
                }
                \yesnum\label{eq:SS:cdfratios:ub}
                \,.
        \end{align*}
    \end{restatable}
    The proof of \cref{lem:SS:cdfratios} 
    {is based on carefully analyzing moments of truncated Gaussians with one-sided truncation and is deferred to} 
    \cref{sec:proofof:lem:SS:cdfratios}.
    In the remainder of the proof, we fix the following values of $R$ and $\gamma$  
    \[
        R = \frac{\sqrt{72}C}{c}\cdot \sqrt{
            20 + 3\log{\frac{k}{c}} + 2\log{C}
        }
        \qquadand
        \gamma = \frac{1}{e}\,.
        \yesnum\label{eq:SS:valueRandGamma}
    \] 
    Note that, since $\Pr[\evE]\geq e^{-{28-6R^2}}$ and $R^2={72c^{-2}C^2}\cdot \inparen{
    \proofconstant{} + \log{\nfrac{k}{c}}+ 2\log{C}}$, it holds that
    \[
        \Pr[\evE] 
        =  e^{-28 - 432C^2 c^{-2} (20 + 3 \log(k/c) + 2 \log C)}
        =  e^{-28} e^{-8640 C^2/c^2} \cdot 
            C^{-864 C^2/c^2} \cdot 
            \inparen{\frac{c}{k}}^{1296 C^2/c^2}\,.
        \yesnum\label{eq:SS:IP:lowerBoundOnProb}
    \] 

    \paragraph{Step 1.1 (Bound on Term~\textsf{C}).}
    First, we will simplify Term~\textsf{C} in \cref{eq:SS:differenceAsProduct}:
    The CDF Ratio Lemma (\cref{lem:SS:cdfratios}) implies that 
    \begin{align*}
        \textsf{C} 
        =
            \prod_{j\neq i}~ \frac{
                \customcdf{-\rho_{j,W}}{j,W}
            }{
                \customcdf{-\rho_{j,V}}{j,V}
        }
        \in
        \inparen{
            1
                \pm 10R\sqrt{\log{\nfrac{1}{\gamma}}}\cdot 
                \inparen{\norm{v_j-w_j}_2+C\norm{v_j-w_j}_2^2}
                \cdot \gamma^{-\frac{c^2R^2}{72C^2}}
        }^{k-1}
        \,.
    \end{align*}
    Substituting the values of $R$ and $\gamma$ from \cref{eq:SS:valueRandGamma} and simplifying implies that 
    \[
        \textsf{C} 
        \in
        \inparen{
            1
                \pm \frac{10\sqrt{72}}{C\cdot e^{20}}
                    \cdot 
                    \sqrt{
                        20 + 3\log{\frac{k}{c}}
                        +2 \log{C}
                    } 
                \cdot 
                \inparen{\norm{v_j-w_j}_2+C\norm{v_j-w_j}_2^2}
                \cdot {\frac{c^{2}}{k^{\proofconstant{}}}}
        }^{k-1}\,.
    \]
    Using that $\sqrt{20+3\log{\nfrac{k}{c}}+2\log{C}}\leq \sqrt{2\log{C}}\cdot \nfrac{5k}{c}$ as $C,\nfrac{k}{c}\geq 1$,  implies that 
    \[
        \textsf{C} 
        \in
        \inparen{
            1
                \pm \frac{50\sqrt{72}}{e^{20}}
                \cdot 
                \frac{\sqrt{2\log{C}}}{C}\cdot
                \inparen{\norm{v_j-w_j}_2+C\norm{v_j-w_j}_2^2}
                \cdot {\frac{c}{k^{2}}}
        }^{k-1}\,.
    \]
    Since $\norm{v_j-w_j}_2\leq \norm{V-W}_F\leq 1/C$ (\cref{eq:SS:distance}) and $\nfrac{\sqrt{2\log{C}}}{C}\leq 1$ (as $\nfrac{\sqrt{2\log{z}}}{z}\leq 1$ for all $z>0$), it follows that 
    \[
        \textsf{C} 
        \in
        \inparen{
            1
                \pm \frac{100\sqrt{72}}{e^{20}} 
                \cdot 
                \norm{v_j-w_j}_2
                \cdot {\frac{c}{k^{2}}}
        }^{k-1}\,.
        \yesnum\label{eq:SS:termC}
    \]
    Next, consider following standard inequalities 
    \begin{align*}
        \begin{split}
            \text{for all $k\in \N$ and $z\geq 1$\,,}\qquad 
                (1-z)^{k-1} &\geq 1-(k-1)z\,,\\
            \text{for all $0\leq z\leq \frac{1}{2(k-1)}$\,,}\qquad 
                \inparen{1+z}^{k-1} &\leq \frac{1}{1-(k-1)z}\leq 1+2(k-1)\,.
        \end{split}
        \yesnum\label{eq:bernouliiInequality}
    \end{align*}
    To use the above inequalities with \cref{eq:SS:termC}, we need the following observation 
    \[
        \frac{100\sqrt{72}}{e^{20}} 
                \cdot 
                \norm{v_j-w_j}_2
                \cdot \frac{c}{k^{2}}
        ~~\Stackrel{\rm\eqref{eq:SS:distance}}{\leq}~~ \frac{1}{2}\frac{c}{k^{2}}
        ~~\quad\Stackrel{(c\leq 1\,,~k\geq 1)}{\leq}\quad~~ \frac{1}{2k}\,,
    \]
    Substituting the aforementioned standard inequalities with the above observation in \cref{eq:SS:termC} implies that 
    \[
        \textsf{C} 
            \in 1 \pm \frac{100\sqrt{72}}{e^{20}} 
                \cdot 
                \norm{v_j-w_j}_2
                \cdot {\frac{c}{k}}\,.
\yesnum\label{eq:SS:termC:final}
    \] 
    
    \paragraph{Step 1.2 (Bound on Term~\textsf{B}).}
    Substituting the values of $R$ and $\gamma$ (\cref{eq:SS:valueRandGamma}) into the first part of \cref{lem:SS:cdfratios} implies that
    \[
        \textsf{B} \notin 1 \pm \frac{1}{100C}\cdot 
                \min\inbrace{
                    1\,,
                    \frac{\sqrt{72}C}{c}\cdot \sqrt{
                        20 + 3\log{\frac{k}{c}}
                        + 2\log{C}
                    }\cdot \norm{v_i-w_i}_2 
                }\,.
            \yesnum\label{eq:SS:termB}
    \]
    Observe that \cref{eq:SS:distance} implies that 
    \[
        \norm{v_i-w_i}_2\leq \norm{V-W}_F\leq \frac{c}{\sqrt{72}C^2}\cdot \frac{1}{10\log{\nfrac{k}{c}}}
        ~~~\quad\Stackrel{(k\geq 2,~ c\leq 1,~ C\geq 1)}{\leq}\quad~~~
        \frac{c}{\sqrt{72}C}\cdot 
        \frac{1}{
            \sqrt{20+3\log{\nfrac{k}{c}}+2\log{C}}
        }\,.
    \]
    Above, we use that $k\geq 2$.
    The edge case of $k=1$ can be handled separately: for $k=1$, the problem reduces to a standard linear regression problem for which strong convexity is well-known.
    Therefore, the minimum in \cref{eq:SS:termB} always evaluates to the second term and, hence,  
    \[
        \textsf{B}\notin 1\pm \frac{6\sqrt{2}}{100c} \cdot \sqrt{20+3\log{\frac{k}{c}}+2\log{C}}\cdot \norm{v_i-w_i}_2\,.
    \]
    Since $20+3\log{\nfrac{k}{c}}+2\log{C}\geq 1$ and $c\leq 1$, it also follows that 
    \[
        \textsf{B}\notin 1\pm \frac{6\sqrt{2}}{100} \cdot \norm{v_i-w_i}_2\,.
        \yesnum\label{eq:SS:termB:final}
    \]
    
    \paragraph{Step 1.3 (Bound on Term~\textsf{A}).}
        Finally, we lower bound Term~\textsf{A} in \cref{eq:SS:differenceAsProduct}.
        Recall that after the translation we described at the start of \cref{sec:proofof:eq:SS:conditionalInfoPreservation}, we have that $\min\inbrace{\rho_{i,V}, \rho_{i,W}}=0$.
        Hence, the separation properties of $\rho$ (\cref{lem:SS:propertiesOfRho}), 
        the choice of $R, \gamma$ (\cref{eq:SS:valueRandGamma}), 
        and the proximity of $V, W, \Wstar$ (\cref{eq:SS:distance}) imply that 
    \[
        \rho_{j,V}\leq 
        -{\sqrt{72}C}\cdot \sqrt{
            \proofconstant{} + \log{\frac{k}{c}}
        }
        \qquadand
        \rho_{i,V}  
        \leq \frac{1}{10}
        \,. 
    \]
    Further, from \cref{eq:SS:covarianceUB} we have that 
    \[
        1\leq \sigma_{j,V}^2\,,~ \sigma_{i,V}^2~\leq~ 2+C^2~~~\Stackrel{(C\geq 1)}{\leq}~~~ 3C^2\,.
    \]
    Furthermore, we have the following standard fact: for all $z,\sigma^2>0$
    \[
        \Phi(z;\sigma^2) 
        = 
            \frac{1}{\sigma}\cdot \Phi(\nfrac{z}{\sigma}; 1)
        \quad\Stackrel{\rm\cref{fact:GaussianTail}}{\geq}\quad 
            1 - \frac{e^{-z^2/(2\sigma^2)}}{\sqrt{2\pi}z}
        \,.
    \]
    Combining the above three statements, 
    it follows that 
    \begin{align*}
        \customcdf{-\rho_{j,V}}{j,V}
            &\geq   1 
                    - \frac{
                        e^{
                            -\frac{72C^2\inparen{3+\log{\nfrac{k}{c}}}}{6C^2}
                        }
                    }{\sqrt{2\pi(2+C^2)}}
            \geq 
                1 - \inparen{\frac{c}{k}}^{12} \frac{e^{-36}}{\sqrt{2\pi(2+C^2)}}
            ~~\quad\Stackrel{(C\geq 1,c\leq 1)}{\geq}\quad~~ 
                1 - {\frac{1}{k^{12}}}\,,
                \yesnum\label{eq:SS:infoPreservation:toReferenceLater1}\\
        \customcdf{-\rho_{i,V}}{i,V}&
            ~~\quad\Stackrel{(\rho_{i,V}\leq \sfrac{1}{10})}{\geq}\quad~~
            \customcdf{-\nfrac{1}{10}}{i,V}
            \quad\Stackrel{(\sigma_{i,V}^2\geq 1)}{\geq}\quad
            \Phi\inparen{-\nfrac{1}{10}; 1}
            \geq  \frac{1}{3}\,.
            \yesnum\label{eq:SS:infoPreservation:toReferenceLater2}
    \end{align*}
    Hence, it follows that 
    \[
        \textsf{A}
        = \prod_{j}~ \customcdf{-\rho_{j,V}}{j,V}
        \geq 
        \frac{1}{3} 
        \inparen{1 - \frac{1}{k^{12}}}^{k-1}
        ~~\Stackrel{\eqref{eq:bernouliiInequality}}{\geq}~~ 
        \frac{1}{3}\inparen{1 - \frac{1}{k^{11}}}
        \quad\Stackrel{(k\geq 2)}{\geq}\quad \frac{1}{4}\,.
        \yesnum\label{eq:SS:termA:final}
    \]
    Again, we use that $k\geq 2$.
    The edge case of $k=1$  reduces to a standard linear regression problem for which strong convexity is well-known.

    
    \paragraph{Step 2 (Completing the proof of \cref{eq:SS:conditionalInfoPreservation}).}
        In this step, we prove \cref{eq:SS:probabilityDiff} which {as proved above} implies \cref{eq:SS:conditionalInfoPreservation} {and, in turn, the information preservation claim in \cref{thm:SS:localConvexity}}.
        We use the following simple fact.
        \begin{fact}\label{fact:simpleInequality}
            Fix any $a,b\in [-1,1]$ with $\abs{b}\leq \nfrac{\abs{a}}{10}$.
            For any $z_1\notin 1\pm a$ and $z_2\in 1\pm b$, $\abs{1 - z_1z_2}\geq \nfrac{\abs{a}}{2}$.
        \end{fact}
        \begin{proof}[Proof of \cref{fact:simpleInequality}]
            Toward a contradiction suppose that $\abs{1 - z_1z_2} < 0.5\abs{a}$ and, hence, in particular, $1-z_1z_2\geq -0.5\abs{a}$.
            On rearranging this gives, $z_1z_2\leq 1+0.5\abs{a}$.
            Further, since $z_2\in 1\pm b$ and $\abs{b}\leq 0.1\abs{a}$, we get the following contradiction:
            \[
                z_1 \leq \frac{1 + 0.5\abs{a}}{1-0.1\abs{a}}
                \leq \inparen{1 + 0.5\abs{a}}\inparen{1 + 0.2\abs{a}}
                \quad\Stackrel{(\abs{a}\leq 1)}{\leq}\quad 1 + 0.8\abs{a}\,.
            \]
            Where in the second inequality we used that $\nfrac{1}{(1-z)}\leq 1+2z$ for any $0\leq z\leq 1$ and $0\leq \abs{a}\leq 1$.
        \end{proof}
        The above fact combined with \cref{eq:SS:termC:final,eq:SS:termB:final} implies that 
    \begin{align*}
        \abs{
            1 - 
            \frac{
                \customcdf{-\rho_{i,W}}{i,W}
            }{
                \customcdf{-\rho_{i,V}}{i,V}
            }
            \times 
            \prod_{j\neq i}~ \frac{
                \customcdf{-\rho_{j,W}}{j,W}
            }{
                \customcdf{-\rho_{j,V}}{j,V}
            }
        }
        \geq  
        \frac{3\sqrt{2}}{100} \cdot \norm{v_i-w_i}_2\,.
        \yesnum\label{eq:SS:secondTermLB}
    \end{align*}
    Substituting this lower bound and the lower bound in \cref{eq:SS:termA:final} into \cref{eq:SS:differenceAsProduct} implies 
    \[
        \abs{
            \Pr\insquare{y_{\max, V} \leq 0\mid \evE}
            -
            \Pr\insquare{y_{\max, W} \leq 0\mid \evE}
        }
        \geq \frac{3\sqrt{2}}{400} \cdot \norm{v_i-w_i}_2\,.
    \]
    This implies \cref{eq:SS:probabilityDiff} due to \cref{eq:SS:IP:lowerBoundOnProb}.

\subsection{Local Convexity}\label{sec:SS:proof:localConvexity}
In this section, we prove the local convexity of the negative log-likelihood (\cref{thm:SS:localConvexity}). We begin by defining the negative log-likelihood and presenting its Hessian. 
Then, we prove local convexity in \cref{sec:SS:proof:localConvexity:main}; see \cref{fig:outline:SS:LC} for an outline of the proof.
\begin{figure}[ht]
    \centering
    \vspace{4mm}
    \begin{tikzpicture}[node distance=2.5cm, auto]
    \node[myNodeNarrow, line width=1pt] (thm) at (0,0) 
        {Local Convexity (Claim 2 of \cref{thm:SS:localConvexity})};
    \node[myNode, text width=7.25cm, line width=1pt] (subres1) at (-4cm, -2cm) 
        {\mbox{Strong Convexity at $\Wstar$} (\cref{eq:SS:LC:hessianLB})};
    \node[myNode, text width=7.25cm, line width=1pt] (subres3) at (4cm, -2cm) 
        {Strong Convexity Near $\Wstar$ (\cref{thm:SS:LC:formal})};

    \node[myNodeNarrow, text width=5cm, line width=1pt] (lem51) at (-4cm, -4cm) {Information Preservation (Claim 1 of \cref{thm:SS:localConvexity})};
    
    \node[myNodeNarrow, draw, line width=1pt] (step1) at (2cm, -4cm) {\cref{lem:SS:LC:step1}};
    \node[myNodeNarrow, draw, line width=1pt] (step2) at (6cm, -4cm) {\cref{lem:SS:LC:step2}};
    
    \node[myNodeNarrow, dashed] (propOfE) at (-1cm, -6cm) {Properties of ``Norm-Bound'' Event $\evE$};
    \node[myNodeNarrow, dashed] (tailBounds) at (5cm, -6cm) {Tail Bounds}; 

    \node[myNodeNarrow, dashed, text width=2.5cm] (tailBoundI) at (3.65cm, -8.225cm) {Bounds on Conditional Expectations (\cref{lem:SS:LC:tailBoundI})}; 
    \node[myNodeNarrow, dashed, text width=2.5cm] (tailBoundII) at (6.55cm, -8.225cm) {Bounds on Unconditional Expectations (\cref{lem:SS:LC:tailBoundII})}; 
    
    \node[myNodeNarrow, dashed, text width=2.75cm] (whp) at (-2.6cm, -8cm) {High Probability Guarantee (\cref{lem:SS:LC:prob})}; 
    \node[myNodeNarrow, dashed, text width=2.75cm] (changeOfMeasure) at (0.6cm, -8cm) {Change of Measure II (\cref{lem:SS:LC:changeOfMeasure:corr})}; 

    \node[myNodeNarrow, dashed, text width=2.75cm] (changeOfMeasureII) at (0.6cm, -10cm) {Change of Measure I (\cref{lem:SS:LC:changeOfMeasure})}; 
    
    \draw[-stealth, line width=0.5mm] (thm) -- (subres1.north); 
    \draw[-stealth, line width=0.5mm] (thm) -- (subres3.north);

    \draw[-stealth, line width=0.5mm] (subres1) -- (lem51); 

    \draw[-stealth, line width=0.5mm] (subres3) -- (step1);
    \draw[-stealth, line width=0.5mm] (subres3) -- (step2);
    
    \draw[-stealth, line width=0.5mm] (step1) -- (propOfE);
    \draw[-stealth, line width=0.5mm] (step1) -- (tailBounds);  

    \draw[-stealth, line width=0.5mm] (tailBounds) -- (tailBoundI);  
    \draw[-stealth, line width=0.5mm] (tailBounds) -- (tailBoundII);  

    \draw[-stealth, line width=0.5mm] (propOfE) -- (whp);
    \draw[-stealth, line width=0.5mm] (propOfE) -- (changeOfMeasure); 

    \draw[-stealth, line width=0.5mm] (changeOfMeasure) -- (changeOfMeasureII); 
    
    \end{tikzpicture}
    \caption{Outline of Proof of Local Convexity for Self-Selection.
    }
    \label{fig:outline:SS:LC}
\end{figure}


\subsubsection{Negative Log-Likelihood and Its Hessian}
    Recall that given an observation $\ymax=m$, $y$ lies in the following coarse set (see \cref{eq:L-shapeSets})
    \[
    P_m \coloneqq \inbrace{z \in \R^k : \max_i z_i = m}\,,
    \]
    For an observation $(x,\ymax)$, the sample negative log-likelihood is
    \[
    \cL(W; x,\ymax) = -\log\int_{P_{\ymax}} \exp\Bigl(-\frac{1}{2}\|z-W^\top x\|_2^2\Bigr)\d z\,,
    \]
    and, hence, the (population) negative log-likelihood is given by
    \[
    \cL(W) = \Ex_{(x,\ymax)}\cL(W; x,\ymax)\,.
    \]
    The gradients and Hessians of the likelihood are as follows (see \cref{sec:hessian:SS} for a proof):
    \begin{align}
    \grad \cL(W; x,\ymax) &= xx^\top W - \Ex_{z\sim \normal{W^\top x}{I}}\Bigl[xz^\top \mid z\in P_{\ymax}\Bigr],\\[1mm]
    \grad \cL(W) &= W - {\Ex_{(x,\ymax)}}\Ex_{z\sim \normal{W^\top x}{I}}\bigl[xz^\top \mid z\in P_{\ymax}\bigr]\,,\\[1mm]
    \grad^2 \cL(W; x,\ymax) &= xx^\top\otimes I_{k} - \cov_{z\sim \normal{W^\top x}{I}}\Bigl[z \mid z\in P_{\ymax}\Bigr]\otimes xx^\top\,, \label{eq:SS:conditionalHess:mainBody}\\[1mm]
    \grad^2 \cL(W) &= I_{dk} - \Ex_{(x,\ymax)}\inparen{\cov_{z\sim \normal{W^\top x}{I}}\bigl[z \mid z\in P_{\ymax}\bigr]\otimes xx^\top}\,.
    \end{align}

\subsubsection{Proof of Local Convexity}\label{sec:SS:proof:localConvexity:main}
        Our goal is to show that 
        \[
            \forall~W\quadtext{such that}
            \norm{W-\Wstar}_2 \leq k^{-O(\sfrac{C^2}{c^2})}\,,
            \quadtext{it holds} \nabla^2 \cL(W) \succeq 0 \,.
            \yesnum\label{eq:SS:LC:sufficientCondition}
        \] 
    
    \paragraph{Strong Convexity at $W^\star$.}
    First, we will prove strong convexity of the NLL at the true parameter $W = W^\star$.
    Since $\alpha$-information preservation implies $\alpha$-quadratic growth (see \cref{thm:quadGrowth}), the information preservation property in \cref{eq:SS:informationPreservation} (proved in \cref{sec:SS:proof:informationPreservation}) yields, for some \textit{absolute} constant $A\ge 1$,
    \[
        \phantom{.}
        \qquad\qquad\qquad\qquad\qquad\quad~~
        \nabla^2 \cL(W^\star) \succeq k^{-A\cdot C^2/c^2}\, I\,. 
        \qquad~~
        \text{(Strong Convexity at $\Wstar$)}\hspace{-10mm}
        \yesnum\label{eq:SS:LC:hessianLB}
    \]
    Indeed, a Taylor expansion of $\cL(W)$ around $W^\star$ shows that, as $W\to W^\star$, the higher-order terms become negligible compared to the Hessian (since $W^\star$ is a stationary point of $\cL$ -- see \cref{fact:SS:stationaryPoint} -- the first-order term is zero), thus establishing \cref{eq:SS:LC:hessianLB}. This completes the left branch of \Cref{fig:outline:SS:LC}.

    \paragraph{Strong Convexity Near $W^\star$.}
    We now move to the right branch of the proof outline (\Cref{fig:outline:SS:LC}).
    Given \cref{eq:SS:LC:hessianLB}, local convexity is implied by the following lower bound on $\nabla^2 \negLL(\cdot)$ around $\Wstar$.
    To simplify the notation,
    define $B\geq 1$ as the following quantity 
    \[
        B\coloneqq O\inparen{\frac{\sqrt{A}C}{c}}\,,
        \yesnum\label{eq:SS:LC:defB}
    \]
    where the hidden constant is appropriately large and can be deduced from the proof of \cref{thm:SS:LC:formal} below.
    \begin{theorem}[Strong Convexity in Neighbourhood of $W^\star$]\label{thm:SS:LC:formal}
    For $\rho\in(0,k^{-\Omega(B^2)})$, $t=\rho\cdot \wt{O}\inparen{\frac{1}{C^2B^4k^{17/2}}}$, and $V\in\R^{d\times k}$ satisfying $\norm{V}_F=1$, 
    define
    \[
        W_t\coloneqq W^\star+tV\,.
    \]
    Then, 
    \[
        \nabla^2 \cL(W_t) 
        \succeq 
        \nabla^2 \cL(W^\star)
            - (\rho + k^{-\Omega(B^2)}) I\,.
    \] 
    \end{theorem}
    First, we pause to observe that \cref{thm:SS:LC:formal} and \cref{eq:SS:LC:hessianLB} imply local convexity.
    To see this, note that any parameter $W$ that is $k^{-\Omega(B^2)}$ close to $W^\star=W_0$ in Frobenius norm can be expressed as $W=W_t$ for a suitable matrix $V$ (with $\norm{V}_2=1$), and constants $\rho=k^{-O(B^2)}$ and $t=\rho\cdot \wt{O}\inparen{{C^{-2}B^{-2}k^{-17/2}}}$. 
    

    In words, we get that at $W^\star$ the Hessian is at least $\poly(\nfrac{1}{k}) \cdot I$ and at any point $W$ which is $\poly(\nfrac{1}{k})$-close to $W^\star$, the Hessian decreases by some quantity of order at most $\poly(\nfrac{1}{k})$ and hence, by tuning the constants appropriately, remains {positive semi-definite}.
    
\bigskip 

\paragraph{Proof of Strong Convexity Near $W^\star$.} Next, we prove \cref{thm:SS:LC:formal}. 
Following \Cref{fig:outline:SS:LC}, we decompose the proof into the following two lemmas.

\begin{lemma}[Lower Bound on Sample NLL's Hessian]\label{lem:SS:LC:step1}
For $\zeta\in(0,1)$ and $t=\zeta\cdot O\inparen{\frac{1}{B^4k^3}}$, 
\[
    \forall_{x\in \R^d}\,,~\forall_{\ymax\in \R}\,,\quad 
    \nabla^2 \cL(W_t;x,\ymax) - \nabla^2 \cL(W_0;x,\ymax) \succeq -\gamma(x,\ymax) \cdot I_{k}\otimes xx^\top\,,
\]
where
\[ 
    \gamma(x,\ymax) \coloneqq 
        80 k^2\cdot
        \Bigl(\zeta + e^{-B^2 k/2}\Bigr)\cdot
        \Bigl(1 + \ymax^4 + \max\sinbrace{\|W_t^\top x\|_\infty, \|W_0^\top x\|_\infty}^4\Bigr)
        \,.
\]
\end{lemma}
The proof of \cref{lem:SS:LC:step1} is given in \cref{sec:proofof:lem:SS:LC:step1}.

\begin{lemma}[Bounds on $\Ex\insquare{\gamma}$]\label{lem:SS:LC:step2}
    For 
    $\zeta\in (0, 1)$, 
    $t=\zeta\cdot O\inparen{\frac{1}{B^4k^3}}$,
    and $\gamma$ as in \Cref{lem:SS:LC:step1},
    \[
        0\preceq \Ex_{x,\ymax}\Bigl[\gamma(x,\ymax) \cdot I_{k}\otimes xx^\top \Bigr] \preceq 
        C
            \cdot \wt{O}(k^{11/2})
            \cdot (\zeta + e^{-B^2 k/2})
            \cdot I_{dk}\,.
    \]
\end{lemma}
The proof of \cref{lem:SS:LC:step2} is given in \cref{sec:proofof:lem:SS:LC:step2}.

\medskip

Given these two lemmas, we are ready to prove the strong convexity property near $W^\star.$
\begin{proof}[Proof of \cref{thm:SS:LC:formal}]
    Fix $V$ with $\norm{V}_F=1$.
    We will prove the desired statement for $\rho\in (0,1)$
    but only need the smaller range to establish \cref{thm:SS:LC:formal}.
    Define
    \[
        t=\rho\cdot \wt{O}\inparen{\frac{1}{C^{2}B^{4}k^{17/2}}}
        \qquadand 
        \zeta=\rho\cdot \frac{1}{k^{11/2} C^{2}}\,.
        \yesnum\label{eq:SS:LC:defZetaVal}
    \]
    By construction,
    \[ 
        t = \zeta \cdot \wt{O}\inparen{\frac{1}{B^{4}k^{3}}}\,.
    \]
    Hence, \cref{lem:SS:LC:step1} is applicable with $(t,\zeta)$ and implies that, for each $x$, we have
    \[
    \Ex_{\ymax}\Bigl[\nabla^2 \cL(W_t;x,\ymax) - \nabla^2 \cL(W_0;x,\ymax)\Bigr] \succeq -\Ex_{\ymax}\bigl[\gamma(x,\ymax)\bigr]\, I_{k}\otimes xx^\top\,.
    \]
Averaging over $x$, we deduce
\[
\Ex_{x,\ymax}\Bigl[\nabla^2 \cL(W_t;x,\ymax) - \nabla^2 \cL(W_0;x,\ymax)\Bigr] \succeq -\Ex_{x,\ymax}\Bigl[\gamma(x,\ymax) \cdot I_{k}\otimes xx^\top\Bigr]\,.
\]
Since $\nabla^2 \cL(W)=\Ex_{x,\ymax}[\nabla^2 \cL(W;x,\ymax)]$ for any $W$, we obtain
\[
\nabla^2 \cL(W_t) - \nabla^2 \cL(W_0) \succeq -\Ex_{x,\ymax}\Bigl[\gamma(x,\ymax) \cdot I_{k}\otimes xx^\top\Bigr]\,.
\]
Substituting the upper bound on $\Ex_{x,\ymax}\Bigl[\gamma(x,\ymax) \cdot I_{k}\otimes xx^\top\Bigr]$ from \cref{lem:SS:LC:step2} implies, 
\[
    \nabla^2 \cL(W_t) - \nabla^2 \cL(W_0) \succeq -\wt{O}(Ck^{11/2})\cdot (\zeta+e^{-B^2k/2})\cdot I_{dk}\,.
\]
Substituting the value of $\rho$ from \cref{eq:SS:LC:defZetaVal}, we deduce that
\[
    \nabla^2 \cL(W_t) - \nabla^2 \cL(W_0) \succeq -(\rho+ \wt{O}(Ck^{11/2}) e^{-B^2k/2})\cdot I_{dk}\,.
\]
Hence, for a suitably large constant in the definition of $B$ (\cref{eq:SS:LC:defB}), we have,
\[
    \nabla^2 \cL(W_t) - \nabla^2 \cL(W_0) \succeq -(\rho + e^{-\Omega(B^2k)})\cdot I_{dk}\,.
\]
This is the required bound since $W_0=\Wstar$ (by construction).
\end{proof}



    \subsubsection{Proof of \cref{lem:SS:LC:step1} (Lower Bound on Sample NLL's Hessian)}\label{sec:proofof:lem:SS:LC:step1}

        We now move to the proof of \Cref{lem:SS:LC:step1}.
        To set up the stage, from the expressions of the Hessian (\cref{eq:SS:conditionalHess:mainBody}) one can verify that for each $(x,\ymax)$,
        \[
            \grad^2 \cL(W_t; x, \ymax) - \grad^2 \cL(W_0; x, \ymax) 
            =
            \inparen{
                \cov_{\normal{W_t^\top x}{I}}\insquare{z\mid z\in P_{\ymax}}
                -
                \cov_{\normal{W_0^\top x}{I}}\insquare{z\mid z\in P_{\ymax}}
            }
            \otimes xx^\top\,.
        \]
        Therefore, it suffices to upper bound the spectral norm of the following difference of covariances 
        \[
            \cov_{\normal{W_0^\top x}{I}}\insquare{z\mid z\in P_{\ymax}}
            -
            \cov_{\normal{W_t^\top x}{I}}\insquare{z\mid z\in P_{\ymax}}\,.
        \]
        This, in turn, is equivalent to upper bounding the following difference, for all unit vectors $v\in \R^k$
        \[
            \abs{
                \var_{\normal{W_t^\top x}{I}}\insquare{ \inangle{z, v} \mid z\in P_{\ymax}}
                -
                \var_{\normal{W_0^\top x}{I}}\insquare{ \inangle{z, v} \mid z\in P_{\ymax}}
            }\,.
            \yesnum\label{eq:SS:LC:diffOfVars}
        \]
            To simplify the notation, henceforth, for any $t\in\R$ and any function $f(\cdot)$, we use $\var\nolimits_{t}[f(z)]$ to denote $\var_{z\sim \normal{W_t^\top x}{I}}\insquare{ f(z) \mid z\in P_{\ymax}}$ and, similarly, use $\Ex\nolimits_{t}[ f(z) ]$ to denote $\Ex_{z\sim \normal{W_t^\top x}{I}}\insquare{ f(z) \mid z\in P_{\ymax}}$.
            
        \bigskip 

        \paragraph{Proof of \Cref{lem:SS:LC:step1}.}
        Using the previous decomposition, we will show the next claim, which directly implies \Cref{lem:SS:LC:step1}.
        \begin{claim}
            Define $\gamma$ as the following function,
            \[ 
                \gamma(x,\ymax) \coloneqq 
                    80 k^2\cdot
                    \Bigl(\zeta + e^{-B^2 k/2}\Bigr)\cdot
                    \Bigl(1 + \ymax^4 + \max\sinbrace{\|W_t^\top x\|_\infty, \|W_0^\top x\|_\infty}^4\Bigr)
                    \,.
            \]
            Fix any unit vector $z\in \R^d$, it holds that
            \[
            \abs{\var\nolimits_{t}\insquare{\inangle{v, z}} 
            -
            \var\nolimits_{0}\insquare{\inangle{v, z}}}
            \leq {\gamma(x,\ymax)}\,.
            \]
            Hence,
            \[
                 \nabla^2 \cL(W_t;x,\ymax) - \nabla^2 \cL(W_0;x,\ymax) \succeq -\gamma(x,\ymax) \cdot I_{k}\otimes xx^\top\,,
            \]
            \label{claim:variance}
        \end{claim}
       Following \Cref{fig:outline:SS:LC}, in order to prove this claim {we follow the following steps}: 
        \begin{enumerate}
            \item Introduce a high-probability event $\evE$ (left branch of the node \Cref{lem:SS:LC:step1} in \Cref{fig:outline:SS:LC}). This event controls the norms of all the random variables appearing in the Hessian.
            \item Conditioned on this event, we will derive the bound of \Cref{claim:variance} using some appropriate tail bounds (right branch of the node \Cref{lem:SS:LC:step1} in \Cref{fig:outline:SS:LC}).
        \end{enumerate} 


        \paragraph{Properties the Event $\evE$.} 
        \noindent Fix any unit vector $v\in \R^k$.
        We consider the following high-probability event, {which controls the norm {of} ${y}, W_0^\top x$ and $V^\top x$.}
        \begin{definition}
        [Good Event on Norm Bounds]
        \label{def:SS:LC:event}
            Define $\evE$ as the following event {over the draw of $(x,y)$ from the self-selection model where $y=(y_1,y_2,\dots,y_k)$ is the vector of all outcomes}\footnote{{Note that in the self-selection model, we only observe $(x,\ymax)$ where $\ymax=\max_{i}y_i$.}}
            \begin{align*}
                \snorm{W_0^\top x}_2\,,~\snorm{V^\top x}_2\,,\norm{y}_2 &~~\leq~~ B^2\cdot O\sinparen{k^{3/2}}\,.
            \end{align*}
            Here $B$ is the constant defined in \cref{eq:SS:LC:defB}.
        \end{definition}
        First, we verify that $\evE$ indeed happens with high probability; the proof appears in \cref{sec:proofof:lem:SS:LC:prob}.
        \begin{lemma}\label{lem:SS:LC:prob}
            $\Pr[\evE]\geq 1-e^{-B^2 k}$.  
        \end{lemma}
        Second, we show that $\evE$ enables us to perform the following change-of-measure; the proof appears in \cref{sec:proofof:lem:SS:LC:changeOfMeasure}. 
        This change-of-measure allows us to multiplicatively relate expectations of $W_t$ and $W_0$. {While, in general, these} expectations are not easily related but we show that they are {similar} conditioned on the good event $\evE.$
        \begin{lemma}[Change of Measure I]\label{lem:SS:LC:changeOfMeasure}
                For function $f\colon \R^k\to \R$ and $t\in [0,1]$,
                \[
                    e^{-O\inparen{t B^4 k^3}}
                    \leq 
                    \frac{
                        \Ex\nolimits_{t}\insquare{f(z)\given  \evE} 
                    }{
                        \Ex\nolimits_{0}\insquare{f(z)\given  \evE}
                    }\leq 
                    e^{O\inparen{t B^4 k^3}}\,.
                \]
        \end{lemma}
        This ability to perform a change-of-measure is the main reason why this event $\evE$ is useful.
        Observe that the following result is an immediate corollary of  \cref{lem:SS:LC:changeOfMeasure}.
        \begin{corollary}[Change of Measure II]\label{lem:SS:LC:changeOfMeasure:corr}
            For $\zeta\in [0,1)$ and $t=\zeta\cdot O\inparen{\frac{1}{B^4k^3}}$, it holds that,
            \[
                \Ex\nolimits_{t}\insquare{z\given \evE}
                \in 
                    \inparen{1\pm \zeta} 
                    \cdot 
                    \Ex\nolimits_{0}\insquare{z\given \evE}
                \quadand
                \Ex\nolimits_{t}\sinsquare{zz^\top \mid  \evE}
                \in 
                    \inparen{1\pm \zeta} 
                    \cdot 
                    \Ex\nolimits_{0}\sinsquare{zz^\top \mid \evE}
                \,.
            \]
        \end{corollary}
        \begin{proof}[Proof of \cref{lem:SS:LC:changeOfMeasure:corr}]
            Since for provided $t$, $O(tB^4k^3)=\zeta\leq 1$, and for any $r\in [0,1]$, $1-r\leq e^{-r}, e^r\leq 1+2r$, the result follows by applying  \cref{lem:SS:LC:changeOfMeasure} to each entry in $\Ex\nolimits_{t}\insquare{z\given \evE}$ and $\Ex\nolimits_{t}\insquare{zz^\top\given \evE}$.
        \end{proof}

        \paragraph{Bounding the Variance.} Given the above good event which allows us to control the norm bounds, we proceed with the proof of \Cref{claim:variance}.
        Now, we are ready to bound \cref{eq:SS:LC:diffOfVars}. 
        First, we recall the following expressions: for $t\geq 0$
        \begin{align*}
            \Ex\nolimits_{t}\sinsquare{\inangle{v, z}^2}
            &= \Pr[\evE]\cdot 
                \Ex\nolimits_{t}\sinsquare{\inangle{v, z}^2\mid \evE}
                ~+~ 
                \Pr[\lnot\evE]\cdot 
                \Ex\nolimits_{t}\sinsquare{\inangle{v, z}^2\mid \lnot\evE}\,,\\
            \Ex\nolimits_{t}\insquare{\inangle{v, z}}
            &= \Pr[\evE]\cdot 
                \Ex\nolimits_{t}\insquare{\inangle{v, z}\mid \evE}
                ~~+~ 
                \Pr[\lnot\evE]\cdot 
                \Ex\nolimits_{t}\insquare{\inangle{v, z}\mid \lnot\evE}\,.
        \end{align*}
        Hence,
        \begin{align*}
            &\eqref{eq:SS:LC:diffOfVars}
            ~=\abs{\var\nolimits_{t}\insquare{\inangle{v, z}} 
            -
            \var\nolimits_{0}\insquare{\inangle{v, z}}}\\
            &\qquad\leq 
            \Pr[\evE]~~\abs{ 
                \Ex\nolimits_{t}\sinsquare{\inangle{v, z}^2\mid \evE}
                - 
                \Ex\nolimits_{0}\sinsquare{\inangle{v, z}^2\mid \evE}
            }\yesnum\label{eq:SS:LC:ub1}\\
            &\qquad + 
            \Pr[\lnot\evE]~~\inparen{ 
                \Ex\nolimits_{t}\sinsquare{\inangle{v, z}^2\mid \lnot\evE}
                - 
                \Ex\nolimits_{0}\sinsquare{\inangle{v, z}^2\mid \lnot\evE}
            }\yesnum\label{eq:SS:LC:ub2}\\
            &\qquad +\Pr[\evE]^2~~\abs{
                {\Ex\nolimits_{t}\insquare{\inangle{v, z}\mid \evE}}^2
                - 
                {\Ex\nolimits_{0}\insquare{\inangle{v, z}\mid \evE}}^2
            }\yesnum\label{eq:SS:LC:ub3}\\
            &\qquad +\Pr[\lnot\evE]^2~~\inparen{
                {\Ex\nolimits_{t}\insquare{\inangle{v, z}\mid \lnot\evE}}^2
                + %
                {\Ex\nolimits_{0}\insquare{\inangle{v, z}\mid \lnot\evE}}^2
            }\yesnum\label{eq:SS:LC:ub4}\\
            &\qquad +2\Pr[\evE]\Pr[\lnot\evE]~~\inparen{
                \Ex\nolimits_{t}\insquare{\inangle{v, z}\mid \evE}
                \Ex\nolimits_{t}\insquare{\inangle{v, z}\mid \lnot\evE}
                + %
                \Ex\nolimits_{0}\insquare{\inangle{v, z}\mid \evE}
                \Ex\nolimits_{0}\insquare{\inangle{v, z}\mid \lnot\evE}
            }\,.\yesnum\label{eq:SS:LC:ub5}
        \end{align*}

        \paragraph{Tail Bounds.} 
        {The last step to prove \Cref{claim:variance} is to bound all the above terms. This reduces to obtaining {suitable} tail bounds {related to} these terms.}
        Our next result upper bounds each of the above terms by quantities that scale with $\zeta$ or $e^{-B^2k/2}$.
        \begin{lemma}[Tail Bound I]
            \label{lem:SS:LC:UB15}
            \label{lem:SS:LC:tailBoundI}
            For $\zeta\in [0,1)$ and $t\leq \zeta \cdot O\inparen{\frac{1}{B^4k^3}}$,
           \begin{align*}
               \abs{\eqref{eq:SS:LC:ub1}}
                    &~~\leq~~ \zeta
                    \cdot \sqrt{\Ex\nolimits_{0}\sinsquare{\inangle{v, z}^4}}\,,\\
                \abs{\eqref{eq:SS:LC:ub2}}
                    &~~\leq~~ 
                    e^{-B^2 k/2}\cdot
                    \inparen{ 
                        \sqrt{\Ex\nolimits_{t}\sinsquare{\inangle{v, z}^4}}
                        +
                        \sqrt{\Ex\nolimits_{0}\sinsquare{\inangle{v, z}^4}}
                    }\,,\\[-1mm]
                \abs{\eqref{eq:SS:LC:ub3}}
                    &~~\leq~~ \zeta
                    \cdot \inparen{\Ex\nolimits_{0}\sinsquare{{\inangle{v, z}}^2}}^2\,,\\
                \abs{\eqref{eq:SS:LC:ub4}}
                    &~~\leq ~~
                    e^{-B^2 k}\cdot
                    \inparen{
                        {\Ex\nolimits_{t}\sinsquare{{\inangle{v, z}}^4}}
                        +
                        {\Ex\nolimits_{0}\sinsquare{{\inangle{v, z}}^4}}
                    }\,,\\
                \abs{\eqref{eq:SS:LC:ub5}}
                    &~~\leq~~ 2e^{-B^2 k/2}\cdot
                    \inparen{
                        {\Ex\nolimits_{t}\sinsquare{\inangle{v, z}^2}}
                        +
                        {\Ex\nolimits_{0}\sinsquare{\inangle{v, z}^2}}
                    }\,.
           \end{align*} 
        \end{lemma}
        To prove \cref{lem:SS:LC:UB15}, intuitively, we control the contributions from \cref{eq:SS:LC:ub1,eq:SS:LC:ub3} via the change-of-measure bound in \cref{lem:SS:LC:changeOfMeasure:corr}, while the remaining terms (namely, \cref{eq:SS:LC:ub2,eq:SS:LC:ub4,eq:SS:LC:ub5}) are bounded using the tail bound $\Pr[\lnot\evE]\le e^{-B^2 k}$ from \cref{lem:SS:LC:prob}.
        The proof of \cref{lem:SS:LC:UB15} appears in \cref{sec:proofof:lem:SS:LC:UB15}.
        
        To use the above bounds, we {further} need to bound expectations of the form $\Ex\nolimits_{t}\sinsquare{\inangle{v,z}^\ell}$ for $\ell\in\{2,4\}$.
        Our next result proves these bounds and its proof appears \cref{sec:proofof:lem:SS:LC:unconditionalMoments}.
        \begin{lemma}[Tail Bound II]
            \label{lem:SS:LC:unconditionalMoments}
            \label{lem:SS:LC:tailBoundII}
            It holds that 
            \begin{align*}
                 \Ex\nolimits_{t}\sinsquare{\inangle{v,z}^2\mid z\in P} &\leq 
                 2k\inparen{1 + \abs{\ymax} + \snorm{W_t^\top x}_\infty}^2
                 ~\text{and}~
                 \Ex\nolimits_{t}\sinsquare{\inangle{v,z}^4\mid z\in P} &\leq 
                 3k^2\inparen{1 + \abs{\ymax} + \snorm{W_t^\top x}_\infty}^4\,.
            \end{align*}
            Substituting $t=0$, gives upper bounds on $\Ex\nolimits_{0}\sinsquare{\inangle{v,z}^2}$ and $\Ex\nolimits_{0}\sinsquare{\inangle{v,z}^4}$.
        \end{lemma}

        \paragraph{Completing the Proof of \Cref{claim:variance}.}
        Combining \cref{lem:SS:LC:UB15,lem:SS:LC:unconditionalMoments} and substituting them in \cref{eq:SS:LC:diffOfVars} implies 
        \begin{align*}
            \abs{\var\nolimits_{t}\insquare{\inangle{v, z}} 
            -
            \var\nolimits_{0}\insquare{\inangle{v, z}}}
            \leq 20 k^2\cdot
            \Bigl(\zeta + e^{-B^2 k/2}\Bigr)\cdot
            \Bigl(1 + |\ymax| + \max\sinbrace{\|W_t^\top x\|_\infty, \|W_0^\top x\|_\infty}\Bigr)^4\,.
        \end{align*}
        Since $(a+b+c)^4 \leq 4(a^4+b^4+c^4)$, we have the following 
        \begin{align*}
            \abs{\var\nolimits_{t}\insquare{\inangle{v, z}} 
            -
            \var\nolimits_{0}\insquare{\inangle{v, z}}}
            \leq 80 k^2\cdot
            \Bigl(\zeta + e^{-B^2 k/2}\Bigr)\cdot
            \Bigl(1 + \ymax^4 + \max\sinbrace{\|W_t^\top x\|_\infty, \|W_0^\top x\|_\infty}^4\Bigr)\,.
        \end{align*}
        {This completes the proof of \Cref{claim:variance}.}
        
        
        





    \subsubsection{Proof of \cref{lem:SS:LC:step2} (Bounds on $\Ex[\gamma]$)}\label{sec:proofof:lem:SS:LC:step2}

        Following \Cref{fig:outline:SS:LC}, our last ingredient for the proof of the strong convexity around $W^\star$ is a bound on the expected value of the function $\gamma$, introduced in \Cref{lem:SS:LC:step1}.
        Recall that 
        \[
            \gamma(x,\ymax) \coloneqq 
                80 k^2\cdot
                \Bigl(\zeta + e^{-B^2 k/2}\Bigr)\cdot
                \Bigl(1 + \ymax^4 + \max\sinbrace{\|W_t^\top x\|_\infty, \|W_0^\top x\|_\infty}^4\Bigr)\,.
            \yesnum\label{eq:SS:LC:defGamma2}
        \]
        The lower bound $\Ex_{x,\ymax}\Bigl[\gamma(x,\ymax) \cdot I_{k}\otimes xx^\top \Bigr]\succeq 0$ is simple: it follows because $\gamma(x,\ymax)\geq 0$ and the other remaining terms are positive semi-definite.
        
        The proof of the upper bound has two steps: 
            the first step bounds $\Ex_{\ymax}{
            \gamma\inparen{x,\ymax}
            }$ and, then, we use this to bound $\Ex_{x,\ymax}\insquare{
            \gamma\inparen{x,\ymax}\cdot I_{k}\otimes xx^\top
            }$.
        We use the following fact, which we proved at the end of this section.
        \begin{fact}[Expectations of Powers of $L_\infty$-norms of Gaussian Vectors]\label{fact:GaussMaxMoments}
            For $1\leq \ell \leq 6$ and $\mu\in \R^k$,
            \[
                \Ex_{v\sim \normal{\mu}{I}}{\norm{v}_\infty^\ell} \leq 
                    12 + 384\cdot \inparen{10+\log{k}}^{\ell/2}
                    + \ell \norm{\mu}_\infty^\ell\,.
            \]
        \end{fact}
        
    \paragraph{Step 1 (Upper bound $\Ex_{\ymax}{
            \gamma\inparen{x,\ymax}
        }$):}~~
        Fix any $x\in \R^d$.
        By linearity of expectation,
        \begin{align*}
            \Ex_{\ymax} \gamma(x,\ymax) 
            = 
                80 k^2\cdot
                \Bigl(\zeta + e^{-B^2 k/2}\Bigr)\cdot
                \Bigl(1 + \Ex_{\ymax}[\ymax^4] + \max\sinbrace{\|W_t^\top x\|_\infty, \|W_0^\top x\|_\infty}^4\Bigr)
                \,.
        \end{align*}
        Since $\ymax = \norm{y}_\infty$ for $y\sim \normal{W_0^\top x}{I_{k}}$, \cref{fact:GaussMaxMoments} implies
        \begin{align*}
            \Ex_{\ymax}[\ymax^4]
            \leq 12 + 384 \cdot \inparen{10+\log{k}}^2 + 2\snorm{W_0^\top x}_\infty^4
            = O\sinparen{1+\log{k}}^2 + 2\snorm{W_0^\top x}_\infty^4
            \,.
        \end{align*}
        Substituting this in the previous expression implies that
        \begin{align*}
            \Ex_{\ymax} \gamma(x,\ymax) 
            \leq 
                O(k^2)\cdot
                (\zeta + e^{-B^2 k/2})\cdot
                \Bigl(1 + \log^2{k} + \max\sinbrace{\|W_t^\top x\|_\infty, \|W_0^\top x\|_\infty}^4\Bigr)
                \,.
                \yesnum\label{eq:SS:LC:step2}
        \end{align*} 
    
    \paragraph{Step 2 (Bound $\Ex_{x,\ymax}\insquare{\gamma\inparen{x,\ymax}\cdot I_{k}\otimes xx^\top }$):}~~
    Since $\normal{0}{I_{d\times d}}$ is rotation invariant and symmetric on each coordinate, we can rotate the space and re-arrange the coordinates without changing the distribution of $x$.
    Perform a rotation and re-arrangement so that the columns of $W_t$ and $W_0$ lie in the span of the first $2k$ standard basis vectors $e_1,\dots,e_{2k}$.
    Let $x_{\leq 2k}$ denote the first $2k$ coordinates of $x$ and $x_{>2k}$ denote the last $2k$ coordinates of $x$.
    

    It holds that 
    \[
        \Ex_{x,\ymax}\insquare{
            \gamma\inparen{x,\ymax}\cdot I_{k}\otimes xx^\top
            }
        ~~=~~
        \Ex_{x_{\leq 2k}}
        \inparen{
            \Ex_{x_{> 2k}}\insquare{
                \Ex_{\ymax}{
                    \gamma\inparen{x,\ymax}
                }\cdot I_{k}\otimes 
                \begin{bmatrix}
                    x_{\leq 2k}x_{\leq 2k}^\top & x_{\leq 2k}x_{>2k}^\top\\
                    x_{>2k}x_{\leq 2k}^\top & x_{>2k} x_{>2k}^\top
                \end{bmatrix}
            }
        }\,.
    \]
    Next, observe that $\Ex_{\ymax}{
                    \gamma\inparen{x,\ymax}
                }$ is a constant independent of $x_{>2k}$. 
    To see this, observe that \cref{eq:SS:LC:defGamma2} tells us that $\Ex_{\ymax}{\gamma\inparen{x,\ymax}}$ is only a function of $\Ex_{\ymax}[\ymax]$ and the inner products $W_t^\top x,W_0^\top x$ -- and further, after the above transformation, $\Ex_{\ymax}[\ymax]$ and $W_t^\top x,W_0^\top x$ are only functions of $x_{\leq 2k}$, which is independent of $x_{>2k}$ as $x\sim \normal{0}{I_d}$.
    Hence,
    \begin{align*}
        \Ex_{x,\ymax}\insquare{
            \gamma\inparen{x,\ymax}\cdot I_{k}\otimes xx^\top
            }
        ~~=~~
        \Ex_{x_{\leq 2k}}
        \inparen{
            \Ex_{\ymax}{
                \gamma\inparen{x,\ymax}
            }\cdot I_{k}\otimes 
            \Ex_{x_{> 2k}}\insquare{
                \begin{bmatrix}
                    x_{\leq 2k}x_{\leq 2k}^\top & x_{\leq 2k}x_{>2k}^\top\\
                    x_{>2k}x_{\leq 2k}^\top & x_{>2k} x_{>2k}^\top
                \end{bmatrix}
            }
        }\,.
    \end{align*}
    Since $x_{\leq 2k}$ and $x_{>2k}$ are independent and have zero mean, and $x_{>2k}\sim \normal{0}{I_{d-2k}}$, it follows that,
    \[
        \Ex_{x,\ymax}\insquare{
            \gamma\inparen{x,\ymax}\cdot I_{k}\otimes xx^\top
            }
        ~~=~~
        \Ex_{x_{\leq 2k}}
        \inparen{
            {
                \Ex_{\ymax}{
                    \gamma\inparen{x,\ymax}
                }\cdot I_{k}\otimes 
                \begin{bmatrix}
                    x_{\leq 2k}x_{\leq 2k}^\top & 0\\
                    0 & I_{d-2k}
                \end{bmatrix}
            }
        }\,.
    \]
    Which can further be rewritten as 
    \[
        \Ex_{x,\ymax}\insquare{
            \gamma\inparen{x,\ymax}\cdot I_{k}\otimes xx^\top
            }
        ~~=~~
        I_{k}\otimes 
        \begin{bmatrix}
            \Ex\limits_{x_{\leq 2k},\ymax}{
            \gamma\inparen{x,\ymax}
        }\cdot  x_{\leq 2k}x_{\leq 2k}^\top & 0\\
            0 & \Ex\limits_{x_{\leq 2k},\ymax}{
            \gamma\inparen{x,\ymax}
        }\cdot  I_{d-2k}
        \end{bmatrix} 
        \,.
        \yesnum\label{eq:SS:LC:part2:spectralMatrix}
    \]
    Next, we upper bound the two non-zero terms.
    We use the following bounds:
    \begin{align*} 
             \snorm{W_0^\top x}_\infty 
        &~~~\Stackrel{\substack{W_0=[w_i^\star]_i\\[1.5mm]}}{\leq}~~~
            \sum\nolimits_i \abs{\inangle{w_i,x_{\leq 2k}}}
        ~\leq~
            \sum\nolimits_i \norm{w^\star_i}_2 \norm{x_{\leq 2k}}_2
        \quad\Stackrel{\substack{\rm\cref{asmp:SS:separationBoundedness}\\[1.5mm]}}{\leq}\quad 
            kC\norm{x_{\leq 2k}}_2
        ~\leq~ k^{3/2}C\norm{x_{\leq 2k}}_\infty\,, \yesnum\label{eq:SS:LC:simplebounds1}\\
        \snorm{W_t^\top x}_\infty 
        &~~~\Stackrel{\substack{W_0=[w_i^\star+tv_i]_i\\[1.5mm]}}{\leq}~~~
            \sum\nolimits_i \abs{\inangle{w_i+tv_i,x_{\leq 2k}}}
        ~\leq~ 
        \sum\nolimits_i \norm{w^\star_i + t v_i}_2 \norm{x_{\leq 2k}}_2\\
        &\qquad\qquad\qquad\qquad\qquad\quad\hspace{0.5mm}
            ~~~\Stackrel{\substack{\rm\cref{asmp:SS:separationBoundedness},\\\norm{V}_F, t\leq 1}}{\leq}\quad 
        k(C+1)\norm{x_{\leq 2k}}_2
        ~\Stackrel{(C\geq 1)}{\leq}~
            2k^{3/2}C \norm{x_{\leq 2k}}_\infty \,. 
        \yesnum\label{eq:SS:LC:simplebounds2}
    \end{align*}
    Substituting these bounds in \cref{eq:SS:LC:step2} implies that 
    \begin{align*}
        \Ex_{x_{\leq 2k}}\Ex_{\ymax}{\gamma\inparen{x,\ymax}}
        &\leq 
        O(k^2)\cdot
            (\zeta + e^{-B^2 k/2})\cdot
            \Bigl(1 + \log^2{k} + 2k^{3/2}C \Ex_{x_{\leq 2k}}\norm{x_{\leq 2k}}_\infty \Bigr)\,.
    \end{align*}
    Next, the upper bound on the $\ell_\infty$-norm of a Gaussian vector from \cref{fact:GaussMaxMoments} implies that 
    \begin{align*}
        \Ex_{x_{\leq 2k}}\Ex_{\ymax}{\gamma\inparen{x,\ymax}}
        &\leq 
        O(k^2)\cdot
                (\zeta + e^{-B^2 k/2})\cdot
            \Bigl(
                1 
                + \log^2{k} 
                + 2k^{3/2}C 
                    \inparen{12+384\cdot\sqrt{10+\log{k}}}
            \Bigr)\\
        &= C\cdot 
            O(k^{5/2}\sqrt{\log{k}}) \cdot
            (\zeta + e^{-B^2 k/2})
        \,.
        \yesnum\label{eq:SS:LC:part2:spectral1}
    \end{align*}
    In the remainder of the proof, we upper bound $\norm{\Ex_{x_{\leq 2k}, \ymax}{
                \gamma\inparen{x,\ymax}
            }\cdot  x_{\leq 2k}x_{\leq 2k}^\top}_2$.
    Since for any matrix $M\in \R^{2k\times 2k}$, $\norm{M}_2\leq (2k)^2\norm{M}_\infty$, it follows that 
    \begin{align*}
        \norm{\Ex_{x_{\leq 2k}}\Ex_{\ymax}{
                \gamma\inparen{x,\ymax}
            }\cdot  x_{\leq 2k}x_{\leq 2k}^\top}_2
        &\leq 4k^2 \norm{\Ex_{x_{\leq 2k}}\Ex_{\ymax}{
                \gamma\inparen{x,\ymax}
            }\cdot  x_{\leq 2k}x_{\leq 2k}^\top}_\infty\,.
    \end{align*} 
    Further, by the triangle inequality, we have that 
    \begin{align*}
        \norm{\Ex_{x_{\leq 2k}}\Ex_{\ymax}{
                \gamma\inparen{x,\ymax}
            }\cdot  x_{\leq 2k}x_{\leq 2k}^\top}_2
        \leq 4k^2 \Ex_{x_{\leq 2k}}\insquare{\abs{\Ex_{\ymax}{
                \gamma\inparen{x,\ymax}
            }} \cdot \norm{x_{\leq 2k}}_\infty^2}\,.
    \end{align*}
    Substituting bounds from \cref{eq:SS:LC:simplebounds1,eq:SS:LC:simplebounds2} and using linearity of expectation implies, 
    \[
        \norm{\Ex_{x_{\leq 2k}}\Ex_{\ymax}{
                \gamma\inparen{x,\ymax}
            }\cdot  x_{\leq 2k}x_{\leq 2k}^\top}_2
        \leq 
        O(k^4)\cdot
            (\zeta + e^{-B^2 k/2})\cdot
            \Bigl(
                (1 + \log^2{k})\norm{x_{\leq 2k}}_\infty^2
                + 
                2k^{3/2}C \norm{x_{\leq 2k}}_\infty^3 
            \Bigr)\,.
    \]
    Next, the upper bound on the $\ell_\infty$-norm of a Gaussian vector from \cref{fact:GaussMaxMoments} implies that 
    \begin{align*}
        \norm{\Ex_{x_{\leq 2k}}\Ex_{\ymax}{
                \gamma\inparen{x,\ymax}
            }\cdot  x_{\leq 2k}x_{\leq 2k}^\top}_2
        &\leq 
        O(k^4)\cdot
            (\zeta + e^{-B^2 k/2})\cdot
            \Bigl(
                1
                +\log^3{k}
                +2k^{3/2}C \log^{3/2}{k}
            \Bigr)\\
        &=C
        \cdot O(k^{11/2}\log^{3/2}{k})
        \cdot (\zeta + e^{-B^2 k/2})\,.
        \yesnum\label{eq:SS:LC:part2:spectral2}
    \end{align*} 
    Therefore, substituting the bounds from \cref{eq:SS:LC:part2:spectral1,eq:SS:LC:part2:spectral2} into \cref{eq:SS:LC:part2:spectralMatrix} implies that 
    \[
        \Ex_{x,\ymax}\insquare{
            \gamma\inparen{x,\ymax}\cdot I_{k}\otimes xx^\top
            }
        ~~\preceq~~
        C
        \cdot O(k^{11/2}\log^{3/2}{k})
        \cdot (\zeta + e^{-B^2 k/2})
        \cdot 
        I_{dk}\,.
    \]
    

    \subsubsection*{Proof of \cref{fact:GaussMaxMoments} (Expectations of Powers of $\ell_\infty$-Norms of Gaussian Vectors)}
    \begin{proof}[Proof of \cref{fact:GaussMaxMoments}]
        The triangle inequality and the fact that $(a+b)^\ell\leq \ell(a^\ell+b^\ell)$ for $a,b\geq 0$ implies
        \[
            \Ex_{v}\norm{v}_\infty^\ell 
            \leq 
            \ell \Ex_{v}\norm{v-\mu}_\infty^\ell 
            + 
            \ell\norm{\mu}_\infty^\ell\,.
        \]
        Define $M\coloneqq \norm{v-\mu}_\infty$.
        It remains to upper bound $ \Ex_{v}[M^\ell]$.
        Observe that,
        \begin{align*}
            \Ex[M^\ell] 
            = \int_0^\infty \Pr[M^\ell > w] \d w 
            = \int_0^\infty \ell r^{\ell-1}\Pr[M > r] \d r\,.
        \end{align*}
        We divide the integral into two parts:
        \[
            \int_{} \ell r^{\ell-1}\Pr[M > r]\d r
            = 
                \int_{0}^{2\sqrt{10+\log{k}}} \ell r^{\ell-1}\Pr[M > r]\d r
                +
                \int_{2\sqrt{10+\log{k}}}^\infty \ell r^{\ell-1}\Pr[M > r]\d r\,.
                \yesnum\label{eq:SS:LC:GaussMaxMoments:twoTermSum}
        \]
        Since $\Pr[M > r]\leq 1$, the first term satisfies the following upper bound 
        \[
            \int_{0}^{2\sqrt{10+\log{k}}} \ell r^{\ell-1}\Pr[M > r]\d r
            \leq \ell \cdot 2^{\ell} \cdot \inparen{10+\log{k}}^{\ell/2}\,.
        \]
        Next observe that, for any $r\geq 0$,
        \[
            \Pr[M > r]
            \leq k\cdot \Pr_{z\sim \cN(0,1)}[z>r] 
            \Stackrel{}{\leq}
                \frac{2k}{r}\cdot e^{-r^2/2}\,.
                \yesnum\label{eq:SS:LC:tailBound}
        \]
        Where the last inequality uses a standard Gaussian tail bound, see \cref{fact:GaussianTail}.
        Therefore, the second term satisfies the following upper bound
        \[
            \int_{2\sqrt{10+\log{k}}}^\infty \ell r^{\ell-1}\Pr[M > r]\d r
            ~~\Stackrel{\eqref{eq:SS:LC:tailBound}}{\leq }~~
            2k \int_{2\sqrt{10+\log{k}}}^\infty \ell r^{\ell-2} e^{-r^2/2}\d r\,.
        \]
        Since $\ell\in [0,6]$, $2\sqrt{10+\log{k}}\geq 6$, $\int_{6}^\infty r^{2\ell-4} e^{-r^2/2}\d r\leq 1$,\footnote{For $z\ge6$ and $0\le x\le6$, we have $z^{2x-4}\le z^8$, so it suffices to show that $\int_{6}^{\infty}z^8e^{-z^2/2}\d z\leq 1.$
        Substituting $u=z^2/2$, transforms the integral into $2^{7/2}\int_{18}^\infty u^{7/2} e^{-u} \d u = 2^{7/2}\Gamma(9/2,18)\approx 0.005 \le 1$; where $\Gamma(x,s)\coloneqq \int_{x}^\infty u^{s-1} e^{-u}\d u$.} and $\int_0^\infty e^{-r^2/2} \d r\leq 2$,
        the Cauchy--Schwarz inequality implies 
        \begin{align*}
            \int_{2\sqrt{10+\log{k}}}^\infty \ell r^{\ell-2} e^{-r^2/2}\d r
            \leq 
            \sqrt{\int_{2\sqrt{10+\log{k}}}^\infty \ell^2 r^{2\ell-4} e^{-
            r^2/2}\d r}
            \cdot 
            \sqrt{\int_{0}^\infty e^{-r^2/2}\d r}
            \leq 2\ell\,.
        \end{align*}
        Substituting the bounds on the first and second term of \cref{eq:SS:LC:GaussMaxMoments:twoTermSum} into \cref{eq:SS:LC:GaussMaxMoments:twoTermSum} implies,
        \[
            \Ex[M^\ell]
            \leq 
            \ell 2^\ell \cdot \inparen{10+\log{k}}^{\ell/2}
            + 2\ell
            ~~~~\Stackrel{(0\leq \ell\leq 6)}{\leq}~~~~
                12 + 384\cdot \inparen{10+\log{k}}^{\ell/2}
            \,.
        \]
    \end{proof}
    
\subsubsection{Proof of \cref{lem:SS:LC:prob} (High Probability Guarantee)}
    \label{sec:proofof:lem:SS:LC:prob}
    \begin{proof}[Proof of \cref{lem:SS:LC:prob}]
        Consider a $d\times k$ matrix 
        \[
        M=\begin{bmatrix} m_1 & m_2 & \dots & m_k \end{bmatrix}\,.
        \]
        Rotate the coordinate system (and re-order if necessary) so that $m_1,\dots,m_k$ lie in the span of the first $k$ standard basis vectors $e_1,\dots,e_k$. 
        Since $\cN(0,I)$ is rotation invariant, $x\sim\normal{0}{I}$ remains unchanged. 
        Consequently, $M^\top x$ depends only on the first $k$ entries of $x$, \ie{}, 
        \[
        x_{1:k}=\inparen{x_1,x_2,\dots,x_k}\,.
        \] 
        Since $\norm{M}_2\leq \norm{M}_F$,
        \[
            \snorm{M^\top x}_2
            \leq  \norm{M}_2\norm{x_{1:k}}_2
            \leq  \norm{M}_F\norm{x_{1:k}}_2\,.
        \]
        Standard concentration for the norm of a $k$-dimensional standard Gaussian shows that, with probability at least $1-(\nfrac{\delta}{3})$,
        \[
        \norm{x_{1:k}}_2 \le \sqrt{k} + O\inparen{\sqrt{\log\nfrac{1}{\delta}}}\,.
        \]
        Now, set $M=W_0$ or $M=V$, and note that by construction $\norm{V}_F=1$ and by \cref{asmp:SS:separationBoundedness} each $\norm{w_i^\star}\le C$. Hence, with probability at least $1-(\nfrac{\delta}{3})$,
        \begin{align*}
            \snorm{V^\top x}_2 \le \sqrt{k} + O\inparen{\sqrt{\log\nfrac{1}{\delta}}}
                \quadand
            \snorm{W_0^\top x}_2 \le kC\inparen{\sqrt{k} + O\inparen{\sqrt{\log\nfrac{1}{\delta}}}}\,.
        \yesnum\label{eq:SS:LC:eventWHP:1}
        \end{align*}
        Since $W_t=W_0+tV$ with $t\le 1$, we obtain
        \[
            \snorm{W_t^\top x}_2 
            ~~\le~~ \snorm{W_0^\top x}_2 + t\,\snorm{V^\top x}_2 
            ~~\Stackrel{(C\geq 1)}{\le}~~ 
                2kC\inparen{\sqrt{k} + O\inparen{\sqrt{\log\nfrac{1}{\delta}}}}\,.
        \yesnum\label{eq:SS:LC:eventWHP:2}
        \]
        Moreover, since 
        \[
        z\sim \cN(W_t^\top x, I_{k})\,,
        \]
        then by Gaussian concentration, with probability at least $1-(\nfrac{\delta}{3})$,
        \[
            \snorm{z-W_t^\top x}_2 \le \sqrt{k} + O\Bigl(\sqrt{\log\nfrac{1}{\delta}}\Bigr)\,.
            \yesnum\label{eq:SS:LC:eventWHP:event2}
        \]
        Thus, combining this with \cref{eq:SS:LC:eventWHP:2}, if the events in \cref{eq:SS:LC:eventWHP:1,eq:SS:LC:eventWHP:event2} hold, then, 
        \[
        \norm{z}_2 
            ~~\Stackrel{(C\geq 1)}{\le}~~ 3kC\Bigl(\sqrt{k} + O\Bigl(\sqrt{\log\nfrac{1}{\delta}}\Bigr)\Bigr)\,.
        \]
        Taking a union bound over the events and setting $\delta=e^{-B^2k}$ completes the proof. 
    \end{proof} 




\subsubsection{Proof of \cref{lem:SS:LC:changeOfMeasure} (Change of Measure)}  
\label{sec:proofof:lem:SS:LC:changeOfMeasure}
    \begin{proof}[Proof of \cref{lem:SS:LC:changeOfMeasure}]
        Define 
        \[
        \mu_t \coloneqq W_t^\top x\,.
        \]
        Then, for any $y$, 
        \begin{align}
        \frac{\normalMass{\mu_t}{I}{y}}{\normalMass{\mu_0}{I}{y}}
          &= \exp\Bigl(\frac{\|\mu_0-y\|_2^2 - \|\mu_t-y\|_2^2}{2}\Bigr)
          \,,
          \label{eq:SS:LC:changeOfMeasure:ratio}
        \end{align} 
        By construction, 
        \[
        \mu_t = W_0^\top x + t\,V^\top x\,,
        \]
        and, conditioned on the event $\evE$, 
        \[
        \snorm{V^\top x}_2 \le B^2\cdot O\sinparen{k^{3/2}}\,.
        \]
        Hence, we have
        \[
        \norm{\mu_t-y} \in \norm{\mu_0-y} \pm \norm{\mu_t-\mu_0} \in \norm{\mu_0-y} \pm t\,B^2\cdot O\sinparen{k^{3/2}}\,.
        \]
        Furthermore, under $\evE$, 
        \[
        \norm{\mu_0}_2,~~~ \norm{y} ~~~\le~~~ B^2\cdot O\sinparen{k^{3/2}}\,.
        \]
        It follows that
        \[
        \norm{\mu_t-y}^2 \in \norm{\mu_0-y}^2 
        \pm 4t\,B^4\cdot O(k^3)
        \pm t^2\,B^4\cdot O(k^3)
        \overset{(t\le 1)}{\in} \norm{\mu_0-y}^2 \pm 5t\,B^4\cdot O(k^3)\,.
        \]
        Substituting this bound into \cref{eq:SS:LC:changeOfMeasure:ratio} gives
        \begin{align}
        \exp\inparen{-\frac{5t\,B^4}{2}\cdot O(k^3)}
        \le \frac{\normalMass{\mu_t}{I}{y}}{\normalMass{\mu_0}{I}{y}}
        \le \exp\inparen{\frac{5t\,B^4}{2}\cdot O(k^3)}\,.
        \label{eq:SS:LC:changeOfMeasure:Ratio2}
        \end{align} 
        Finally, for any function $f(y)$, using \cref{eq:SS:LC:changeOfMeasure:Ratio2} we obtain
        \begin{align*}
        \Ex\nolimits_t\insquare{f(y)}
          &= \frac{\int_{S(\ymax)} f(y)\,\normalMass{\mu_t}{I}{y}\,\d y}{\normalMass{\mu_t}{I}{S(\ymax)}}\\[1mm]
          &\in e^{5t\,B^4\cdot O(k^3)}\cdot \frac{\int_{S(\ymax)} f(y)\,\normalMass{\mu_0}{I}{y}\,\d y}{\normalMass{\mu_0}{I}{S(\ymax)}}\\
          &= e^{5t\,B^4\cdot O(k^3)}\cdot \Ex\nolimits_0\insquare{f(y)}\,.
        \end{align*}
        This completes the change-of-measure argument.
    \end{proof} 
     

    
    \subsubsection{Proof of \cref{lem:SS:LC:UB15} (Tail Bounds on Conditional Expectations)}
        \label{sec:proofof:lem:SS:LC:UB15}
    \begin{proof}[Proof of \cref{lem:SS:LC:UB15}]
            
        We divide the proof into two parts. 
        In the first part, we upper bound \cref{eq:SS:LC:ub1,eq:SS:LC:ub3}; in the second part, we control \cref{eq:SS:LC:ub2,eq:SS:LC:ub4,eq:SS:LC:ub5}.
        
    \paragraphit{Part A (Bounding terms \cref{eq:SS:LC:ub1,eq:SS:LC:ub3}):}~~
            By \cref{lem:SS:LC:changeOfMeasure:corr}, we have
            \begin{align*}
                \abs{\eqref{eq:SS:LC:ub1}}
                &~~\leq~~
                \zeta \cdot \Pr[\evE]
                \cdot \Ex\nolimits_{0}\sinsquare{\inangle{v, z}^2\mid \evE}
                \qquadand
                \abs{\eqref{eq:SS:LC:ub3}}
                ~~\leq~~
                \zeta \cdot \Pr[\evE]^2
                \cdot \inparen{\Ex\nolimits_{0}\insquare{\inangle{v, z}\mid \evE}}^2\,.
            \end{align*}
            Using the Cauchy--Schwarz inequality, we can ``remove'' the conditioning on the expectations 
           \begin{align*}
                \abs{\eqref{eq:SS:LC:ub1}}
                &~~\leq~~
                \zeta \cdot \sqrt{\Pr[\evE] \cdot \Ex\nolimits_{0}\insquare{\inangle{v, z}^4}}
                \qquadand
                \abs{\eqref{eq:SS:LC:ub3}}
                ~~\leq~~
                \zeta \cdot \Pr[\evE]
                \cdot {\Ex\nolimits_{0}\insquare{{\inangle{v, z}}^2}}\,.
           \end{align*}
           Since $\Pr[\evE]\le 1$, the upper bounds further simplify to 
           \begin{align*}
                \abs{\eqref{eq:SS:LC:ub1}}
                &~~\leq~~
                \zeta \cdot 
                    \sqrt{\Ex\nolimits_{0}\insquare{\inangle{v, z}^4}}
                \qquadand
                \abs{\eqref{eq:SS:LC:ub3}}
                ~~\leq~~
                \zeta \cdot 
                    \inparen{\Ex\nolimits_{0}\insquare{{\inangle{v, z}}^2}}^2\,.
           \end{align*} 

        \paragraphit{Part B (Bounding terms \cref{eq:SS:LC:ub2,eq:SS:LC:ub4,eq:SS:LC:ub5}):}~~
            First, note that
            \[
                \abs{\eqref{eq:SS:LC:ub2}}
                \leq 
                \Pr[\lnot\evE]\cdot 
                \inparen{ 
                    \Ex\nolimits_{t}\sinsquare{\inangle{v, z}^2\mid \lnot\evE}
                    +
                    \Ex\nolimits_{0}\sinsquare{\inangle{v, z}^2\mid \lnot\evE}
                }
                \,.
            \]
            Applying the Cauchy--Schwarz inequality yields
            \[
                \abs{\eqref{eq:SS:LC:ub2}}
                \leq 
                \sqrt{\Pr[\lnot\evE]}\inparen{ 
                    \sqrt{\Ex\nolimits_{t}\insquare{\inangle{v, z}^4}}
                    +
                    \sqrt{\Ex\nolimits_{0}\insquare{\inangle{v, z}^4}}
                }
                \,.
            \]
            Importantly, all expectations in the above are now unconditional. Repeating the same two steps for the remaining terms, we obtain
            \begin{align*}
                \abs{\eqref{eq:SS:LC:ub4}}
                    &\leq 
                    \Pr[\lnot\evE]
                    \inparen{
                        {\Ex\nolimits_{t}\sinsquare{{\inangle{v, z}}^4}}
                        +
                        {\Ex\nolimits_{0}\sinsquare{{\inangle{v, z}}^4}}
                    }\,,\\
                \abs{\eqref{eq:SS:LC:ub5}}
                    &\leq 
                    2\sqrt{\Pr[\evE]\Pr[\lnot\evE]}
                    \inparen{
                        {\Ex\nolimits_{t}\sinsquare{\inangle{v, z}^2}}
                        +
                        {\Ex\nolimits_{0}\sinsquare{\inangle{v, z}^2}}
                    }\,.
            \end{align*}
            Finally, since $\Pr[\evE]\le 1$ and, by \cref{lem:SS:LC:prob}, $\Pr[\lnot\evE]\le e^{-k}$, it follows that
            \begin{align*}
                \abs{\eqref{eq:SS:LC:ub2}}
                    &\leq 
                    e^{-k/2}\cdot \inparen{ 
                        \sqrt{\Ex\nolimits_{t}\sinsquare{\inangle{v, z}^4}}
                        +
                        \sqrt{\Ex\nolimits_{0}\sinsquare{\inangle{v, z}^4}}
                    }\,,\\
                \abs{\eqref{eq:SS:LC:ub4}}
                    &\leq 
                    e^{-k}\cdot 
                    \inparen{
                        {\Ex\nolimits_{t}\sinsquare{{\inangle{v, z}}^4}}
                        +
                        {\Ex\nolimits_{0}\sinsquare{{\inangle{v, z}}^4}}
                    }\,,\\
                \abs{\eqref{eq:SS:LC:ub5}}
                    &\leq 
                    2e^{-k/2}\cdot 
                    \inparen{
                        {\Ex\nolimits_{t}\sinsquare{\inangle{v, z}^2}}
                        +
                        {\Ex\nolimits_{0}\sinsquare{\inangle{v, z}^2}}
                    }\,.
            \end{align*}
        \end{proof}
        
        
    
    \subsubsection{Proof of \cref{lem:SS:LC:unconditionalMoments} (Tail Bounds on Unconditional Expectations)}\label{sec:proofof:lem:SS:LC:unconditionalMoments}
        \begin{proof}[Proof of \cref{lem:SS:LC:unconditionalMoments}]
            Divide $P_{\ymax}$ into $k$ parts $P_1,P_2,\dots,P_k$ where, for each $1\leq i\leq k$,
            \[
                P_i = \inbrace{ z\in \R^k \colon z_i = \ymax ~~\text{and}~~z_{-i}\leq \ymax}\,.
            \]
            See \cref{fig:self-selection-3D-distinct} for an illustration of the partition part of the parts $P_1,P_2,\dots,P_k$ with $k=3$.
            \begin{figure}[tbh!]
                \centering
                \includegraphics[width=0.5\linewidth]{fig/self_selection_3D_distinct.pdf}
                \caption{
                    This figure illustrates the partitions $P_1,P_2,\dots,P_k$ of the set $P_{\ymax}$ constructed in the proof of \cref{lem:SS:LC:unconditionalMoments} for $k=3$. 
                    Each colored slab in the figure corresponds to a set $P_i$ for $1\leq i\leq 3$, and the union of them corresponds to $P_{\ymax}$.
                    Compare this figure with \cref{fig:self-selection-3D} which illustrates $P_{\ymax}$.
                    \label{fig:self-selection-3D-distinct}
                }
            \end{figure} 
            Further, to simplify the notation, define 
            \[
                P\coloneqq P_{\ymax} 
                \qquadand
                \mu(t) \coloneqq W_t^\top x\,.
            \]
            Also recall that we use the short hands $\Ex_{t}[\cdot]$ and $\Pr_{t}[\cdot]$ to denote $\Ex_{z\sim \normal{\mu(t)}{I}}[\cdot]$ and $\Pr_{z\sim \normal{\mu(t)}{I}}[\cdot]$ respectively.
            First, observe that from the Cauchy--Schwarz inequality and the fact that $\norm{v}_2=1$
            \begin{align*}
                \begin{split}
                    \Ex\nolimits_{t}\insquare{\inangle{v,z}^2\mid z\in P}
                    &\leq 
                    \sum_i \Ex\nolimits_{t}\insquare{z_i^2 \mid z\in P}\,,\\
                    \Ex\nolimits_{t}\insquare{\inangle{v,z}^4\mid z\in P}
                    &\leq 
                    \Ex\nolimits_{t}\insquare{\norm{z}_2^4\mid z\in P}
                    \leq k \sum_i  \Ex\nolimits_{t}\insquare{z_i^4 \mid z\in P}\,.
                \end{split} 
                \yesnum\label{eq:SS:LC:unconditionalMoment:CS}
            \end{align*}
            Where we use Cauchy--Schwarz in the last inequality. 
            Moreover, for any $1\leq i\leq k$,  
            \begin{align*}
                \begin{split}
                    \Ex\nolimits_{t}\insquare{z_i^2  \mid z\in P}
                &= 
                \sum_j
                \Pr\nolimits_{t}\insquare{z\in P_j\mid z\in P}
                \cdot \Ex\nolimits_{t}\insquare{z_i^2 \mid z\in P_j}
                \leq \max_j \Ex\nolimits_{t}\insquare{z_i^2 \mid z\in P_j}\,,\\
                \Ex\nolimits_{t}\insquare{z_i^4 \mid z\in P}
                &= 
                \sum_j
                \Pr\nolimits_{t}\insquare{z\in P_j\mid z\in P}
                \cdot \Ex\nolimits_{t}\insquare{z_i^4 \mid z\in P_j}
                \leq \max_{j} \Ex\nolimits_{t}\insquare{z_i^4 \mid z\in P_j} \,.
                \end{split}
                \yesnum\label{eq:SS:LC:unconditionalMoment:AcrossLegs}
            \end{align*} 
            Fix any $1\leq i, j\leq k$.
            In the remainder of the proof, we upper bound 
            \[
                \Ex\nolimits_{t}\insquare{z_i^2 \mid z\in P_j}
                \quadand
                \Ex\nolimits_{t}\insquare{z_i^{\smash{4}} \mid z\in P_j}\,.
            \]
            We divide the upper bound into two cases.
            \begin{enumerate}
                \item \textbf{Case A ($i=j$):}~~ In this case, $z_i=\ymax$ and, hence,
                \[
                    \Ex\nolimits_{t}\insquare{z_i^2 \mid z\in P_j} = \ymax^2
                    \quadand
                    \Ex\nolimits_{t}\insquare{z_i^{\smash{4}} \mid z\in P_j} = \ymax^4 \,.
                \]
                \item \textbf{Case B ($i\neq j$):}~~ Observe that conditioned on $z\in P_j$, the distribution of $z_i$ is $\truncatedNormal{\mu_i(t)}{1}{(-\infty, \ymax]}$, \ie{}, the truncation of $\normal{\mu_i(t)}{1}$ to the set $(-\infty, \ymax]$.
                Next, we use the following fact which is proved in \cref{sec:proofof:fact:SS:LC:boundOnTruncatedMoments}.
                \begin{restatable}[]{fact}{factboundOnTruncatedMoments}
                    \label{fact:SS:LC:boundOnTruncatedMoments}
                    Fix $\mu,b\in \R$. It holds that 
                    \[
                        \Ex_{z\sim \normal{\mu}{1}}\insquare{z^2\mid z\leq b}
                        \leq 2\inparen{1+\abs{b}+\abs{\mu}}^2
                        \qquadand
                        \Ex_{z\sim \normal{\mu}{1}}\insquare{z^{\smash{4}}\mid z\leq b}
                        \leq 3\inparen{1+\abs{b}+\abs{\mu}}^4\,.
                    \]
                \end{restatable} 
                In our setting this implies the following bounds
                \begin{align*}
                    \Ex\nolimits_{t}\insquare{z_i^2 \mid z\in P_j} &\leq 
                        2 \inparen{1 + \abs{\ymax} + \abs{\mu_i(t)}}^2
                    \quadand
                    \Ex\nolimits_{t}\insquare{z_i^4 \mid z\in P_j} &\leq 
                        3 \inparen{1 + \abs{\ymax} + \abs{\mu_i(t)}}^4\,.
                \end{align*}
                Substituting $\mu_i(t)=\inangle{W_t^\top x}_i\leq \norm{W_t^\top x}_\infty$ implies that 
                \begin{align*}
                    \Ex\nolimits_{t}\insquare{z_i^2 \mid z\in P_j} &\leq 
                        2 \inparen{1 + \abs{\ymax} + \snorm{W_t^\top x}}^2
                    \quadand
                    \Ex\nolimits_{t}\insquare{z_i^4 \mid z\in P_j} &\leq 
                        3 \inparen{1 + \abs{\ymax} + \snorm{W_t^\top x}}^4\,.
                \end{align*} 
            \end{enumerate}
            Substituting the bounds from both cases into \cref{eq:SS:LC:unconditionalMoment:CS,eq:SS:LC:unconditionalMoment:AcrossLegs} implies that 
            \begin{align*}
                 \Ex\nolimits_{t}\sinsquare{\inangle{v,z}^2\mid z\in P} &\leq 
                 2k\inparen{1 + \abs{\ymax} + \snorm{W_t^\top x}_\infty}^2\,,\\\Ex\nolimits_{t}\sinsquare{\inangle{v,z}^4\mid z\in P} &\leq 
                 3k^2\inparen{1 + \abs{\ymax} + \snorm{W_t^\top x}_\infty}^4\,.
            \end{align*}
        \end{proof}

        
\subsection{Projection Set}
Given a $D$-warm start $W^{(0)}$ for $D = O(\nfrac1{\poly(k)})$,
we define the projection set
\[
    K\coloneqq \inbrace{W: \norm{W-W^{(0)}}_F\leq D\land \forall i\in [k], \norm{w_i}_2\leq C}\,.
    \yesnum\label{eq:SS:projectionSet}
\]
Here $C>0$ is the constant from \Cref{asmp:SS:separationBoundedness}.
We certainly have $W^\star\in K$.
When the hidden constant is sufficiently large,
we are also guaranteed that $\grad^2\negLL_\hyP \succeq 0$ over $K$ by \Cref{thm:SS:localConvexity}.

\subsection{Stochastic Gradient Oracle and Second Moment}
    The gradient of the negative log-likelihood function (\Cref{eq:SS:gradientExpression}) at $W\in K$ is given by
    \begin{align*}
        \grad \cL(W)
            = \Ex_x \Ex_{\ymax}
                \inparen{
                    \Ex_{z\sim \normal{W^\top x}{I} \mid z\in P(\ymax)}\insquare{xz^\top }
                    - 
                    xx^\top W
                }
            \,.
    \end{align*}
    
    \begin{lemma}\label{lem:SS:gradientOracle}
        Consider an instance of the max-self-selection problem
        and suppose \Cref{asmp:SS:gaussianity,asmp:SS:separationBoundedness} hold.
        Fix $W$ in the projection set $K$ (\Cref{eq:SS:projectionSet}).
        Suppose $(x, \ymax)$ is drawn from the max-self-selection model
        and $z\sim \normal{W^\top x}{I}\mid z\in P(\ymax)$.
        The following hold:
        \begin{enumerate}[(i)]
            \item $g(W)\coloneqq xz^\top - xx^\top W$ is an unbiased estimate of $\grad \negLL(W)$ \label{part:SS:gradientOracle:unbiased}
            \item $\E\insquare{\norm{g(W)}_F^2} = d\cdot O\inparen{k^4C^2 \log^3{k}}$ \label{part:SS:gradientOracle:secondMoment} 
            \item For any $W\in K$ in the projection set (\Cref{eq:SS:projectionSet}),
            there is an algorithm that consumes a single observation $(x, \ymax)$ 
            from the self-selection model
            to compute a random vector $\tilde g(W)$ such that $\tv{\tilde g(W)}{g(W)}\leq \delta$
            in time $O(kd\cdot \polylog(\nfrac1{\delta}))$. \label{part:SS:gradientOracle:algorithm}
        \end{enumerate}
    \end{lemma}

    \begin{proof}
        \sloppy
        \Cref{part:SS:gradientOracle:unbiased} follows by definition.
        The formal proof of \Cref{part:SS:gradientOracle:secondMoment} is through an involved calculation
        and we defer it to \cref{sec:SS:secondMoment}.

        In order to see \Cref{part:SS:gradientOracle:algorithm},
        we note that it suffices to approximately sample \mbox{$z\sim \normal{W^\top x}{I}\mid z\in P(\ymax)$}
        and output $\tilde g(W) = xz^\top - xx^\top W$.
        We can express the event $z\in P(\ymax)$ as union of events $\cup_{i\in [k]} P_i(\ymax)$ where
        \[
            P_i(\ymax) \coloneqq \inbrace{z_i=\ymax\land z_j\leq \ymax, \forall j\neq i}.
        \]
        Now, given $z\in P(\ymax)$,
        the conditional probability of $P_i(\ymax)$ is given by
        \[
            \Pr\insquare{z\in P_i(\ymax)\mid z\in P(\ymax)}
            = \frac{\Pr[z_i=\ymax]\cdot \prod_{j\neq i} \Pr[z_j\leq \ymax]}{\sum_{i=1}^k \Pr[z_i=\ymax]\cdot \prod_{j\neq i} \Pr[z_j\leq \ymax]}.
        \]
        Furthermore,
        given $z\in P_i(\ymax)$,
        we know that $z_i = \ymax$ and the rest of the coordinates $z_{-i}$ follows a truncated Gaussian distribution
        \[
            z_{-i}\mid z\in P_i(\ymax) \sim \normal{(W^\top x)_{-i}}{I}\mid z_{-i}\leq \ymax\,.
        \]
        
        The mixture probabilities can be expressed in closed form using the density and cumulative distribution function of the standard Gaussian distribution.
        We can thus estimate the mixture probabilities up to $\xi$-accuracy in $O(k\log\nfrac1\xi)$ time using an evaluation oracle for the error function (\eg{}, \cite[Proposition 3]{chevillard2012erf})
        and approximately sample from the mixture probability in $\xi$-TV distance.
        
        Suppose we sample $i$ from the mixture probabilities.
        We set $z_i = \ymax$ and sample the rest of the coordinates from the truncated Gaussian $z_{-i}\sim \normal{(W^\top x)_{-i}}{I}\mid z_{-i}\leq \ymax$.
        Since each coordinate of the truncated distribution is independent,
        we can perform the sampling for each coordinate individually.

        If $\ymax\geq w_j^\top x$,
        then we can simply perform rejection sampling with probability success at least $\nfrac12$.
        Suppose now that $\ymax< w_j^\top x$.
        \cite[Lemma 12]{daskalakis2019computationally} demonstrate how to approximately sample in $\xi$-TV distance from a truncated 1-dimensional Gaussian distribution $\truncatedNormal{\mu}{1}{[a, b]}$
        using inverse transform sampling.
        If $\alpha \coloneqq \normalMass{[a, b]}{\mu}{1}$,
        then their algorithm requires $\poly\log(\nfrac1\alpha, \nfrac1\xi)$ time.
        Fix a coordinate $z_j, j\neq i$.
        By \Cref{asmp:SS:separationBoundedness},
        $\norm{w_j}_2, \norm{\wstar_j}_2\leq C$.
        Since $\wstar_j^\top x\sim \normal{0}{\wstar_j^\top\wstar_j}$,
        $w_j\top x\sim \normal{0}{w_j^\top w_j}$,
        we see that $\ymax \geq -C-O(\sqrt{\log(\nfrac{1}{\delta})})$ and $w_j^\top x\leq C+O(\sqrt{\log(\nfrac{k}{\delta})})$ for all $j$
        with probability $1-\delta$ over the covariate $x$.
        In particular,
        $\abs{\ymax - w_j^\top x}\leq O(\sqrt{\log(\nfrac{k}{\delta})})$ with high probability
        so that by standard bounds on the Gaussian cumulative distribution function (\eg{}, \cite[Section 2.3]{kantor2015mathematics}),
        the survival probability $\alpha_j\coloneqq \normalMass{w_j^\top x}{1}{(-\infty, \ymax]}$
        is at least
        \[
            \Omega\inparen{\frac1{\sqrt{\log(\nfrac{k}{\delta})}} \exp\inparen{-\log(\nfrac{k}{\delta})}}
            = \Omega(\poly(\nfrac\delta{k}))\,.
        \]
        For any $a>0$,
        consider the event $z_j < \ymax-a \leq w_j^\top x - a$.
        We have
        \begin{align*}
            \Pr\insquare{z_j \leq \ymax - a\mid z_j\leq \ymax}
            = \frac{\Pr\insquare{z_j\leq \ymax-a}}{\Pr\insquare{z_j\leq \ymax}}
            \leq \frac1{\alpha_j} \Pr\insquare{z_j\leq w_j^\top x-a}
            \leq \poly(\nfrac{k}\delta) e^{-\frac{a^2}{2C^2}}\,.
        \end{align*}
        Choosing $a = \Omega(\log(\nfrac{k}{\delta\xi}))$ reduces the tail probability of the truncated Gaussian to at most $\xi$.
        We can thus run the inverse transform sampling algorithm of \cite{daskalakis2019computationally}
        on the distribution $\truncatedNormal{w_j^\top x}{1}{[\ymax-a, \ymax]}$ to approximately sample in $\xi$-TV distance from $\truncatedNormal{w_j^\top x}{1}{(-\infty, \ymax]}$.
        The running time is $\poly\log(\nfrac{k}\delta, \nfrac1\xi)$.

        Repeat this sampling algorithm for each coordinate $j\neq i$.
        Suppose we wish to bound the joint TV distance of every coordinate by $\xi$
        with probability $1-\delta$ over the observed covariate.
        This leads to a total running time of $k\cdot \polylog(\nfrac{k}\delta, \nfrac{k}\xi)$ for all sampling procedures.
        In the $\delta$ probability event of failure,
        we can simply output an arbitrary vector in $P(\ymax)$, and the TV distance is at most $\xi+\delta$.

        Once we have the desired sample $z$,
        we can perform the matrix multiplications $x(z^\top - x^\top W)$ in the order given by the parenthesis in $O(kd)$ arithmetic operations.
        This yields the final running time guarantees. 
    \end{proof}
    If we wish to run PSGD for $T$ iterations,
    we remark that it suffices to estimate each gradient to $\nfrac\xi{T}$-TV distance with probability $1-\nfrac\delta{T}$
    and the algorithm cannot distinguish between the approximate gradients and true gradients.


    \subsubsection{Bound on the Second Moment}\label{sec:SS:secondMoment}
        In this section, we bound the second moment of the gradient, as promised in the proof \cref{lem:SS:gradientOracle}.
        Fix any $V$ with $\norm{V}=1$.
        For each $t\in \R$, define $W_t = \Wstar + tV$.
        Our goal is to prove the following upper bound (for $0\leq t\leq 1$)%
        \[
            \Ex_x \Ex_{\ymax}
                    \inparen{
                        \Ex_{z\sim \normal{W_t^\top x}{I}}
                        \insquare{
                            \snorm{xz^\top - xx^\top W_t}_F^2 
                            \mid z\in P(\ymax)
                        }
                }
                \leq d\cdot O(k^4 C^2 \log^3{k})\,.
        \]
        Since $\snorm{xz^\top 
            - xx^\top W_t}_F^2 
            \leq 
            \snorm{x}_2^2 \snorm{z 
            - W_t^\top x}_2^2$, it suffices to prove that
        \begin{align*}
            \Ex_x\insquare{
                \Ex_{\ymax}
                \inparen{
                    \Ex_{z\sim \normal{W_t^\top x}{I}}
                    \insquare{
                        \snorm{z - W_t^\top x}_2^2
                        \mid z\in P(\ymax)
                    }
                }
                \cdot \norm{x}_2^2 
            }
            \leq d\cdot O(k^4 C^2 \log^3{k})\,.
            \yesnum\label{eq:SS:SecondMoment:sufficient}
        \end{align*}
        We divide the proof into three steps.

        \paragraph{Step 1 (Upper bound on $\Ex_z\sinsquare{
            \snorm{z - W_t^\top x}_2^2  \mid z\in P(\ymax)
        }$).}
            Observe that 
            \[
                \Ex_z\insquare{
                    \snorm{z - W_t^\top x}_2^2  \mid z\in P(\ymax)
                }
                \leq 
                2\Ex_z\insquare{
                    \norm{z}_2^2 \mid z\in P(\ymax)
                }
                + 2 \snorm{W_t^\top x}_2^2  \,.
            \]
            Substituting $v$ as $\wh{z}$ (\ie{}, the unit vector along $z$) in \cref{lem:SS:LC:unconditionalMoments} implies that 
            \[
                \Ex_z\insquare{
                    \norm{z}_2^2 \mid z\in P(\ymax)
                }\leq 
                    2k \inparen{1+\abs{\ymax} + \snorm{W_t^\top x}_\infty}^2
                    \leq 4k \inparen{1+\ymax^2 + \snorm{W_t^\top x}_\infty^2}
                    \,.
            \]
            Therefore, it follows that
            \[
                \Ex_z\insquare{
                    \snorm{z - W_t^\top x}_2^2  \mid z\in P(\ymax)
                }
                \leq 
                8k\inparen{1+{\ymax^2} + \snorm{W_t^\top x}_\infty^2}
                + 2 \snorm{W_t^\top x}_2^2  \,.
                \yesnum\label{eq:SS:secondMoment:ub1}
            \]

        \paragraph{Step 2 (Upper bound on $\Ex_{\ymax}\Ex_z[\snorm{z - W_t^\top x}_2^2 \mid z\in P(\ymax)]$).}
            Next, we take the expectation with respect to $\ymax$ and use \cref{fact:GaussMaxMoments} to upper bound $\Ex[\ymax^2]$ to obtain 
            \[
                \Ex_{\ymax}
                \Ex_z\insquare{
                    \snorm{z - W_t^\top x}_2^2  \mid z\in P(\ymax)
                }
                \leq 
                8k\cdot \inparen{
                    O\inparen{1 + \log{k}}
                    +\snorm{W_t^\top x}_\infty^2
                 }
                + 2 \snorm{W_t^\top x}_2^2  \,.
                \yesnum\label{eq:SS:SecondMoment:1}
            \]


        \paragraph{Step 3 (Upper bound on $\Ex_{x}[\Ex_{\ymax,z}[\snorm{z - W_t^\top x}_2^2 \mid z\in P(\ymax)]\cdot \norm{x}_2^2]$).}
            Finally, we will take the expectation over $x$.
            For this, we will apply the transformation explained in Step 2 in the proof of \cref{lem:SS:LC:step2} (in \cref{sec:proofof:lem:SS:LC:step2}).
            After this transformation, $W_t^\top x$ only depends on the first $2k$ coordinates of $x$, \ie{}, on $x_{\leq 2k}$.
            To make use of this transformation, we split $\norm{x}_2^2$ between $x_{\leq 2k}$ and $x_{> 2k}$ to yield
            \begin{align*}
                &\Ex_{x}\insquare{
                \Ex_{\ymax}
                \Ex_z\insquare{
                    \snorm{z - W_t^\top x}_2^2  \mid z\in P(\ymax)
                }
                \cdot 
                \norm{x}_2^2 
                }\\
                &\quad=\quad
                \Ex_{x_{>2k},~x_{\leq 2k}}\insquare{
                    \Ex_{\ymax}
                    \Ex_z\insquare{
                        \snorm{z - W_t^\top x}_2^2  \mid z\in P(\ymax)
                    }
                    \cdot 
                    \inparen{
                        \norm{x_{\leq 2k}}_2^2 
                        + 
                        \norm{x_{> 2k}}_2^2 
                    }
                }\,.
                \yesnum\label{eq:SS:SecondMoment:3}
            \end{align*}
            Observe that
            \begin{align*}
                \snorm{W_t^\top x}_\infty 
                \leq \snorm{W_t^\top x}_2
                &\leq \sum_i \norm{w^\star_i + t v_i}_2 \norm{x_{\leq 2k}}_2
                \leq 
                k(C+1)\norm{x_{\leq 2k}}_2
                ~~\Stackrel{(C\geq 1)}{\leq}~~ 2k^{3/2}C \norm{x_{\leq 2k}}_\infty \,.
                \yesnum\label{eq:SS:SecondMoment:2}
            \end{align*}
            Hence, \cref{eq:SS:SecondMoment:1,eq:SS:SecondMoment:3,eq:SS:SecondMoment:2} imply that 
            \begin{align*}
                &\Ex_{x}\insquare{
                \Ex_{\ymax}
                \Ex_z\insquare{
                    \snorm{z - W_t^\top x}_2^2  \mid z\in P(\ymax)
                }
                \cdot 
                \norm{x}_2^2 
                }\\
                &\quad=\quad
                {O(k^4C^2\log{k})}
                \cdot \Ex_{x_{>2k},~x_{\leq 2k}}\insquare{
                    {
                        \inparen{
                            1
                            + \norm{x_{\leq 2k}}_\infty^2
                        }
                        \cdot 
                        \inparen{
                            \norm{x_{\leq 2k}}_2^2 
                            + 
                            \norm{x_{> 2k}}_2^2 
                        }
                    }
                }\,.
            \end{align*}
            This can further be simplified to 
            \[
                {O(k^4C^2\log{k})}
                \cdot \Ex_{x_{>2k}}
                    \insquare{
                        \Ex_{x_{\leq 2k}}\insquare{
                            \norm{x_{\leq 2k}}_\infty^2
                            + \norm{x_{\leq 2k}}_\infty^4
                        }
                        + 
                        \norm{x_{> 2k}}_2^2 \cdot 
                        \Ex_{x_{\leq 2k}}\insquare{
                            1
                            + \norm{x_{\leq 2k}}_\infty^2
                        }
                    }\,.
            \]
            Next, \cref{fact:GaussMaxMoments} implies the following upper bound
            \begin{align*}
                &\Ex_{x}\insquare{
                \Ex_{\ymax}
                \Ex_z\insquare{
                    \snorm{z - W_t^\top x}_2^2  \mid z\in P(\ymax)
                }
                \cdot 
                \norm{x}_2^2 
                }
                ~~\leq~~
                    {O(k^4C^2\log{k})}\cdot \Ex_{x_{> 2k}}\insquare{
                    O\inparen{\log^2{k}}
                    + \norm{x_{> 2k}}_2^2 \cdot O\inparen{\log{k}}
                }\,.
            \end{align*}
            Finally, using $\Ex_{x_{>2k}}[\norm{x_{>2k}}_2^2]=(d-2k)\,\Ex_{u\sim \normal{0}{1}}[z^2]=O(d-2k)$, we get that  
            \begin{align*}
                \Ex_{x}\insquare{
                \Ex_{\ymax}
                \Ex_z\insquare{
                    \snorm{z - W_t^\top x}_2^2  \mid z\in P(\ymax)
                }
                \cdot 
                \norm{x}_2^2 
                }
                &~~\leq~~
                (d-2k)\cdot O\inparen{{k^4C^2} \log^3{k}}
                \leq d\cdot O\inparen{{k^4C^2} \log^3{k}}\,.
            \end{align*}
            
        

    
\subsection{Projected Stochastic Gradient Descent}
    We are now ready to prove \Cref{thm:SS:localConvergence} by applying the iterative PSGD algorithm (\Cref{thm:SGD-local-growth-convergence}).
    Our complete algorithm (\Cref{thm:SS}) follows by combining this result with an appropriate warm start.
    We restate the theorem below for convenience.
    \thmSelfSelectionLocalConvergence*

    \begin{proof}
        By \Cref{thm:SS:localConvexity},
        $\negLL_\hyP$ is convex over $K$.
        Define
        \[
            \alpha \coloneqq \inparen{\frac{c}{ek}}^{O(C^2/c^2)} = \poly(\nfrac1k).
        \]
        Combining \Cref{thm:SS:localConvexity} and \Cref{thm:info-pres-implies-strong-conv}
        shows that $\negLL_\hyP$ satisfies a $\Omega(\alpha)$-local growth condition (\Cref{def:local-growth}) over $K$. 

        \Cref{lem:SS:gradientOracle} shows that
        we have an (approximate) stochastic gradient oracle
        with second moment $G^2 = \tilde O(k^4d)$ and running time $\tilde O(kd \log(\nfrac1\delta))$.
        Thus the given $D$-warm start has an objective value at most $DG = O(\sqrt{d}\cdot \poly(k))$.

        As we would like an $O(\alpha^2\eps^2)$-optimal solution in order to recover the parameters to accuracy $\eps$,
        \Cref{thm:SGD-local-growth-convergence} ensures that it suffices to take
        \[
            m = \tilde O\inparen{\frac{G^2}{\alpha^4 \eps^2}}
            = \tilde O\inparen{\frac{d\cdot \poly(k)}{\eps^2}}
        \]
        gradient steps with each step consuming one sample from the self-selection model.
        Each gradient computation takes $\tilde O(kd\log(\nfrac1\eps))$ running time 
        in order to reduce the joint TV-distance sampling error over all gradients to $0.01$.
        Moreover,
        each gradient step requires $O(kd)$ arithmetic operations.
        Thus the total running time becomes
        \[
            \tilde O(mkd \log(\nfrac1\eps))
            = \tilde O\inparen{\frac{d^2\cdot \poly(k)}{\eps^2}}.
        \] 
        Our algorithm recovers $w_1^\star, \dots, w_k^\star$ in $\eps$-Euclidean distance with probability $0.98$ up to some permutation.
        As described in \Cref{sec:SGD:clusteringBoosting},
        we wish to repeat the algorithm $O(\log(\nfrac1\delta))$ times and output the solution close to at least 50\% of the points.
        In our setting,
        we can interpret closeness under permutation as a metric over a quotient space over $\R^{d\times k}$
        and perform the clustering trick from \Cref{sec:SGD:clusteringBoosting} using this metric.
    \end{proof}
% --- end self_selection.tex ---

\section{Linear Regression with Second-Price Auction Data}\label{sec:2ndPrice}



% --- begin second_price.tex ---



In this section, we prove \cref{thm:2ndPrice:localConvexity}, which we restate below.
\thmsecondPriceLocalConvexity*


\subsection{Information Preservation}\label{sec:2ndPrice:proof:informationPreservation}
    
    In this section, we prove the first claim in \cref{thm:2ndPrice:localConvexity}.
    This proof borrows results and analysis from the proof of the first claim in \cref{thm:SS:localConvexity}, which lower bounds the information preservation for self-selection.
    Hence, it would be helpful for the reader to familiarize themselves with the proof of the information preservation claim in \cref{thm:SS:localConvexity} first (see  \cref{sec:SS:proof:informationPreservation}).


    Observe that the only difference between the two instances is the observations:
    In the max-self-selection problem, one observes $(x,\ymax)$, and in the self-selection problem arising from Second Price auctions, one observes $(x,\ysmax,\imax)$.
    Toward proving information preservation, 
    fix any parameters $V=[v_1,\dots,v_k]$ and $W=[w_1,\dots,w_k]$ close to each other and $\Wstar$ in the following sense  
    \[
        \norm{V - W^\star}_F\,,~\norm{W - W^\star}_F
        \leq \frac{c}{{400C\log{\nfrac{k}{c}}}}\,.
        \yesnum\label{eq:2ndPrice:distance}
    \]
    Since the only difference between the two problems is in the observations, we can reproduce the proof in \cref{sec:SS:proof:informationPreservation} until \cref{eq:SS:probabilityDiff}.
    At this point, we get that it is sufficient to prove that, for some $1\leq m\leq k$
    \begin{align*}
        &\abs{
            \begin{array}{cc}
                &\Pr\insquare{y_{\smax, V} \leq 0, \imax=m\mid \evE}\\
                - &\Pr\insquare{y_{\smax, W} \leq 0, \imax=m\mid \evE}
            \end{array}
        }
        \geq %
        \frac{1}{\Pr[\evE]}
        \cdot \inparen{\frac{c}{ek}}^{O(C^2/c^2)}
        \cdot \norm{v_i-w_i}_2\,.
        \yesnum\label{eq:2ndPrice:probAsDiff}
    \end{align*}
    Where $\evE=\evE_{i,\gamma,R}$ is the event from \cref{def:SS:event} and $i$ is any coordinate in $\argmax_{1\leq j\leq k} \norm{v_j-w_j}_2$ (as defined in \cref{eq:SS:i}).
    After the translation performed at the start of \cref{sec:proofof:eq:SS:conditionalInfoPreservation} (see \cref{eq:SS:translation}), it holds that 
    \[
        \rho_{i,V}\,, ~\rho_{i,W}\geq 0\,.
        \yesnum\label{eq:2ndPrice:2}
    \]
    Toward proving the above, we re-write the two terms on the left-hand side:
    
    \begin{align*}
        \begin{split}
        \Pr\insquare{y_{\smax, V} \leq 0, \imax=i\mid \evE}
            &= \inparen{
                1-\Phi\inparen{-\rho_{i,v}; \sigma_{i,V}^2}
              }\cdot 
              \prod_{j \neq i}
                \Phi\inparen{-\rho_{j,V}; \sigma_{j,V}^2}
            \,,
        \\
        \Pr\insquare{y_{\smax, W} \leq 0, \imax=i\mid \evE} 
            &= \inparen{
                    1-\Phi\inparen{-\rho_{i,W}; \sigma_{i,W}^2}
                  }\cdot 
                  \prod_{j \neq i}
                    \Phi\inparen{-\rho_{j,W}; \sigma_{j,W}^2}
        \,. 
        \end{split} 
    \end{align*}
    Where, for any $\sigma^2\geq 0$, $\customcdf{\cdot}{}$ is the cumulative density function of the one-dimensional normal distribution $\cN(0,\sigma^2)$.
    Since $1-\customcdf{-r}{}=\customcdf{r}{}$, it follows that 
    \begin{align*}
        \begin{split}
        \Pr\insquare{y_{\smax, V} \leq 0, \imax=i\mid \evE}
            &= {
                \Phi\inparen{\rho_{i,v}; \sigma_{i,V}^2}
              }\cdot 
              \prod_{j \neq i}
                \Phi\inparen{-\rho_{j,V}; \sigma_{j,V}^2}
            \,,
        \\
        \Pr\insquare{y_{\smax, W} \leq 0, \imax=i\mid \evE} 
            &= {
                    \Phi\inparen{\rho_{i,W}; \sigma_{i,W}^2}
                  }\cdot 
                  \prod_{j \neq i}
                    \Phi\inparen{-\rho_{j,W}; \sigma_{j,W}^2}
        \,. 
        \end{split}
        \yesnum\label{eq:2ndPrice:probabilityExpressions}
    \end{align*}
    Observe that
    \begin{align*}
        &\abs{ 
            \Pr\insquare{y_{\smax, V} \leq 0,\imax=i\mid \evE}
            - \Pr\insquare{y_{\smax, W} \leq 0,\imax=i\mid \evE} 
        }\\
        &~~= 
        \underbrace{\inparen{
            {
                \customcdf{\rho_{i,V}}{i,V}
            }\cdot \prod_{j\neq i}~ \customcdf{-\rho_{j,V}}{j,V}
        }}_{\textsf{A}}
        \cdot  
        \biggabs{
            1 - 
            \underbrace{\frac{
                \customcdf{\rho_{i,W}}{i,W}
            }{
                \customcdf{\rho_{i,V}}{i,V}
            }}_{\textsf{B}}
            \cdot 
            \underbrace{\prod_{j\neq i}~ \frac{
                \customcdf{-\rho_{j,W}}{j,W}
            }{
                \customcdf{-\rho_{j,V}}{j,V}
            }}_{\textsf{C}}
        }\,.
        \yesnum\label{eq:2ndPrice:differenceAsProduct}
    \end{align*}
    Next, we bound Terms~\textsf{A}, \textsf{B}, and \textsf{C} respectively.
    \begin{enumerate}
        \item Term~\textsf{A} is very similar to Term~\textsf{A} in \cref{eq:2ndPrice:differenceAsProduct}.
        Toward bounding it, recall that \cref{eq:SS:infoPreservation:toReferenceLater1} implies that 
        \[
            \prod_{j\neq i}~ \customcdf{-\rho_{j,V}}{j,V}
            \geq \inparen{1 - \frac{1}{k^{12}}}^{k-1}
            \geq 1 - \frac{1}{k^{11}}\,.
        \]
        Moreover, since $\rho_{i,V}\geq 0$,
        \[
            \customcdf{\rho_{i,V}}{i,V}
            \geq \frac{1}{2}\,.
        \]
        Therefore, it follows that 
        \[
            \textsf{A}
            \geq \frac{1}{2}\inparen{1 - \frac{1}{k^{11}}}
            ~~\Stackrel{(k\geq 2)}{\geq} ~~ \frac{1}{4}\,.
        \]
        \item Recall that after the translation described earlier $\min\inbrace{\rho_{i,V}, \rho_{i,W}}=0$.
        Without loss of generality, let $\rho_{i, W}=0$.
        By \cref{lem:SS:propertiesOfRho} we have that 
        \[
            \abs{\rho_{i, V}} = \abs{\rho_{i, W} - \rho_{i, V}}\geq \frac{5R}{6} \cdot \sqrt{\log{\nfrac{1}{\gamma}}}\cdot \norm{v_i - w_i}_2\,.
        \]
        Substituting $\mu=-\rho_{i,V}$, $\sigma^2=\sigma_{i,V}^2$, and $\nu^2=\sigma_{i,W}^2$ in \cref{lem:SS:cdfRatiosLB} implies that 
        \[
            \frac{
                \customcdf{-\rho_{i,V}}{i,V}
            }{
                \customcdf{-\rho_{i,W}}{i,W}
            }
            \notin 1 \pm \frac{1}{\sqrt{2\pi e}}
                \min\inbrace{
                    1\,, 
                    \frac{5R}{6\sigma_{i,V}} \cdot \sqrt{\log{\nfrac{1}{\gamma}}}\cdot \norm{v_i - w_i}_2
                } 
            \,.
        \]
        Substituting $\sigma_{i,V}^2=1+C^2+(16c^4/C^2)\leq 17+C^2\leq 18C^2$ (from \cref{eq:SS:distributionOfZeta}), 
        \[
            \frac{
                \customcdf{\rho_{i,V}}{i,V}
            }{
                \customcdf{\rho_{i,W}}{i,W}
            }
            \notin 1 \pm \frac{1}{6C\sqrt{\pi e}}\cdot 
                \min\inbrace{
                    1\,,
                    \frac{5R}{6}\cdot \sqrt{\log{\nfrac{1}{\gamma}}}\cdot \norm{v_i-w_i}_2 
                }
            \,.
        \] 
        \item Further, Term~\textsf{C} above is the same as Term~\textsf{C}  in \cref{eq:SS:differenceAsProduct} and, \cref{eq:SS:termC:final} implies that
        \[
            \textsf{C}\in 1\pm \frac{100\sqrt{72}}{e^{20}} \cdot \norm{v_j-w_j}_2\cdot \frac{c}{k}\,.
        \]
    \end{enumerate}
    Since all three terms (Terms~\textsf{A}, \textsf{B}, and \textsf{C}) in \cref{eq:2ndPrice:differenceAsProduct} satisfy the same conditions as the corresponding terms (\cref{eq:2ndPrice:differenceAsProduct}) in the proof of (max) self-selection problem, \cref{eq:SS:secondTermLB} and subsequent arguments in the proof imply that 
    \[
        \tv{\cM(V)}{\cM(W)}~~
        \geq \inparen{\frac{c}{ek}}^{{O({C^2/c^2}){}}}\cdot \norm{V-W}_F\,.
    \]
    This completes the proof of information preservation for the self-selection problem arising from second-priced auctions.
    
\subsection{Local Convexity}\label{sec:2ndPrice:localConvexity}
    In this section, we prove the second claim in \cref{thm:2ndPrice:localConvexity}.
    As with the proof of the first claim, this proof borrows results and analysis from the proof of the corresponding claim in \cref{thm:SS:localConvexity}.
    Hence, it would be helpful for the reader to familiarize themselves with the proof of the local convexity in \cref{thm:SS:localConvexity} first (which is presented in \cref{sec:SS:proof:localConvexity}).

    First, we begin by defining the negative log-likelihood and presenting its Hessian.
    Then we present analogs of the lemmas from \cref{sec:SS:proof:localConvexity} that are sufficient to prove convexity for the self-selection problem arising from second-price auction data.
    Finally, we explain how these lemmas can be proved by adapting the proofs of lemmas in \cref{sec:SS:proof:localConvexity}.
    

    \paragraph{Negative Log-Likelihood and Its Hessian.}
        To state the negative log-likelihood, define the following set: given a value of the second max $s$ and an index of the winner $i$, define 
        \[
            P(s,i)\coloneqq 
            \inbrace{
                z\in \R^k\colon 
                    z_{\smax} = s ~~\text{and}~~ z_{\imax} = z_{\max}
            }\,.
        \]
        Now, the conditional negative log-likelihood and its Hessian are (see \cref{sec:hessian:2ndPrice} for a proof):
        \begin{align*}
            \cL(W; x, \ysmax, \imax) &= -\log{
                \int_{P(\ysmax, \imax)} \exp\inparen{-\frac{1}{2}\norm{z-W^\top x}_2^2}\d z 
            }\,,\\
            \grad^2 \cL(W; x, \ysmax, \imax) &= 
            xx^\top \otimes I_{k} - {
                \cov_{z\sim \normal{W^\top x}{I}}\insquare{z\mid z\in P\inparen{\ysmax,\imax}}
                \otimes xx^\top
            } 
            \,.
            \yesnum\label{eq:2ndPrice:conditionalHe2ndPrice:mainBody}
        \end{align*}
        Taking the expectation over the observation implies that population negative log-likelihood and its Hessian are as follows:
        \begin{align*}
            \cL(W) &= \Ex_{x,\ysmax,\imax}~\cL(W; x, \ysmax, \imax)\,,\\
            \grad^2 \cL(W) &= I_{dk} - \Ex_{x,\ysmax,\imax}\inparen{
                \cov_{z\sim \normal{W^\top x}{I}}\insquare{z\mid z\in P\inparen{\ysmax,\imax}}
                \otimes xx^\top
            }\,.
        \end{align*}

    \paragraph{Proof of Local Convexity.}
        Our goal is to show that 
        \[
            \forall~W\quadtext{such that}
            \norm{W-\Wstar}_2 \leq k^{-O(\sfrac{C^2}{c^2})}\,,
            \quadtext{it holds} \nabla^2 \cL(W) \succeq 0 \,.
            \yesnum\label{eq:2ndPrice:LC:sufficientCondition}
        \] 
    
        \paragraph{Strong Convexity at $W^\star$.}
            First, we will prove strong convexity of the NLL at the true parameter $W = W^\star$.
        Since $\alpha$-information preservation implies $\alpha$-quadratic growth (see \cref{thm:quadGrowth}), the information preservation property in \cref{eq:2ndPrice:informationPreservation} (proved in \cref{sec:2ndPrice:proof:informationPreservation}) yields, for some \textit{absolute} constant $A\ge 1$,
        \[
            \phantom{.}
            \qquad\qquad\qquad\qquad\qquad\quad~~
            \nabla^2 \cL(W^\star) \succeq k^{-A\cdot C^2/c^2}\, I\,. 
            \qquad~~
            \text{(Strong Convexity at $\Wstar$)}\hspace{-10mm}
            \yesnum\label{eq:2ndPrice:LC:hessianLB}
        \]
        Indeed, a Taylor expansion of $\cL(W)$ around $W^\star$ shows that, as $W\to W^\star$, the higher-order terms become negligible compared to the Hessian (since $W^\star$ is a stationary point of $\cL$ -- see \cref{fact:2ndPrice:stationaryPoint} -- the first-order term is zero), thus establishing \cref{eq:2ndPrice:LC:hessianLB}. 

    \paragraph{Strong Convexity Near $W^\star$.}
    Given \cref{eq:2ndPrice:LC:hessianLB}, local convexity is implied by the following lower bound on $\nabla^2 \negLL(\cdot)$ around $\Wstar$.
    To simplify the notation define $B\geq 1$ as the following quantity 
    \[
        B\coloneqq O\inparen{\frac{\sqrt{A}C}{c}}\,,
        \yesnum\label{eq:2ndPrice:LC:defB}
    \]
    where the hidden constant is appropriately large and can be deduced from the proof of \cref{thm:2ndPrice:LC:formal} below.
    \begin{theorem}[Strong Convexity in Neighbourhood of $W^\star$]\label{thm:2ndPrice:LC:formal}
    For $\rho\in(0,k^{-\Omega(B^2)})$, $t=\rho\cdot \wt{O}\inparen{\frac{1}{C^2B^4k^{17/2}}}$, and $V\in\R^{d\times k}$ satisfying $\norm{V}_F=1$, 
    \[
        W_t\coloneqq W^\star+tV\,.
    \]
    Then, 
    \[
        \nabla^2 \cL(W_t) 
        \succeq 
        \nabla^2 \cL(W^\star)
            - (\rho + k^{-\Omega(B^2)}) I\,.
    \] 
    \end{theorem}
    First, we pause to observe that \cref{thm:2ndPrice:LC:formal} and \cref{eq:2ndPrice:LC:hessianLB} imply local convexity.
    To see this, note that any parameter $W$ that is $k^{-\Omega(B^2)}$ close to $W^\star=W_0$ in Frobenius norm can be expressed as $W=W_t$ for a suitable matrix $V$ (with $\norm{V}_2=1$), and constants $\rho=k^{-\Omega(B^2)}$ and $t=\rho\cdot \wt{O}\inparen{{C^{-2}B^{-2}k^{-17/2}}}$. 
    In words, we get that at $W^\star$ the Hessian is at least $\poly(\nfrac{1}{k}) \cdot I$ and at any point $W$ which is $\poly(\nfrac{1}{k})$-close to $W^\star$, the Hessian decreases by some quantity of order at most $\poly(\nfrac{1}{k})$ and hence, by tuning the constants appropriately, remains {positive semi-definite}.

    As in \cref{sec:SS:proof:localConvexity}, we divide the proof of this result into two lemmas.
    \begin{lemma}[Lower Bound on Sample NLL's Hessian]\label{lem:2ndPrice:LC:step1}
        For $\zeta\in(0,1)$ and $t=\zeta\cdot O\inparen{\frac{1}{B^4k^3}}$, 
        \begin{align*}
            \forall_{x\in \R^d}\,,~\forall_{\ysmax\in \R}~\forall_{\imax\in [k]}\,,\quad 
            \nabla^2 \cL(W_t;x,\ysmax,\imax) 
            &- \nabla^2 \cL(W_0;x,\ysmax,\imax)\\
            &\qquad\succeq -\gamma(x,\ysmax,\imax) \cdot I_{k}\otimes xx^\top\,,
        \end{align*}
        where
        \[ 
            \gamma(x,\ysmax,\imax) \coloneqq 
                \Theta(k^2)
                \cdot
                \Bigl(\zeta + e^{-B^2 k/2}\Bigr)\cdot
                \Bigl(1 + \ysmax^4 + \max\sinbrace{\|W_t^\top x\|_\infty, \|W_0^\top x\|_\infty}^4\Bigr)
                \,.
        \]
    \end{lemma} 
    
    \begin{lemma}[Bounds on $\Ex\insquare{\gamma}$]\label{lem:2ndPrice:LC:step2}
        Let $\gamma$ as in \Cref{lem:2ndPrice:LC:step1}.
        For $0\leq \zeta,t\leq 1$ 
        \[
            0\preceq \Ex_{x,\ysmax,\imax}\Bigl[\gamma(x,\ysmax,\imax) \cdot I_{k}\otimes xx^\top \Bigr] \preceq 
            C
                \cdot \wt{O}(k^{11/2})
                \cdot (\zeta + e^{-B^2 k/2})
                \cdot I_{dk}\,.
        \]
    \end{lemma}
    The proofs of \cref{lem:2ndPrice:LC:step1,lem:2ndPrice:LC:step2} are identical to the proofs of their analogs (\cref{lem:SS:LC:step1,lem:SS:LC:step2}): the only change required is suitably changing the observations from $(x,\ymax)$ to $(x,\ysmax,\imax)$ throughout the proof.
    Note that this is only a syntactic change and does not require any new arguments; we explain one exception to this in the remark below.
    \begin{remark}
        One change required is in the proof of \cref{lem:SS:LC:unconditionalMoments}, which upper bounds 
        \[
            \Ex_{z\sim \normal{W_t^\top x}{1}}\insquare{\sinangle{v,z}^2\mid z\in P(\ymax)}
            \quadand
            \Ex_{z\sim \normal{W_t^\top x}{1}}\insquare{\sinangle{v,z}^4\mid z\in P(\ymax)}\,.
        \]
        Its proof divides the moment into $k$-parts one corresponding to each coordinate of $z$ and, further, breaks each conditioning $z\in P(\ymax)$ into $k$-subparts corresponding to the $k$-parts of $P_{\ymax}$; see \cref{fig:self-selection-3D-distinct}. 
        One can repeat the same proof for $P(\ysmax,\imax)$ since it has $(k-1)$-parts and these parts have a similar structure as the parts of $P_{\ymax}$: they fix two coordinate's value and restrict the remaining $(k-2)$-coordinates to lie in an interval; see \cref{fig:2nd-price-observation-distinct}. 
        One difference is that one of the slabs (which corresponds to $\imax$) pleases a constraint of the form $z_i\geq \ysmax$ instead of $z_i\leq \ysmax$, but this is easily handled by considering $w=-z_i$ and using that $\Ex[w^2]=\Ex[z^2]$ and $\Ex[w^4]=\Ex[z^4]$.
    \end{remark}
    Finally, \cref{lem:2ndPrice:LC:step1,lem:2ndPrice:LC:step2} imply the desired local convexity: the proof of this is also analogous to the one in \cref{sec:SS:proof:localConvexity} and again follows by replacing $\ymax$ by $(\ysmax,\imax)$.

\subsection{Projection Set}
The projection set is identical to that of self-selection (\Cref{eq:SS:projectionSet}).
Given a $D$-warm start $W^{(0)}$ for $D = O(\nfrac1{\poly(k)})$,
we define the projection set
\[
    K\coloneqq \inbrace{W\colon \norm{W-W^{(0)}}_F\leq D
        \quadand
        \forall i\in [k]\,,~~ \forall \norm{w_i}_2\leq C}\,.
    \yesnum\label{eq:2ndPrice:projectionSet}
\]
Here $C>0$ is again the constant from \Cref{asmp:SS:separationBoundedness}.
Moreover, $W^\star\in K$ and $\grad^2\negLL_\hyP \succeq 0$ over $K$ by \Cref{thm:2ndPrice:localConvexity}.

\subsection{Stochastic Gradient Oracle and Second Moment}
    The gradient of the negative log-likelihood function (see \Cref{sec:hessian:2ndPrice}) at $W\in K$ is given by
    \begin{align*}
        \grad \cL(W)
            = \Ex_x \Ex_{\ymax}
                \inparen{
                    \Ex_{z\sim \normal{W^\top x}{I} \mid z\in P(\ysmax,\imax)}\insquare{xz^\top }
                    - 
                    xx^\top W
                }
            \,.
    \end{align*}
    
    \begin{lemma}\label{lem:2ndPrice:gradientOracle}
        Consider an instance of the {second price auction} (\cref{def:secondPrice}) problem
        and suppose \Cref{asmp:SS:gaussianity,asmp:SS:separationBoundedness} hold.
        Fix $W\in K$ in the projection set (\Cref{eq:2ndPrice:projectionSet}).
        Suppose $(x, \ymax)$ is drawn from the {second price auction} model (\cref{def:secondPrice}) 
        and $z\sim \normal{W^\top x}{I}\mid z\in P(\ysmax, \imax)$.
        The following hold:
        \begin{enumerate}[(i)]
            \item $g(W)\coloneqq xz^\top - xx^\top W$ is an unbiased estimate of $\grad \negLL(W)$; \label{part:2ndPrice:gradientOracle:unbiased}
            \item $\E\insquare{\norm{g(W)}_F^2} = d\cdot O\inparen{k^4C^2  \log^3{k}}$; and  \label{part:2ndPrice:gradientOracle:secondMoment}  
            \item For any $W\in K$ in the projection set (\Cref{eq:2ndPrice:projectionSet}),
            there is an algorithm that consumes a single observation $(x, \ymax, \imax)$ 
            from the self-selection model
            to compute a random vector $\tilde g(W)$ such that $\tv{\tilde g(W)}{g(W)}\leq \delta$
            in time $O(kd\cdot \polylog(\nfrac1{\delta}))$. \label{part:2ndPrice:gradientOracle:algorithm}
        \end{enumerate}
    \end{lemma}

    \begin{proof}
        \sloppy
        \Cref{part:2ndPrice:gradientOracle:unbiased} follows by definition.
        The formal proof of \Cref{part:2ndPrice:gradientOracle:secondMoment} is through an involved calculation
        and we defer it to \cref{sec:2ndPrice:secondMoment}.
        The only necessary change from between \Cref{part:SS:gradientOracle:algorithm} in \Cref{lem:SS:gradientOracle} and \Cref{part:2ndPrice:gradientOracle:algorithm} is sampling from $z\sim \normal{W^\top x}{I}\mid z\in P(\ysmax, \imax)$,
        which we now sketch.

        Decompose the set $P(\ysmax, \imax)$ as
        $
            \cup_{\imax\neq j [k]} P_{j}(\ysmax, \imax)
        $
        where
        \[
            P_{j}(\ysmax, \imax) \coloneqq \inbrace{
                z_{\imax}\geq \ysmax,~~z_j=\ysmax~~\text{and}~~\forall \ell\neq j, \imax, z_\ell\leq \ysmax
            }\,.
        \]
        Similar to \Cref{lem:SS:gradientOracle},
        we can compute the probability $z\in P_j(\ysmax, \imax)$ given that $z\in P(\ysmax, \imax)$
        by evaluating the density and cdf of the standard Gaussian distribution.
        Thus we can compute the conditional probability that $z_j = \ysmax$.
        Upon fixing the second max coordinate $j$,
        we sample the rest of the coordinates $z_{-j, -\imax}$ independently using inverse transform sampling
        in an identical fashion to \Cref{lem:SS:gradientOracle}.
        Finally,
        we must sample $z_{\imax}\sim \truncatedNormal{w_{\imax}^\top x}{1}{[\ysmax, \infty)}$.
        This can also be done using inverse transform sampling.
    \end{proof}

    \subsubsection{Bound on the Second Moment}\label{sec:2ndPrice:secondMoment}
        \begin{figure}[tbh!]
            \centering
            \includegraphics[width=0.5\linewidth]{fig/2nd_price_one_observation_distinct.pdf}
            \caption{
            This figure illustrates the partitions $P_1,P_2,\dots,P_{k-1}$ of the set $P(\ysmax,\imax)$ mentioned in  \cref{footnote:secondPriceBounds} for $k=3$.
            Each colored slab in the figure corresponds to a set $P_i$ for $1\leq i\leq 3$, and the union of them corresponds to $P(\ysmax,\imax)$. 
            Compare this figure with \cref{fig:2nd-price-observation} which illustrates $P(\ysmax,\imax)$.
            \label{fig:2nd-price-observation-distinct}
            }
        \end{figure}  
        \noindent In this section, we bound the second moment of the gradient, as promised in the proof of \cref{lem:2ndPrice:gradientOracle}.
        Fix any $V$ with $\norm{V}=1$.
        For each $t\in \R$, define $W_t = \Wstar + tV$.
        Our goal is to prove the following upper bound (for $0\leq t\leq 1$)
        \[
            \Ex_x \Ex_{\ymax}
                    \inparen{
                        \Ex_{z\sim \normal{W_t^\top x}{I}}
                        \insquare{
                            \norm{xz^\top - xx^\top W_t}_F^2 
                            \mid z\in P(\ysmax,\imax)
                        }
                    }
                \le d\cdot O(k^4C^2\log^2{k})\,.
        \]
        Since $\norm{xz^\top 
            - xx^\top W_t}_F^2 
            \leq 
            \norm{x}_2^2  \norm{z 
            - W_t^\top x}_2^2$, it suffices to prove the following 
        \begin{align*}
            \Ex_x\insquare{
                \Ex_{\ymax}
                \inparen{
                    \Ex_{z\sim \normal{W_t^\top x}{I}}
                    \insquare{
                        \snorm{z - W_t^\top x}_2^2
                        \mid z\in P(\ysmax,\imax)
                    }
                }
                \cdot \norm{x}_2^2 
            }
            \le d\cdot O(k^4C^2\log^2{k})
            \,.
            \yesnum\label{eq:2ndPrice:SecondMoment:sufficient}
        \end{align*}
        We divide the proof into three steps.

        \paragraph{Step 1 (Upper bound on $\Ex_z\sinsquare{
            \snorm{z - W_t^\top x}_2^2  \mid z\in P(\ysmax,\imax)
        }$).}
            Toward this, observe that 
            \[
                \Ex_z\insquare{
                    \snorm{z - W_t^\top x}_2^2  \mid z\in P(\ysmax,\imax)
                }
                \leq 
                2 \Ex_z\insquare{
                    \norm{z}_2^2 \mid z\in P(\ysmax,\imax)
                }
                + 2 \snorm{W_t^\top x}_2^2  \,.
            \]
            Substituting $v$ as $\wh{z}$ (\ie{}, the unit vector along $z$) in a simple analog of  \cref{lem:SS:LC:unconditionalMoments}\footnote{One can deduce this analog as follows: The proof of \cref{lem:SS:LC:unconditionalMoments} divides the moment into $k$-parts one corresponding to each coordinate and, further, breaks the conditioning $z\in P_{\ymax}$ into $k$-subparts corresponding to the $k$-parts of $P_{\ymax}$. One can repeat the same proof for $P(\ysmax,\imax)$ since it has $(k-1)$-parts and these parts have a similar structure as the parts of $P_{\ymax}$: they fix two coordinate's value and restrict the remaining $(k-2)$-coordinates to lie in an interval; see \cref{fig:2nd-price-observation-distinct}. One difference is that one of the slabs (which corresponds to $\imax$) pleases a constraint of the form $z_i\geq \ysmax$ instead of $z_i\leq \ysmax$, but this is easily handled by considering $w=-z_i$ and using that $\Ex[w^2]=\Ex[z^2]$ and $\Ex[w^4]=\Ex[z^4]$.\label{footnote:secondPriceBounds}} for the set $P(\ysmax,\imax)$, implies that 
            \[
                \Ex_z\insquare{
                    \norm{z}_2^2 \mid z\in P(\ysmax,\imax)
                }\leq 
                    2k\cdot 
                    {
                    \inparen{
                        1
                        +\abs{\ysmax}
                        +\snorm{W_t^\top x}_\infty
                    }^2
                    }
                    \leq 
                    4k\cdot 
                    {
                    \inparen{
                        1
                        +\ysmax^2
                        +\snorm{W_t^\top x}^2_\infty
                    }
                    }\,.
            \]
            Therefore, it follows that 
            \[
                \Ex_z\insquare{
                    \snorm{z - W_t^\top x}_2^2  \mid z\in P(\ysmax,\imax)
                }
                \leq 
                8k\cdot {
                    \inparen{
                        1
                        +\ysmax^2
                        +\snorm{W_t^\top x}^2_\infty
                    }
                 }
                + 2 \snorm{W_t^\top x}_2^2  \,.
                \yesnum\label{eq:2ndPrice:secondMoment:ub1}
            \]

        \paragraph{Step 2 (Upper bound on $\Ex_{\ysmax,\imax}\Ex_{z}[\snorm{z - W_t^\top x}_2^2 \mid z\in P(\ysmax,\imax)]$).}
            Next, we take the expectation with respect to $\ysmax,\imax$ and use \cref{fact:GaussMaxMoments} to upper bound $\Ex[\ysmax^2]$ to obtain
            \[
                \Ex_{\ysmax,\imax}
                \Ex_z\insquare{
                    \snorm{z - W_t^\top x}_2^2  \mid z\in P(\ysmax,\imax)
                }
                \leq 
                8k\cdot \inparen{
                    O\inparen{1+\log{k}}
                    + \snorm{W_t^\top x}_\infty^2
                 }
                + 2 \snorm{W_t^\top x}_2^2  \,.
                \yesnum\label{eq:2ndPrice:SecondMoment:1}
            \]


        \paragraph{Step 3 (Upper bound on $\Ex_{x}[\Ex_{\ysmax,\imax,z}[\snorm{z - W_t^\top x}_2^2 \mid z\in P(\ysmax,\imax)]\cdot \norm{x}_2^2]$).}
            Finally, we will take the expectation over $x$.
            For this, we will apply the transformation explained in Step 2 in the proof of \cref{lem:2ndPrice:LC:step2} (or rather \cref{lem:SS:LC:step2}).
            After this transformation, it follows that $W_t^\top x$ depends on the first $2k$ coordinates of $x$, \ie{}, on $x_{\leq 2k}$. 
            To make use of the above transformation, we split $\norm{x}_2^2$ between $x_{\leq 2k}$ and $x_{>2k}$ to yield
            \begin{align*}
                &\Ex_{x}\insquare{
                \Ex_{\ysmax,\imax}
                \Ex_z\insquare{
                    \snorm{z - W_t^\top x}_2^2  \mid z\in P(\ysmax,\imax)
                }
                \cdot 
                \norm{x}_2^2 
                }\\
                &\quad=\quad
                \Ex_{x_{>2k},~x_{\leq 2k}}\insquare{
                    \Ex_{\ysmax,\imax}
                    \Ex_z\insquare{
                        \snorm{z - W_t^\top x}_2^2  \mid z\in P(\ysmax,\imax)
                    }
                    \cdot 
                    \inparen{
                        \norm{x_{\leq 2k}}_2^2 
                        + 
                        \norm{x_{> 2k}}_2^2 
                    }
                }\,.
                \yesnum\label{eq:2ndPrice:SecondMoment:3}
            \end{align*}
            Observe that 
            \begin{align*}
                \snorm{W_t^\top x}_\infty 
                \leq \snorm{W_t^\top x}_2
                &\leq \sum_i \snorm{w^\star_i + t v_i}_2 \snorm{x_{\leq 2k}}_2
                \leq 
                {k(C+1)}\snorm{x_{\leq 2k}}_2
                ~~\Stackrel{(C\geq 1)}{\leq}~~ 2k^{3/2}C \snorm{x_{\leq 2k}}_\infty \,.
                \yesnum\label{eq:2ndPrice:SecondMoment:2}
            \end{align*}
            Hence,  \cref{eq:2ndPrice:SecondMoment:3,eq:2ndPrice:SecondMoment:2}, imply that 
            \begin{align*}
                &\Ex_{x}\insquare{
                \Ex_{\ysmax,\imax}
                \Ex_z\insquare{
                    \snorm{z - W_t^\top x}_2^2  \mid z\in P(\ysmax,\imax)
                }
                \cdot 
                \norm{x}_2^2 
                }\\
                &\quad=\quad
                {O(k^4C^2\log{k})}
                \cdot \Ex_{x_{>2k},~x_{\leq 2k}}\insquare{
                    {
                        \inparen{
                            1
                            + \norm{x_{\leq 2k}}_\infty^2
                        }
                        \cdot 
                        \inparen{
                            \norm{x_{\leq 2k}}_2^2 
                            + 
                            \norm{x_{> 2k}}_2^2 
                        }
                    }
                }\,.
            \end{align*}
            This can be further simplified to  
            \[
                {O(k^4C^2\log{k})}
                \cdot \Ex_{x_{>2k}}
                    \insquare{
                        \Ex_{x_{\leq 2k}}\insquare{
                            \snorm{x_{\leq 2k}}_\infty^2
                            + \snorm{x_{\leq 2k}}_\infty^4
                        }
                        + 
                        \norm{x_{> 2k}}_2^2 \cdot 
                        \Ex_{x_{\leq 2k}}\insquare{
                            1
                            + \snorm{x_{\leq 2k}}_\infty^2
                        }
                    }\,.
            \]
            Next, \cref{fact:GaussMaxMoments}, implies the following upper bound 
            \begin{align*}
                &\Ex_{x}\insquare{
                \Ex_{\ysmax,\imax}
                \Ex_z\insquare{
                    \snorm{z - W_t^\top x}_2^2  \mid z\in P(\ysmax,\imax)
                }
                \cdot 
                \norm{x}_2^2 
                }\\
                &\qquad\qquad\leq
                    {O(k^4C^2\log{k})}
                    \cdot \Ex_{x_{> 2k}}\insquare{
                    O\sinparen{\log^2{k}}
                    + \norm{x_{> 2k}}_2^2 \cdot O\inparen{\log{k}}
                }\,.
            \end{align*}
            Finally, using $\Ex_{x_{>2k}}[\norm{x_{>2k}}_2^2]=(d-2k)\,\Ex_{u\sim \normal{0}{1}}[u^2]=O(d-2k)$, we get that 
            \begin{align*}
                \Ex_{x}\insquare{
                \Ex_{\ysmax,\imax}
                \Ex_z\insquare{
                    \snorm{z - W_t^\top x}_2^2  \mid z\in P(\ysmax,\imax)
                }
                \cdot 
                \norm{x}_2^2 
                }
                &\leq
                (d-2k)\cdot {O(k^4C^2\log^3{k})}
                \leq d\cdot {O(k^4C^2\log^3{k})}\,.
            \end{align*}
            
        

    
\subsection{Projected Stochastic Gradient Descent}
    We are now ready to prove \Cref{thm:2ndPrice:localConvergence} by applying the iterative PSGD algorithm (\Cref{thm:SGD-local-growth-convergence}).
    We restate the theorem below for convenience.
    \thmSecondPriceLocalConvergence*

        \noindent The proof follows by putting together all the lemmas in \Cref{sec:2ndPrice} similar to the proof of \Cref{thm:SS:localConvergence}.



% --- end second_price.tex ---


\section{Coarse Gaussian Mean Estimation with Convex Partitions}
\label{sec:coarseGaussianEstimation}
% --- begin mean_estimation.tex ---
In this section, we provide our efficient algorithm for coarse Gaussian mean estimation under convex partitions.
We restate the desired result below for convenience.
\thmConvexPartitionMeanEstimation*

    In addition to information preservation (\cref{def:informationPreservation}), our algorithm relies on the following assumptions.
    \begin{assumption}[Convex Partitions]\label{asmp:convexPartitions}
        Each $P\in \hyP$ is a convex subset of $\R^d$.
    \end{assumption}
    The convexity of the sets will be crucial in order to argue about the convexity of the negative log-likelihood objective.
    \begin{assumption}[Bounded Mean]\label{asmp:boundedMean}
        We are given a parameter $D > 0$ such that $\muStar\in B(0, D)$ lies in the Euclidean ball of radius $D$.
    \end{assumption}
    We remark that our algorithm is able to handle the case $\muStar\in B(\mu^{(0)}, D)$ via a translation of the space.
    Moreover,
    our algorithm can tolerate a $\muStar\in B_\infty(0, D)$ but for simplicity,
    we state our results using the Euclidean distance.
    

    \begin{assumption}[Sampling Oracle]\label{asmp:samplingOracle}
        \sloppy
        There is an efficient sampling oracle that, 
        given $R > 0$,
        $P\in \hyP$,
        and a parameter $\mu\in \R^d$ describing a Gaussian distribution,
        outputs an unbiased sample $y\sim \truncatedNormal{\mu}{I}{P\cap B_\infty(0, R)}$.
    \end{assumption}
    {We emphasize that \Cref{asmp:samplingOracle} is \emph{not} a necessary assumption and we state it only for the sake of simplifying our presentation.
    Indeed,
    in order for all partitions to be convex,
    each $P\in \hyP$ must be a polyhedron.
    In this case,
    we can remove \Cref{asmp:samplingOracle} by implementing a sampling oracle that terminates in polynomial time with respect to the complexity of the input samples.
    See \Cref{apx:sampling-polytopes:removing-oracle-assumption} for more details.


\subsection{Projection Set}
Our algorithm is projected SGD. To this end, we define our
 projection set $K$ to be
\[
    K\coloneqq \inbrace{\mu\in \R^d\colon \norm{\mu}_2\leq D}.
    \yesnum\label{eq:mean:projectionSet}
\]
We certainly have $\mu^\star\in K$
and we can efficiently project onto $K$.

\subsection{Convexity and Local Growth of Log-Likelihood Under Convex Partitions}
It can be verified (\Cref{sec:meanEstimation:likelihoodGradientHessianCalculations}) that the coarse negative log-likelihood of the mean under canonical parameterization
is given by the following
\[
    \negLL_{\hyP}(\mu) = \sum_{P\in \hyP} \normalMass{\muStar}{I}{P} \negLL_{P}(\mu)\,,
    \quadwhere
    \negLL_{P}(\mu)\coloneqq -\log\inparen{
        \normalMass{\mu}{I}{P}
    }\,.
    \yesnum\label{eq:convexPartitionNLL}
\]
We also have the following expressions for the gradient and Hessian of $\negLL_\hyP$.
\begin{align*}
    \grad\negLL_\hyP\inparen{\mu}~~
    &=~~ 
        \mu
        -
        \Ex_{P\sim \coarseNormal{\muStar}{I}{\hyP}}
        {
            \Ex_{\truncatedNormal{\mu}{I}{P}}{
                \begin{bmatrix}
                    x
                \end{bmatrix}
            }
        }\,,\\
    \grad^2\negLL_\hyP\inparen{\mu}~~
    &=~~ 
        I
        - 
        \Ex_{P\sim \coarseNormal{\muStar}{I}{\hyP}}
        \cov_{\truncatedNormal{\mu}{I}{P}}{
            \begin{bmatrix}
                x
            \end{bmatrix}
        }
    \,.
\end{align*}
\cite{fotakis2021coarse} show that
$\nabla^2 \negLL_\hyP(\mu) \succeq 0$
when each $P\in \hyP$ is convex
by leveraging the Brascamp--Lieb Inequality (see \eg{}, \cite{guionnet2009large})
and showing that each $\negLL_P$ is convex.
\begin{proposition}[Lemma 15 in \cite{fotakis2021coarse}]\label{prop:convexPartitionVarianceReduction}
    Let $P\sset \R^d$ be convex.
    Then
    $
        \cov_{\truncatedNormal{\mu}{I}{P}}\insquare{
            x
        }
        \preceq I,
    $
    and in consequence,
    $\negLL_P(\mu)$ is convex as a function of $\mu$.
\end{proposition}
To recover $\muStar$,
one sufficient condition is to check that $\negLL_\hyP$ satisfies a local growth condition (\Cref{def:local-growth}). We show this in the next lemma.

\begin{lemma}\label{lem:convexPartitionLocalGrowth}
    Let $\hyP$ be a convex $\alpha$-information preserving partition of $\R^d$.
    The log-likelihood function $\negLL_\hyP(\mu)$ (\Cref{eq:convexPartitionNLL}) satisfies an $\inparen{\sqrt2\alpha, 2}$-local growth condition.
\end{lemma}

\begin{proof}
    Assume $\alpha$-information preservation (\Cref{def:model:informationPreservingDistortion}) holds.
    By \Cref{lem:model:informationPreservingImpliesGrowth},
    we have
    \begin{align*}
        \negLL_\hyP(\mu) - \negLL_\hyP(\muStar)
        &\geq \min\inparen{2, 2\alpha^2\norm{\muStar-\mu}_2^2}\,.
    \end{align*}
This concludes the proof.
\end{proof}

\subsection{Local Partitions}
\label{sec:localPartitions}
To be able to run PSGD, 
we need to obtain stochastic gradients for $\negLL_\hyP$
as well as bound the second moment of the stochastic gradients.
This is not straightforward since the inner and outer expectations in the gradient expressions
are not over the same means.
To overcome this, we first analyze an idealized class of partitions and derive an algorithm that recovers the mean
under this ideal situation.
Then we show that we can implement this algorithm using samples from the actual class of observed partitions.


\begin{definition}[$R$-Local Partition]\label{def:local partition}
    Let $R > 0$.
    We say a partition $\hyP$ of $\R^d$ is $R$-local if for any $P\in \hyP$ that is not a singleton,
    $P\sset B_\infty(0, R)$.
\end{definition}



\subsection{Stochastic Gradient Oracle and Second Moment} 
    \label{sec:SGD-local}\begin{lemma}\label{lem:convexPartitionGradientOracle}
        Suppose $\hyP$ is an $R$-local partition of $\R^d$ and $\norm{\muStar}_2\leq D$.
        Given $v\in K$ in the projection set (\Cref{eq:mean:projectionSet}),
        there is an algorithm that consumes a sample $P\sim \coarseNormal{\muStar}{I}{\hyP}$
        and makes a single call to the sampling oracle (\Cref{asmp:samplingOracle})
        to compute an unbiased estimate $g(\mu)$ of $\grad \negLL_\hyP(\mu)$ (\Cref{eq:convexPartitionNLL})
        such that $\E\insquare{\norm{g(\mu)}_2^2}= O(D^2 + dR^2)$.
    \end{lemma}
    Recall that the gradient of the log-likelihood function is of the form
    \[ 
        \mu
        -
        \Ex_{P\sim \coarseNormal{\muStar}{I}{\hyP}}
        {
            \Ex_{\truncatedNormal{\mu}{I}{P}}\insquare{
                x
            }
        }\,.
    \] 

    \begin{proof}
        The first term in the gradient expression is simply our current iterate for the mean
        and we can obtain an unbiased estimate of the second term by sampling $y\sim \truncatedNormal{\mu}{I}{P}$
        where $P$ is a fresh observation.

        The second moment of the gradient oracle can be upper bounded by
        \[
            2\norm{\mu}_2^2 
            +
            2\Ex_{P\sim \coarseNormal{\muStar}{I}{\hyP}}
            {
                \Ex_{\truncatedNormal{\mu}{I}{P}}\insquare{
                    \norm{x}_2^2
                }
            }\,.
        \]
        The first term is at most $D^2$
        and the second term can be upper bounded as follows.
        \begin{align*}
            &\Ex_{P\sim \coarseNormal{\muStar}{I}{\hyP}}
            {
                \Ex_{\truncatedNormal{\mu}{I}{P}}\insquare{
                    \norm{x}_2^2
                }
            } \\
            &\qquad= \E_{P\sim \coarseNormal{\muStar}{I}{\hyP}} \E_{x\sim \truncatedNormal{\mu}{I}{P}} \insquare{\norm{x}_2^2\cdot \ones_{B_\infty(0, R)}}
            + \E_{P\sim \coarseNormal{\muStar}{I}{\hyP}} \E_{x\sim \truncatedNormal{\mu}{I}{P}} \insquare{\norm{x}_2^2\cdot \ones_{B_\infty(0, R)^c}} \\
            &\qquad= \E_{P\sim \coarseNormal{\muStar}{I}{\hyP}} \E_{x\sim \truncatedNormal{\mu}{I}{P}} \insquare{\norm{x}_2^2\cdot \ones_{B_\infty(0, R)}}
            + \E_{P\sim \coarseNormal{\muStar}{I}{\hyP}} \E_{x\sim \truncatedNormal{\muStar}{I}{P}} \insquare{\norm{x}_2^2\cdot \ones_{B_\infty(0, R)^c}}\,.
        \end{align*}
        {In the last step,
        we use the fact that $\hyP$ is an $R$-local partition so that every set 
        such that $P\cap B_\infty(0, R)^c\neq \varnothing$ is a singleton.
        This allows us to replace the expectation over each $\truncatedNormal{\mu}{I}{P}$
        in the second term
        with an expectation over $\truncatedNormal{\muStar}{I}{P}$,
        as the distribution consists of a single point.}
        We can then bound the first term using a deterministic bound
        and the second term using the second moment of $\normal{\muStar}{I}$.
        \begin{align*}
            &\E_{P\sim \coarseNormal{\muStar}{I}{\hyP}} \E_{x\sim \truncatedNormal{\mu}{I}{P}} \insquare{\norm{x}_2^2\cdot \ones_{B_\infty(0, R)}}
            + \E_{P\sim \coarseNormal{\muStar}{I}{\hyP}} \E_{x\sim \truncatedNormal{\muStar}{I}{P}} \insquare{\norm{x}_2^2\cdot \ones_{B_\infty(0, R)^c}}  \\
            &\qquad\qquad\leq\quad  dR^2 + \E_{x\sim \normal{\muStar}{I}}\insquare{\norm{x}_2^2} \\
            &\qquad\qquad=\quad  dR^2 + d + \norm{\muStar}_2^2 \\
            &\qquad\qquad=\quad  O(dR^2 + D^2)\,.
        \end{align*}
    \end{proof}

\subsection{Projected Stochastic Gradient Descent}
We run the iterative PSGD algorithm for functions satisfying local growth (\Cref{alg:SGD-local-growth-convergence}).
The initial value bound can be taken as $\eps_0 = DG$ with any initial value $\mu_0\in B(0, D)$,
where $G^2 = O(D^2 + dR^2)$ is an upper bound on the second moment of the gradient oracle (\Cref{lem:convexPartitionGradientOracle}).
In order to recover $\muStar$ up to $\eps$-Euclidean distance,
we would like an $O(\alpha^2\eps^2)$-optimal solution in function value.
For $R$-local partitions,
the following result is obtained.
In order to obtain the high probability bound,
we repeat \Cref{alg:SGD-local-growth-convergence} a few times and apply the clustering trick described in \Cref{sec:SGD:clusteringBoosting}.

\begin{proposition}\label{thm:convexLocalPartitionMeanEstimation}
    Let $\eps\in (0, 1]$.
    Suppose $\hyP$ is a convex $R$-local $\alpha$-information preserving partition of $\R^d$ and $\norm{\muStar}_2\leq D$.
    There is an algorithm that outputs an estimate $\tilde \mu$ satisfying
    \[
        \norm{\tilde \mu - \muStar}_2 \leq \eps
    \]
    with probability $1-\delta$. Moreover,
    the algorithm requires
    \[
        m = O\inparen{\frac{dR^2 + D^2}{\alpha^4 \eps^2} \log^3\inparen{\frac{dRD}{\eps \delta}}}
    \]
    samples from $\coarseNormal{\muStar}{I}{\hyP}$ and $\poly(m, T_s)$ time,
    where $T_s$ is the time complexity of sampling from a Gaussian distribution truncated $P\cap B_\infty(0, R)$
    for some $P\in \hyP$.
\end{proposition}

\noindent We now describe how to remove the assumption of the partition being $R$-local.
\begin{proposition}\label{cor:convexPartitionMeanEstimation}
    Let $\eps\in (0, 1]$.
    Suppose $\hyP$ is a convex $\alpha$-information preserving partition of $\R^d$ and $\norm{\muStar}_2\leq D$.
    There is an algorithm that outputs an estimate $\tilde \mu$ satisfying
    \[
        \norm{\tilde \mu - \muStar}_2 \leq \eps
    \]
    with probability $1-\delta$. Moreover,
    the algorithm requires
    \[
        m = \tilde O\inparen{\frac{dD^2 \log(\nfrac1\delta)}{\alpha^4 \eps^2}}
    \]
    samples from $\coarseNormal{\muStar}{I}{\hyP}$ and $\poly(m, T_s)$ time,
    where $T_s$ is the time complexity of sampling from a Gaussian distribution truncated to a partition $P\in \hyP$.
\end{proposition}

\begin{proof}
    Consider a general $\alpha$-information preserving convex partition $\hyP$ and $P\sim \coarseNormal{\muStar}{I}{\hyP}$.
    Since $\norm{\muStar}_\infty\leq \norm{\muStar}_2\leq D$,
    setting $R = D + O(\log\nfrac{md}{\delta})$ means that any $m$-sample algorithm will not observe 
    a sample $P$ such that $P\cap B_\infty(0, R) = \varnothing$ with probability $1-\delta$.
    Define the partition $\hyP(R)$ of $\R^d$ given by
    \[
        \hyP(R)\coloneqq 
        \inbrace{P\cap B_\infty(0, R): P\in \hyP, P\cap B_\infty(0, R)\neq \varnothing}
        \cup \inbrace{\set{x}\colon x\notin B_\infty(0, R)}\,.
    \]
    Since $\hyP(R)$ is a refinement of $\hyP$,
    it must also be $\alpha$-information preserving.
    Consider the set-valued algorithm $F\colon  \hyP\to \hyP(R)$ given  by
    \[
        P \mapsto 
        \begin{cases}
            P\cap B_\infty(0, R), &P\cap B_\infty(0, R)\neq \varnothing\,, \\
            \text{$\inbrace{x}$ for an arbitrary $x\in P$}, &P\cap B_\infty(0, R) = \varnothing\,.
        \end{cases}
    \]
    By the choice of $R$,
    with probability $1-\delta$,
    any $m$-sample algorithm will not distinguish between 
    samples from $\coarseNormal{\muStar}{I}{\hyP(R)}$ and $F(P)$ for $P\sim \coarseNormal{\muStar}{I}{\hyP}$.
    Moreover,
    we can sample from $P\cap B_\infty(0, R)$ as per \Cref{asmp:samplingOracle}.
    Thus we can simply run the algorithm from \Cref{thm:convexLocalPartitionMeanEstimation}
    with input samples $F(P)$ for $P\sim \coarseNormal{\muStar}{I}{\hyP}$.
\end{proof}
We can further improve the sample complexity by running the algorithm in two stages.
The first stage aims to obtain an $O(1)$-distance warm start
and the second stage aims to recover the mean up to accuracy $\eps$.
This yields the proof of \Cref{thm:convexPartitionMeanEstimation}.
Indeed,
the only missing detail is an efficient implementation of a sampling oracle (\Cref{asmp:samplingOracle}).
As mentioned,
we defer these details to \Cref{apx:sampling-polytopes:removing-oracle-assumption}.

    
% --- end mean_estimation.tex ---






 
\section{PSGD Convergence for Convex Functions Satisfying a Local Growth Condition}
\label{sec:sgd-analysis}
% --- begin quadratic_growth.tex ---
In this section,
we state and prove the analysis for the variant of gradient descent we use to optimize the various likelihood functions in our work.
Consider a convex function $F\colon  K\to \R$ with a global minimizer $w^\star$
on a convex subset $K\sset \R^d$.
We write
\[
    S_\rho \coloneqq \inbrace{w\in K\colon F(w) - F(w^\star)\leq \rho}
\]
to denote the \emph{$\eps$-sublevel set} of a function $F$
where $F$ is clear from context.

We now state a useful definition.
\begin{definition}[$\eta$-Local Growth Condition]\label{def:local-growth}
    We say that $F\colon K\to \R$ satisfies an \emph{$(\eta, \rho)$-local growth condition} if
    \[
        \norm{w - w^\star}_2
        \leq \frac{(F(w) - F(w^\star))^\frac12}{\eta}
    \]
    for every $w\in S_{\rho}$.
\end{definition}
We remark that a function satisfying a $(\eta_0, \rho_0)$-local growth condition
also satisfies $(\eta, \rho)$-local growth for every $\eta\in (0, \eta_0], \rho\in (0, \rho_0]$.
We also note that by \Cref{thm:quadGrowth},
our log-likelihood functions for $\alpha$-information preserving distortion mechanisms at radius $R$ satisfy a $(\sqrt2 \alpha, 2{\alpha^2})$-local growth condition
{if $K\sset B_R(w^\star)$}.

Recall the following standard convergence result from convex optimization.
\begin{theorem}[Theorem 9.7 in \cite{garrigos2023handbook}]\label{thm:SGD-convergence}
    Let $F\colon S\sset \R^d\to \R$ be convex and $w\in S$.
    Suppose we have acccess to an unbiased gradient oracle $g(w)$ such that $\E\insquare{\norm{g(w)}_2^2}\leq G^2$.
    Then PSGD with initial point $w^{(0)}\in \R^d$
    and constant step-size $\gamma > 0$ 
    outputs an average iterate $\bar w$ satisfying
    \[
        \E\insquare{F(\bar w) - F(w)}
        \leq \frac{\gamma G^2}{2} + \frac{\snorm{w^{(0)} - w}_2^2}{2\gamma T}\,.
    \]
\end{theorem}
One straightforward way to recover the parameters from optimizing a function satisfying a $(\eta, \eps)$-local growth condition
is to find an $\eta^2\eps^2$-optimal point using \Cref{thm:SGD-convergence}.
This yields a $O(\nfrac1{\eta^4 \eps^4})$ rate of convergence which may not be desirable.
We design a novel iterative refinement algorithm that is able to leverage the local growth condition in a more clever way.
Our algorithm is inspired by \citet{xu2019accelerate} but their results only hold under the restrictive condition on the gradient oracle that there is a deterministic bound $\norm{g(w)}\leq G$.
This assumption is necessary to obtain high probability bounds in their algorithm.
On the other hand,
our algorithm will converge with constant probability
but relies on a much weaker bound on the second moment of the gradient oracle.

We begin with a key insight due to \citet{yang2018beating} as stated in \cite{xu2019accelerate}.
For the sake of analysis only,
it will be useful to define the following notation with respect to $F$.
\[
    w_\rho^\dagger
    \coloneqq \argmin_{v\in S_\rho} \norm{v-w}_2^2\,.
\]
\begin{proposition}[Lemma 1 in \cite{xu2019accelerate}]\label{prop:level-set-dist}
    Suppose $F\colon \R^d\to \R$ satisfies a $(\eta, \rho)$-local growth condition.
    Then for any $w\in \R^d$,
    \[
        \snorm{w-w_\rho^\dagger}_2
        \leq \frac{F(w) - F(w_\rho^\dagger)}{\eta \sqrt{\rho}}\,.
    \]
\end{proposition}
In light of \Cref{prop:level-set-dist},
we see now that it is possible to run PSGD for a few iterations,
then convert the improvement in function value to an improvement in parameter distance to the $\eps$-sublevel set.
We can then restart PSGD from this improved point.

\begin{algorithm}[t]
\caption{Iterative-PSGD($K, w_0, \epsilon_0, g, \eps$)}\label{alg:SGD-local-growth-convergence}
\begin{algorithmic}[1]
    \State \textbf{Input:} Projection access to feasible region $K$,
    initial point $w^{(0)}\in K$, 
    $\epsilon_0\geq F(w^{(0)}) - F(\wstar)$,
    gradient oracle $g$ with $\E\insquare{\norm{g}_2^2}\leq G^2$,
    local growth rate $\eta > 0$,
    desired accuracy $\eps > 0$
    \State \textbf{Output:} $2\eps$-optimal solution with probability $0.99$
    \vspace{0.2em}
    \State Set $\tau \gets \ceil{\log_2(\nfrac{\eps_0}\eps)}$
    \vspace{0.2em}
    \State Set $D_0 \gets \frac{2\eps_0}{\eta \sqrt{\eps}}$
    \vspace{0.2em}
    \State Set $\gamma_0 \gets \frac{\eps_0}{100G^2\tau}$
    \vspace{0.2em}
    \State Set $T \gets \frac{40000 G^2 \tau^2}{\eta^2\eps}$
    \vspace{0.2em}
    \For{$\ell=1,\ldots, \tau$}
        \State $\gamma_\ell \gets 2^{-\ell}\gamma_0$
        \State $D_\ell\gets 2^{-\ell}D_0$
        \State $w^{(\ell, 0)} \gets w^{(\ell-1)}$
        \For{$t=1, \dots, T$}
            \State $w^{(\ell, t)} \gets \Pi_{K\cap B(w^{(\ell-1)}, D_k)} \inparen{w^{(\ell, t-1)} - \gamma_\ell\cdot g(w^{(\ell, t-1)})}$
        \EndFor
        \State $w^{(\ell)} \gets \frac1T \sum_{t=1}^T w^{(\ell, t)}$
    \EndFor
    \vspace{1em}
    \State \textbf{Return} $w^{(\tau)}$
\end{algorithmic}
\end{algorithm}


\begin{theorem}\label{thm:SGD-local-growth-convergence}
    Let $F\colon S\sset \R^d\to \R$ be convex and satisfy a $(\eta, \eps)$-local growth condition (\Cref{def:local-growth}).
    Suppose we have acccess to an unbiased gradient oracle $g(w)$ such that $\E\insquare{\norm{g(w)}_2^2}\leq G^2$.
    Suppose further that we have access to an $\eps_0$-optimal solution $w^{(0)}\in S$.
    \Cref{alg:SGD-local-growth-convergence} queries the gradient oracle
    \[
        O\inparen{\frac{G^2}{\eta^2 \eps} \cdot \log^3\inparen{\frac{\eps_0}{\eps}}}
    \]
    times and outputs a $2\eps$-optimal solution with probability $0.99$.
\end{theorem}

\begin{proof}
    Define $\eps_\ell \coloneqq 2^{-\ell} \eps_0$.
    We will argue that $F(w^{(\ell)}) - F(w^\star) \leq \eps_{\ell} + \eps$
    with probability $1-\nfrac1{100\tau}$ conditioned on $F(w^{(\ell-1)}) - F(w^\star) \leq \eps_{\ell-1} + \eps$.
    The claim follows then by a union bound over $\tau$ stages.

    The base case for $\ell=0$ holds by assumption.
    We consider some $\ell\geq 1$.
    Remark that
    \[
        D_\ell = \frac{\eps_{\ell-1}}{\eta \sqrt{\eps}}
        \qquadand
        \gamma_\ell = \frac{\eps_\ell}{100G^2\tau}\,.
    \]
    By \Cref{prop:level-set-dist},
    we have
    \[
        \norm{(w^{(\ell-1)})_\eps^\dagger - w^{(\ell-1)}}_2
        \leq \frac{1}{\eta \sqrt\eps} (F(w^{(\ell-1)}) - F(w^{(\ell-1)})_\eps^\dagger)
        \leq \frac{\eps_{\ell-1}}{\eta\sqrt\eps}
        \leq D_\ell\,.
    \]
    Then \Cref{thm:SGD-convergence} on the $\ell$-th stage of PSGD yields
    \[
        \E[F(w^{(\ell)}) - F((w^{(\ell-1)})_\eps^\dagger)]
        \leq \frac{\gamma_\ell G^2}{2} + \frac{D_\ell^2}{2\eta_\ell T}
        = \frac{\eps_\ell}{200\tau} + \frac{400G^2\tau \eps_\ell}{2\eta^2 \eps T}
        = \frac{\eps_\ell}{100\tau}\,.
    \]
    An application of Markov's inequality yields the desired result.
\end{proof}
In order to recover the parameters up to $\eps$-accuracy by optimizing the log-likelihood function using \Cref{thm:SGD-local-growth-convergence},
we need $O(\alpha^2\eps^2)$-optimal solutions.
This incurs a sample complexity of
\[
    O\inparen{\frac{G^2}{\alpha^4 \eps^2} \cdot \log^3\inparen{\frac{\eps_0}{\eps}}}
\]
which matches the sample-complexity of PSGD even when the function is $\alpha^2$-strongly convex over the entire feasible region,
up to logarithmic terms.
We note that that $F$ is at most $G$-Lipschitz.
Hence if we are given a warm-start in distance,
we can set $\eps_0\leq G\norm{w^{(0)} - w^\star}_2$.
Thus we may alternatively present the sample complexity above as
\[
    O\inparen{\frac{G^2}{\alpha^4 \eps^2} \cdot \log^3\inparen{\frac{G}{\eps}\cdot \snorm{w^{(0)} - w^\star}_2}}\,.
\]

\begin{remark}[Boosting without Function Evaluation via Clustering]
    \label{sec:SGD:clusteringBoosting}
    In order to boost the probability of successfully recovering $\wstar$,
    a standard trick is to repeat the algorithm $O(\log(\nfrac1\delta))$ times and output the solution with smallest objective value.
    However,
    we do not have exact evaluation access to the log-likelihood functions we wish to optimize.
    \cite[Section 3.4.5]{daskalakis2018efficient} demonstrate a ``clustering'' trick to avoid function evaluation.
    Suppose we repeat the algorithm in \Cref{thm:SGD-convergence} $O(\log(\nfrac1\delta))$ times.
    A Chernoff bound yields that with high probability,
    at least $\nfrac23$ of the outputted points are $\eps$-close to $\wstar$
    and thus are $2\eps$-close to each other.
    Thus outputting any point which is at most $2\eps$-close to at least 50\% of the points must be at most $3\eps$-close to $\wstar$.
\end{remark}







% --- end quadratic_growth.tex ---



    
    
\ifanon
\else
\section*{Acknowledgements}
    We thank Manolis Zampetakis for feedback on an initial draft of this work which helped improve its presentation, discussions about the literature on estimation from auction data, and for sharing the reference to \citet{meilijson1981autopsy}.
    Alkis Kalavasis was supported by the Institute for Foundations of Data Science at Yale.
    Felix Zhou acknowledges the support of the Natural Sciences and Engineering Research Council of Canada (NSERC).
\fi

\newpage
\printbibliography

\newpage

\appendix 

\addtocontents{toc}{\protect\setcounter{tocdepth}{1}}

\section{Computations of Gradients of Negative Log-Likelihoods}\label{sec:hessian}
    \subsection{Linear Regression with Self-Selection Bias}\label{sec:hessian:SS}
        Consider an estimate $W\in \R^{d\times k}$ of the parameter $W^\star$ defining the self-selection instance.
        Given an observation $(x,\ymax)$, the likelihood of the estimate $W$ (conditioned on a fixed value of $x$ and $\ymax$) can be written as the likelihood of observing ($x,\ymax$).
        In particular,
        {ignoring normalizing constants,
        we have}
        \[
            {\Pr_{z\sim \normal{W^\top x}{I}} \insquare{\max_i z_i = \ymax}
            \propto} 
            \int_{\max_i {z_i} = \ymax} \exp\inparen{-\frac{1}{2}\snorm{z- W^\top x}_2^2} \d z\,.
        \]
        Here, the integral is over a set of the following form:
        given a value $v\in \R$, define
        \[
            P(v)\coloneqq \inbrace{z\in \R^d\colon \max_i z_i = v}\,.
        \]
        Hence, the negative log-likelihood of observing $(x,\ymax)$ is 
        \[
            \cL\inparen{W; x,\ymax} = -\log\int_{P({\ymax})} \exp\inparen{-\frac{1}{2}\snorm{z- W^\top x}_2^2} \d z\,.
        \]
        Thus, the population negative log-likelihood is
        \[
            \cL\inparen{W} 
            = \Ex_{x,\ymax}\cL\inparen{W; x,\ymax} 
            = -\Ex_{x,\ymax}\log\int_{P({\ymax})} \exp\inparen{-\frac{1}{2}\snorm{z- W^\top x}_2^2} \d z\,.
        \]
        One of the most important properties of the negative log-likelihood is that the true parameters $\Wstar$ is a stationary point of $\cL(\cdot)$.
        \begin{restatable}[]{fact}{SSStationary}\label{fact:SS:stationaryPoint}
            It holds that $\grad \cL(W^\star)=0$.    
        \end{restatable}
        \paragraph{Gradient of $\cL(\cdot)$.}
        Next, we compute $\grad\cL(\cdot)$, which, in particular, enables us to verify \cref{fact:SS:stationaryPoint}.
        \begin{fact}[Direct Calculation]
            It holds that 
            \[
                \grad \cL(W)
                = W - \Ex_{x, \ymax}
                    {
                        \Ex_{z\sim \normal{W^\top x}{I} {\mid z\in P(\ymax)}}\insquare{xz^\top }
                    }
                    \,.
                \yesnum\label{eq:SS:gradientExpression}
            \]
        \end{fact}
        \begin{proof}
            The expression follows by verifying the following expressions via direct computation and using that $\Ex[xx^\top]=I$
            \begin{align*}
                \grad \cL(W; x,\ymax)
                    &~~=~~ \frac{1}{2}\Ex_{z\sim \normal{W^\top x}{I} \mid z\in P(\ymax)}{
                        \grad_W \snorm{z-W^\top x}_2^2
                    }\,,\\
                \grad_W \snorm{z-W^\top x}_2^2 
                    &~~=~~ -2\inparen{xz^\top -xx^\top W}\,.
            \end{align*}
        \end{proof}
        Now we are ready to prove \cref{fact:SS:stationaryPoint}.
        \begin{proof}[Proof of \cref{fact:SS:stationaryPoint}]
            This holds because $\Pr[\ymax = m] = \normalMass{{W^\star}^\top x}{I}{{P(m)}}$ and, conditioned on $(x,\ymax)$, 
            \[
                \Pr[z=v\mid x,\ymax]= \frac{\mathds{1}_{P({\ymax})}(v)}{\normalMass{{W^\star}^\top x}{I}{P({\ymax})}}\cdot \normalMass{{W^\star}^\top x}{I}{v}\,.  
            \]
            Therefore, it follows that 
            \[
                \Pr[z=v\mid x] = \normalMass{{W^\star}^\top x}{I}{v}\,.  
            \]
            Substituting this in the expression for $\grad\cL(\cdot)$ implies that
            \[
                \grad \cL(W^\star) 
                = W^\star - \Ex_x \insquare{x \inparen{\Ex_{z\sim \normal{\Wstar^\top x}{I}} z^\top}}
                = \Wstar - \Ex_{x}\inparen{ xx^\top W^\star} = 0\,.
            \]
        \end{proof}
        
        \paragraph{Hessian of $\cL(\cdot)$.}
        Next, we compute the Hessian of $\cL(\cdot)$. 
        Differentiating $\grad \cL(\cdot)$ once more, gives us the following expression
        \begin{align*}
            \grad^2 \cL(W)
                &= \grad_W \inparen{
                    W - \Ex_{x, \ymax}
                        {
                            \Ex_{z\sim \normal{W^\top x}{I} \mid z\in P(\ymax)}\insquare{xz^\top }
                        } 
                }
                \\
                &= I_{dk}
                    -
                    \Ex_{x,\ymax}\grad_W \inparen{
                            \Ex_{z\sim \normal{W^\top x}{I} }\insquare{
                        xz^\top\mid z\in P({\ymax})
                    } 
                }
                \,.
                \yesnum\label{eq:hessian:SS:1}
        \end{align*}
        Next, for a fixed $(x,\ymax)$,  we compute the above gradient.
        To simplify the calculations, let us define (for fixed $x$ and $\ymax$)
        \[
            {f(W) = f(W; z)}\coloneqq \snorm{z - W^\top x}_2^2\,.
        \]
        {For the purpose of this calculation,
        we identiy $f\colon \R^{dk}\to \R$ as a function of the flattened vector.}
        Now, we can re-write the aforementioned flattened expectation 
        \[
            \Ex_{z\sim \normal{W^\top x}{I} }\insquare{
                        (xz^\top)^\flat\mid z\in P({\ymax})
            }
            = \frac{
                \int_{P({\ymax})} \inparen{xz^\top}^\flat e^{-\frac{1}{2}f(W)} \d z
            }{
                \int_{P({\ymax})} e^{-\frac{1}{2}f(W)} \d z
            }\,.
        \]
        With this expression in hand, we can compute the gradient as follows
        \begin{align*}
           &\grad_W \frac{
            \int_{P({\ymax})} (xz^\top)^\flat e^{-\frac{1}{2}f(W)} \d z 
           }{
            \int_{P({\ymax})} e^{-\frac{1}{2}f(W)} \d z 
           }\\
           &\quad= 
           -\frac{
               \int_{P({\ymax})} (xz^\top)^\flat \grad f(W)^\top e^{-\frac{1}{2}f(W)} \d z 
           }{
                2\int_{P({\ymax})} e^{-\frac{1}{2}f(W)} \d z 
           }
           +
           \frac{
                \inparen{
                    \int_{P({\ymax})} (xz^\top)^\flat e^{-\frac{1}{2}f(W)} \d z} 
                \inparen{
                    \int_{P({\ymax})} \grad f(W)^\top e^{-\frac{1}{2}f(W)} \d z 
                }
           }{
                2\inparen{\int_{P({\ymax})} e^{-\frac{1}{2}f(W)} \d z }^2
           }
           \,.
        \end{align*}
        Substituting $\grad f(W) = 2(xx^\top W- xz^\top)$ implies that 
        \begin{align*}
           \grad_W \frac{
            \int_{P({\ymax})} {(xz^\top)}^\flat e^{-\frac{1}{2}f(W)} \d z 
           }{
            \int_{P({\ymax})} e^{-\frac{1}{2}f(W)} \d z 
           }
           &= 
           \Ex_z\insquare{{(xz^\top)}^\flat {\sinparen{(xz^\top)^\flat - (xx^\top W)^\flat}}^\top \mid z\in P({\ymax})}\\
           &\quad - \Ex_z\insquare{(xz^\top)^\flat \mid z\in P({\ymax})}
           ~\Ex_z\insquare{
                {(xz^\top)^\flat - (xx^\top W)^\flat}
                \mid z\in P({\ymax})
           }^\top\,.
        \end{align*}
        Since $xx^\top W$ is a constant with respect to $z$, the above simplifies to 
        \begin{align*}
            \grad_W \frac{
             \int_{P({\ymax})} {(xz^\top)}^\flat e^{-\frac{1}{2}f(W)} \d z 
            }{
             \int_{P({\ymax})} e^{-\frac{1}{2}f(W)} \d z 
            }
            &= \Cov\insquare{{(xz^\top)}^\flat\mid z\in P({\ymax})} \\
            &= \Cov\insquare{z\otimes x\mid z\in P({\ymax})} \\
            &= \Cov\insquare{z\mid z\in P({\ymax})}\otimes xx^\top
            \,.
        \end{align*}
        Substituting this in \cref{eq:hessian:SS:1} implies that 
        \[
            \grad^2 \cL(W) =  I_{dk} 
                - \Ex_{x,~\ymax}\insquare{\Cov\insquare{z\mid z\in P({\ymax})}\otimes xx^\top}\,.
        \]
         
    

        

    
    
    \subsection{Linear Regression with Second-Price Auction Data}\label{sec:hessian:2ndPrice}
        Consider an estimate $W\in \R^{d\times k}$ of the parameter $W^\star$ defining the second-price auction model (\cref{def:secondPrice}).
        For clarity, in this section, we use $(x,\ysmax,\imax)$ to denote observations $(x, y, i)$ in the second-price auction model.
        Given values $(\ysmax,\imax)$, define the following set 
        \[
            P(\ysmax, \imax)
            \coloneqq 
            \inbrace{
                z\in \R^d\colon 
                \imax\in \argmax_i z_i 
                    ~~\text{and}~~
                \ysmax = \max_{i\neq \imax} z_i
            }\,.
        \]
        Given an observation $(x,\ysmax,\imax)$, the likelihood of the estimate $W$ (conditioned on a fixed value of $x$, $\ysmax$, and $\imax$) can be written as the likelihood of observing ($x,\ysmax,\imax$).
        In particular,
        again ignoring the normalizing constant,
        we have
        \[
            \Pr_{z\sim \normal{W^\top x}{I}}
                \insquare{
                    \imax\in \argmax_i z_i
                    ~~\text{and}~~
                    \ysmax {=} \max_{i\neq \imax} z_i
                }
            \propto
            \int_{P(\ysmax, \imax)}
                \exp\inparen{-\frac{1}{2}\snorm{z- W^\top x}_2^2} \d z\,.
        \]
        Hence, the negative log-likelihood of observing $(x,\ysmax,\imax)$ is 
        \[
            \cL\inparen{W; x,\ysmax,\imax} = -\log\int_{P({\ysmax,\imax})} \exp\inparen{-\frac{1}{2}\snorm{z- W^\top x}_2^2} \d z
        \]
        and the population negative log-likelihood is
        \[
            \cL\inparen{W} 
            = \Ex_{x,\ysmax,\imax}\cL\inparen{W; x,\ysmax,\imax} 
            = -\Ex_{x,\ysmax,\imax}\log\int_{P({\ysmax,\imax})} \exp\inparen{-\frac{1}{2}\snorm{z- W^\top x}_2^2} \d z\,.
        \]
        \paragraph{Gradient of $\cL(\cdot)$.}
        Next, we compute $\grad\cL(\cdot)$. %
        \begin{fact}[Direct Calculation]
            It holds that 
            \[
                \grad \cL(W)
                = W - \Ex_{x, \ymax}
                    {
                        \Ex_{z\sim \normal{W^\top x}{I} \mid z\in P(\ysmax,\imax)}\insquare{xz^\top }
                    }
                    \,.
                \yesnum\label{eq:2ndPrice:gradientExpression}
            \]
        \end{fact}
        \begin{proof}
            The expression follows by verifying the following expressions via direct computation and using that $\Ex[xx^\top]=I$
            \begin{align*}
                \grad \cL(W; x,\ysmax,\imax)
                    &~~=~~ \frac{1}{2}\Ex_{z\sim \normal{W^\top x}{I} {\mid z\in P(\ysmax,\imax)}}{
                        \grad_W \snorm{z-W^\top x}_2^2
                    }\,,\\
                \grad_W \snorm{z-W^\top x}_2^2 
                    &~~=~~ -2\inparen{xz^\top -xx^\top W}\,.
            \end{align*}
        \end{proof}
        In particular, substituting $W=\Wstar$ implies the following.
        \begin{restatable}[]{fact}{2ndPriceStationary}\label{fact:2ndPrice:stationaryPoint}
            It holds that $\grad \cL(W^\star)=0$.    
        \end{restatable}
        \begin{proof}[Proof of \cref{fact:2ndPrice:stationaryPoint}]
            This holds because $\Pr[\ysmax = m, \imax=i] = \normalMass{{W^\star}^\top x}{I}{{P(m,i)}}$ and, conditioned on $(x,\ymax)$, 
            \[
                \Pr[z=v\mid x,\ysmax,\imax]= \frac{\mathds{1}_{P({\ysmax,\imax})}(v)}{\normalMass{{W^\star}^\top x}{I}{P({\ysmax,\imax})}}\cdot \normalMass{{W^\star}^\top x}{I}{v}\,.  
            \]
            Therefore, it follows that 
            \[
                \Pr[z=v\mid x] = \normalMass{{W^\star}^\top x}{I}{v}\,.  
            \]
            Substituting this in the expression for $\grad\cL(\cdot)$ implies that
            \[
                \grad \cL(W^\star) 
                = W^\star - \Ex_x \insquare{x \inparen{\Ex_{z\sim \normal{\Wstar^\top x}{I}} z^\top}}
                = \Wstar - \Ex_{x}\inparen{ xx^\top W^\star} = 0\,.
            \]
        \end{proof}
        
        \paragraph{Hessian of $\cL(\cdot)$.}
        Next, we compute the Hessian of $\cL(\cdot)$.
        Differentiating $\grad \cL(\cdot)$ once more, gives us the following expression
        \begin{align*}
            \grad_W^2 \cL(W)
                &= I_{dk}
                    -
                    \Ex_{x,\ysmax,\imax}\grad_W \inparen{
                            \Ex_{z\sim \normal{W^\top x}{I} }\insquare{
                        xz^\top\mid z\in P({\ysmax,\imax})
                    } 
                }
                \,.
                \yesnum\label{eq:hessian:2ndPrice:1}
        \end{align*}
        Next, for a fixed $(x,\ysmax,\imax)$,  we compute the above gradient.
        To simplify the calculations, let us define (for fixed $x$, $\ysmax$, and $\imax$)
        \[
            f(W; z)\coloneqq \snorm{z - W^\top x}_2^2\,.
        \]
        Again, for the sake of this computation,
        we think of $f\colon \R^{dk}\to \R$ as a function of the flattened vector.
        Now, we can re-write, the aforementioned expectation 
        \[
            \Ex_{z\sim \normal{W^\top x}{I} }\insquare{
                        xz^\top\mid z\in P({\ysmax,\imax})
            }
            = \frac{
                \int_{P({\ysmax,\imax})} \inparen{xz^\top}^\flat e^{-\frac{1}{2}f(W)} \d z
            }{
                \int_{P({\ysmax,\imax})} e^{-\frac{1}{2}f(W)} \d z
            }\,.
        \]
        With this expression in hand, we can compute the gradient as follows
        \begin{align*}
           &\grad_W \frac{
            \int_{P({\ysmax,\imax})} (xz^\top)^\flat e^{-\frac{1}{2}f(W)} \d z 
           }{
            \int_{P({\ysmax,\imax})} e^{-\frac{1}{2}f(W)} \d z 
           }\\
           &\qquad= 
           -\frac{
               \int_{P({\ysmax,\imax})} (xz^\top)^\flat \grad f(W)^\top e^{-\frac{1}{2}f(W)} \d z 
           }{
                2\int_{P({\ysmax,\imax})} e^{-\frac{1}{2}f(W)} \d z 
           }\\
           &\qquad\qquad\qquad+
           \frac{
                \inparen{
                    \int_{P({\ysmax,\imax})} (xz^\top)^\flat e^{-\frac{1}{2}f(W)} \d z} 
                \inparen{
                    \int_{P({\ysmax,\imax})} \grad f(W)^\top e^{-\frac{1}{2}f(W)} \d z 
                }
           }{
                2\inparen{\int_{P({\ysmax,\imax})} e^{-\frac{1}{2}f(W)} \d z }^2
           }
           \,.
        \end{align*}
        Substituting $\grad f(W) = 2(xx^\top W- xz^\top)$ implies that 
        \begin{align*}
           &\grad_W \frac{
            \int_{P({\ysmax,\imax})} {(xz^\top)}^\flat e^{-\frac{1}{2}f(W)} \d z 
           }{
            \int_{P({\ysmax,\imax})} e^{-\frac{1}{2}f(W)} \d z 
           }\\
           &\qquad= 
           \Ex_z\insquare{{(xz^\top)}^\flat {\sinparen{(xz^\top)^\flat - (xx^\top W)^\flat}}^\top \mid z\in P({\ysmax,\imax})}\\
           &\qquad\qquad\qquad - \Ex_z \insquare{
                {(xz^\top)}^\flat \mid z\in P({\ysmax,\imax})
            }
           ~\Ex_z\insquare{
                (xz^\top)^\flat - (xx^\top W)^\flat
                \mid z\in P({\ysmax,\imax})
           }^\top\,.
        \end{align*}
        Since $xx^\top W$ is a constant with respect to $z$, the above simplifies to 
        \begin{align*}
            \grad_W \frac{
             \int_{P({\ysmax,\imax})} {(xz^\top)}^\flat e^{-\frac{1}{2}f(W)} \d z 
            }{
             \int_{P({\ysmax,\imax})} e^{-\frac{1}{2}f(W)} \d z 
            }
            &= \Cov\insquare{{(xz^\top)}^\flat\mid z\in P({\ysmax,\imax})}\\
            &= \Cov\insquare{z\mid z\in P({\ysmax,\imax})}\otimes xx^\top
            \,.
        \end{align*}
        Substituting this in \cref{eq:hessian:2ndPrice:1} implies that 
        \[
            \grad^2 \cL(W) =  I_{dk} 
                - \Ex_{{x,\ysmax,\imax}}\insquare{\Cov\insquare{z\mid z\in P({\ysmax,\imax})}\otimes xx^\top}\,.
        \]

    

    \subsection{Gaussian Mean Estimation}\label{sec:meanEstimation:likelihoodGradientHessianCalculations} 
    Let $\negLL\colon \R^{d}\to \R_{\geq 0}$ denote the negative log-likelihood function for an instance of the coarse Gaussian mean estimation problem.
    $\negLL$ is defined as follows
    \[
        \negLL\inparen{\mu}
        ~\coloneqq~
        \Ex_{P\sim \coarseNormal{\muStar}{I}{\hyP}}
        \insquare{
            -\log\inparen{\normalMass{\mu}{I}{P}}
        }
        \,.
    \]
    To show that $\negLL$ is strongly convex around $\inparen{\muStar,\SigmaStar}$, we compute its derivatives.
    \begin{fact}\label{fact:gradientsNLL}
        It holds that 
        \begin{align*}
            \grad\negLL{}\inparen{\mu}~~
            &=~~ 
                \Ex_{\normal{\mu}{I}}{
                    [x]
                }
                -
                \Ex_{P\sim \coarseNormal{\muStar}{I}{\hyP}}
                {
                    \Ex_{\truncatedNormal{\mu}{I}{P}}{
                        [x]
                    }
                }\,,\\
            \grad^2\negLL{}\inparen{\mu}~~
            &=~~ 
                \cov_{\normal{\mu}{I}}{
                    [x]
                }
                - 
                \Ex_{P\sim \coarseNormal{\muStar}{I}{\hyP}}
                \cov_{\truncatedNormal{\mu}{I}{P}}{
                    [x]
                }
            \,,
        \end{align*}
    \end{fact}
    \begin{proof}
        We can write the log-likelihood as follows
        \[
            \negLL(\mu) = \sum_{P\in \hyP} \normalMass{\muStar}{I}{P} \negLL_{P}(\mu)\,,
            \quadwhere
            \negLL_{P}(\mu)\coloneqq -\log\inparen{
                \normalMass{\mu}{I}{P}
            }\,.
            \yesnum\label{eq:gradOfNLL:expansion}
        \]
        Due to the linearity of gradients, it suffices to compute $\grad\negLL_P(v, T)$ and $\grad^2\negLL_P(\mu)$ for each $P\in \hyP$
        to obtain the gradient and Hessian of $\negLL(\cdot)$.
        Toward this, fix any $P\in \hyP$.
        Observe that 
        \[
            \negLL_{P}(\mu) = 
            \log\frac{
                \int_{x\in \R^d} e^{
                    -\frac{1}{2}(x-\mu)^\top(x-\mu)
                }\d x
            }{
                \int_{x\in P} e^{
                    -\frac{1}{2}(x-\mu)^\top(x-\mu)
                }\d x
            }
            = 
            \log\frac{
                \int_{x\in \R^d}
                e^{
                    -\frac12 \norm{x}_2^2 + x^\top \mu
                }\d x
            }{
                \int_{x\in P}
                e^{
                    -\frac12 \norm{x}_2^2 + x^\top \mu
                }\d x
            }
            \,.
        \]
        Write $f(x; \mu) \coloneqq \exp(-\frac12 \norm{x}_2^2 + x^\top \mu)$. 
        It follows that 
        \begin{align*}
            \grad_\mu \negLL_P(\mu)
            &= \frac{\int_P f}{\int_{\R^d} f}\cdot \left[ \frac{(\int_{\R^d} \grad f) (\int_P f) - (\int_{\R^d} f)(\int_P \grad f)}{(\int_P f)^2} \right] \\
            &= \frac{\int_{\R^d} \grad f}{\int_{\R^d} f} - \frac{\int_P \grad f}{\int_P f}\,.
            \yesnum\label{eq:gradOfNLL:compressed}
        \end{align*}
        Simplifying the expression and substituting the values of $f, \grad f$ gives
        \[
            \grad \negLL_{P}(\mu) 
            = 
            \Ex_{\normal{\mu}{I}}\insquare{
                    x
            }
            - 
            \Ex_{\truncatedNormal{P}{\mu}{I}}\insquare{
                    x
            }\,.
        \]
        Substituting this in \cref{eq:gradOfNLL:expansion} gives the desired expression for $\grad \negLL(v,T)$.

        To compute $\grad^2 \negLL_{P}(\mu)$,
        we differentiate \Cref{eq:gradOfNLL:compressed}.
        \begin{align*}
            \grad_{\mu}^2 \negLL_P(\mu)
            &= \frac{(\int_{\R^d} \grad^2 f)(\int_{\R^d} f) - (\int_{\R^d} \grad f)(\int_{\R^d} \grad f)^\top}{(\int_{\R^d} f)^2}
            - \frac{(\int_{P} \grad^2 f)(\int_{P} f) - (\int_{P} \grad f)(\int_{P} \grad f)^\top}{(\int_{P} f)^2} \\
            &= \frac{\int_{\R^d} \grad^2 f}{\int_{\R^d} f} - \frac{(\int_{\R^d} \grad f)(\int_{\R^d} \grad f)^\top}{(\int_{\R^d} f)^2}
            \qquad\quad~~~ - \left[ \frac{\int_P \grad^2 f}{\int_P f} - \frac{(\int_P \grad f)(\int_P \grad f)^\top}{(\int_P f)^2} \right]\,.
        \end{align*}
        This simplifies to the following 
        \begin{align*}
            \grad^2 \negLL_{P}(\mu) 
            &= 
            \Ex_{\normal{\mu}{I}}\insquare{
                    x
                    x^\top
            }
            - \Ex_{\normal{\mu}{I}}\insquare{
                    x
            }
            \Ex_{\normal{\mu}{I}}\insquare{
                    x
            }^\top
            - 
            \Ex_{\truncatedNormal{P}{\mu}{I}}\insquare{
                    x
                    x^\top
            }
            + \Ex_{\truncatedNormal{P}{\mu}{I}}\insquare{
                    x
            }
            \Ex_{\truncatedNormal{P}{\mu}{I}}\insquare{
                    x
            }^\top 
            \,.
        \end{align*}
        A final simplification gives 
        \[
            \grad^2 \negLL_{P}(\mu) 
            =
            \cov_{\normal{\mu}{I}}\insquare{
                    x
            }
            -
            \cov_{\truncatedNormal{P}{\mu}{I}}\insquare{
                    x
            }
        \]
        Substituting this in \cref{eq:gradOfNLL:expansion} gives the desired expression of $\grad^2 \negLL(\mu)$.
    \end{proof}

    
\section{Proofs Deferred From \texorpdfstring{\cref{sec:SS}}{Section 5}}
    In this section, we prove certain technical results whose proofs were deferred from \cref{sec:SS}.
        \subsection{Proof of \cref{lem:SS:probability} (Constant Probability Guarantee)}
            \label{sec:proofof:lem:SS:probability}
            In this section, we prove \cref{lem:SS:probability}, which we restate below.
            Recall that for any vector $u$, we use $\wh{u}$ to denote the unit vector parallel to $u$.
            \lemSSprobability*
            \begin{proof}[Proof of \cref{lem:SS:probability}]
                Recall that, for some $\gamma\in (0,\nfrac{1}{2}]$ and $R\geq 2$, $\evE=\evE_{i,\gamma,R}$ is the following event:
                \[
                    1 ~\leq~ 
                    \frac{\inangle{x, \wh{w}_i}}{
                        R\sqrt{\log{\nfrac{1}{\gamma}}}
                    }\,,
                    ~
                    \frac{\inangle{x, \wh{u}}}{
                        2\sqrt{\log{\nfrac{1}{\gamma}}}
                    }
                    ~\leq~ 2\,.
                \]
                Here $u\coloneqq v_i - \inangle{v_i,\wh{w}_i}\wh{w}_i$ is the component of $v_i$ orthogonal to $w_i$ {and $\hat u$ the unit vector parallel to $u$}.
                In the special case, where $v_i$ is parallel to $w_i$ (and, hence, $u=0$), the definition of $\evE$ omits the bound on $\inangle{x, \wh{u}}.$
                
                In this section, we prove that 
                \[
                    \Pr[\evE] = \gamma^{6+4R^2+\log{R}}\,.
                \]
                We will use the following standard facts.
                \begin{fact}[E.g., Proposition 2.1.2 in \citet{vershynin2018high}]\label{fact:GaussianTail}
                    The following tail bounds hold for all $z>0$
                    \[
                        \frac{e^{-z^2/2}}{\sqrt{2\pi}}\inparen{\frac{1}{z} - \frac{1}{z^3}}
                        \leq 
                        \Pr_{x\sim \cN(0,1)}[x\geq z]
                        \leq \frac{e^{-z^2/2}}{\sqrt{2\pi}}\frac{1}{z}\,.
                    \]
                \end{fact}
                We divide the proof into two parts.
                First, we prove the statement when $v_i$ is not parallel to $w_i$.
                Then, we extend the proof to the special case where $v_i$ is parallel to $w_i$.

                \paragraph{Case A ($v_i$ is not parallel to $w_i$).}
                    Since $\wh{u}$ is a projection of $v_i$ in a space orthogonal to $w_i$, $\wh{u}$ is orthogonal to $w_i$.
                    Further since $x\sim \cN(0,I)$ and $\wh w_i, \wh u$ are orthonormal, 
                    $\sinangle{x, \wh{w}_i}$ and $\sinangle{x, \wh{u}}$ are independent random variables with distributions $\sinangle{x, \wh{w}_i}\sim \cN(0,1)$ and $\sinangle{x, \wh{u}}\sim \cN(0,1)$.
                    Now, \cref{fact:GaussianTail} implies that 
                    \begin{align*}
                        \Pr\insquare{
                            \frac{
                                \sinangle{x, \wh{w}_i}
                            }{
                                R\sqrt{\log{\nfrac{1}{\gamma}}}
                            }
                            >
                            2
                        }
                        &\leq \frac{\gamma^{4R^2}}{\sqrt{2\pi}} \cdot 
                        \frac{1}{2R\sqrt{\log{\nfrac{1}{\gamma}}}}\,,\\
                        \Pr\insquare{
                            \frac{
                                \sinangle{x, \wh{w}_i}
                            }{
                                R\sqrt{\log{\nfrac{1}{\gamma}}}
                            } \geq 1
                        }
                        &\geq 
                        \frac{\gamma^{4R^2}}{\sqrt{2\pi}} \cdot 
                        \frac{1}{R\sqrt{\log{\nfrac{1}{\gamma}}}} 
                        \inparen{
                            1 - \frac{1}{\inparen{
                                R\sqrt{\log{\nfrac{1}{\gamma}}}
                            }^2}
                        }
                        \quad~~ \Stackrel{R\geq 2\,,~\gamma\leq \nfrac{1}{2}}{\geq}\quad~~ 
                        \frac{\gamma^{4R^2}}{\sqrt{2\pi}} \cdot 
                        \frac{6}{10R\sqrt{\log{\nfrac{1}{\gamma}}}} 
                        \,.
                    \end{align*} 
                    Hence, 
                    \[
                        \Pr\insquare{
                            \frac{\sinangle{x, \wh{w}_i}}{
                                R\sqrt{\log{\nfrac{1}{\gamma}}}
                            } \in [1,2]
                        }
                        \geq \frac{\gamma^{4R^2}}{\sqrt{2\pi}} \cdot \frac{1}{10R\sqrt{\log{\nfrac{1}{\gamma}}}}
                        \,.
                        \yesnum\label{eq:SS:event:lb1}
                    \]
                    Analogously, \cref{fact:GaussianTail} implies that 
                    \[
                        \Pr\insquare{
                            \frac{\sinangle{x, \wh{u}}}{
                                2\sqrt{\log{\nfrac{1}{\gamma}}}
                            } \in [1,2]
                        }
                        \geq \frac{\gamma^{16}}{\sqrt{2\pi}} \cdot \frac{1}{20\sqrt{\log{\nfrac{1}{\gamma}}}}
                        \,.
                        \yesnum\label{eq:SS:event:lb2}
                    \]
                    Since $\sinangle{x, \wh{w}_i}$ and $\sinangle{x, \wh{u}}$ are independent random variables, it holds that 
                    \[
                        \Pr[\evE]
                        = 
                        \Pr\insquare{
                            1 ~\leq~ 
                            \frac{\inangle{x, \wh{w}_i}}{
                                R\sqrt{\log{\nfrac{1}{\gamma}}}
                            }\,,
                            ~
                            \frac{\inangle{x, \wh{u}}}{
                                2\sqrt{\log{\nfrac{1}{\gamma}}}
                            }
                            ~\leq~ 2
                        }
                        ~~\quad\Stackrel{
                            \eqref{eq:SS:event:lb1}\,,~
                            \eqref{eq:SS:event:lb2}
                        }{\geq}\quad~~
                        \frac{\gamma^{16+4R^2}}{2\pi}\cdot \frac{1}{200R\log{\nfrac{1}{\gamma}}}\,.
                    \]
                    Since $\gamma \leq \nfrac{1}{2}$, $2^{11}\geq {400\pi}$, and $z\leq \inparen{{\log{\nfrac{1}{z}}}}^{-1}$ for any $z\in (0,1)$ it follows that 
                    \[
                        \Pr[\evE]
                        \geq 
                        \gamma^{28+4R^2+\log{R}}
                        ~~~\Stackrel{(R\geq 1)}{\geq}~~~
                        \gamma^{28+5R^2}\,.
                    \]
                    
                \paragraph{Case B ($v_i$ is parallel to $w_i$).}
                    In this case, $\evE$ is defined to be the following event
                    \[
                        1 ~\leq~ 
                        \frac{\inangle{x, \wh{w}_i}}{
                            R\sqrt{\log{\nfrac{1}{\gamma}}}
                        }
                        ~\leq~ 2\,.
                    \]
                    Calculations from the previous case imply that 
                    \[
                        \Pr[\evE]
                        \geq 
                        \frac{\gamma^{4R^2}}{\sqrt{2\pi}}\cdot \frac{1}{10R\sqrt{\log{\nfrac{1}{\gamma}}}}\,.
                    \]
                    Since $\gamma \leq \nfrac{1}{2}$, $2^5\geq \sqrt{2\pi}\cdot 10$, and $z\leq \inparen{\sqrt{\log{\nfrac{1}{z}}}}^{-1}$ for any $z\in (0,1)$ it follows that 
                    \[
                        \Pr[\evE]
                        \geq 
                        \gamma^{6+4R^2+\log{R}}
                        ~~~\Stackrel{(R\geq 1)}{\geq}~~~
                        \gamma^{6(1+R^2)}\,.
                    \]
                    
            \end{proof}

        
        \subsection{Proof of \cref{lem:SS:propertiesOfRho} (Properties of $\rho$)}
            \label{sec:proofof:lem:SS:propertiesOfRho}
            
            In this section, we prove \cref{lem:SS:propertiesOfRho}, which we restate below.
            \lemSSpropertiesOfRho*
            \begin{proof}[Proof of \cref{lem:SS:propertiesOfRho}]    
                Suppose $v_i$ is not parallel to $w_i$; we handle the edge case where $v_i$ is parallel to $w_i$ separately.
                Let $z_i$ be the component of $v_i$ orthogonal to $w_i$.
                Recall that for each $1\leq j\leq k$
                \begin{align*}
                    \rho_{j,V}
                    &= \proj_{{\rm span}(v_i,w_i)}(x)^\top v_j 
                    = 
                    \inangle{x,\wh{z}_i} \sinangle{v_j,\wh{z}_i}
                    +
                    \inangle{x,\wh{w}_i} \sinangle{v_j,\wh{w}_i}\,,\\
                    \rho_{j,W}
                    &= \proj_{{\rm span}(v_i,w_i)}(x)^\top w_j 
                    = 
                    \inangle{x,\wh{z}_i} \sinangle{w_j,\wh{z}_i}
                    +
                    \inangle{x,\wh{w}_i} \sinangle{w_j,\wh{w}_i}\,.
                \end{align*}
            We divide the proof into three parts corresponding to the three claims.
            
            \paragraph{Proof of Claim 1.}
                Observe that
                \begin{align*}
                    \rho_{i,W} - \rho_{j,W}
                    &=
                    \inangle{x,\wh{w}_i}\sinangle{w_i - w_j, \wh{w}_i}
                    +
                    \inangle{x,\wh{z}_i}\sinangle{w_i - w_j, \wh{z}_i}\,.
                \end{align*}
                First, we lower-bound the first term conditioned on $\evE$, 
                \begin{align*}
                    \inangle{x,\wh{w}_i}\sinangle{w_i - w_j, \wh{w}_i}
                    = \inangle{x,\wh{w}_i}
                    \cdot 
                    \frac{\norm{w_i}^2 - \sinangle{w_j,w_i}}{
                        {\norm{w_i}_2}
                    }
                    \quad~~~\Stackrel{\rm \cref{asmp:SS:separationBoundedness}}{\geq}\quad~~~ \inangle{x,\wh{w}_i}\cdot c 
                    ~~\quad\Stackrel{\rm\cref{def:SS:event}}{\geq}\quad~~ cR\sqrt{\log{\nfrac{1}{\gamma}}}\,.
                \end{align*}
                Next, we upper-bound the absolute value of the second term conditioned on $\evE$, 
                \begin{align*}
                    \abs{\inangle{x,\wh{z}_i}\sinangle{w_i - w_j, \wh{z}_i}}
                    \leq \abs{\inangle{x,\wh{z}_i}}\norm{w_i - w_j}
                    \quad~~~\Stackrel{\rm\cref{asmp:SS:separationBoundedness}}{\leq}\quad~~~ \abs{\inangle{x,\wh{z}_i}}\cdot 2C  
                    \quad~~\Stackrel{\rm\cref{def:SS:event}}{\leq}\quad~~ 
                    8C \cdot \sqrt{\log{\nfrac{1}{\gamma}}}\,.
                \end{align*}
                Therefore, it follows that, conditioned on $\evE$,  
                \[
                    \rho_{i,W} - \rho_{j,W} 
                    \geq
                        4(cR-8C)\cdot\sqrt{\log{\nfrac{1}{\gamma}}}\,.
                \]
                Further, observe that 
                \begin{align*}
                    \inparen{\rho_{i,V} - \rho_{j,V}}
                    - \inparen{\rho_{i,W} - \rho_{j,W}}
                    &= 
                    \inangle{x,\wh{z}_i}
                    \inparen{
                        \sinangle{v_i - w_i, \wh{z}_i}
                        + \sinangle{w_j - v_j, \wh{z}_i}
                    }
                    + \inangle{x,\wh{w}_i}
                    \inparen{
                        \sinangle{v_i - w_i, \wh{w}_i}
                        + \sinangle{w_j- v_j, \wh{w}_i}
                    }\,.
                \end{align*}
                Since $\norm{V-W}_F\leq \sfrac{c^3}{(3C)}$, the Cauchy--Schwarz inequality implies that
                \[
                    \inparen{\rho_{i,V} - \rho_{j,V}}
                    - \inparen{\rho_{i,W} - \rho_{j,W}}
                    \leq \frac{4c^3}{3C}\,.
                \]
                Therefore, it follows that, conditioned on $\evE$
                \[
                    \rho_{i,V} - \rho_{j,V} 
                    \geq 4\inparen{cR-8C-\frac{c^3}{3C}}\sqrt{\log{\nfrac{1}{\gamma}}}
                    ~~\quad \Stackrel{c\leq 1,~C\geq 1}{\leq}\quad~~ 
                    4\inparen{cR-9C}\sqrt{\log{\nfrac{1}{\gamma}}}\,.
                \]
                Finally, if $v_i$ is parallel to $w_i$, the result follows as a special case of the above since in that case $\rho_{i, W} - \rho_{j, W}
                    =
                    \inangle{x,\wh{w}_i}\sinangle{w_i - w_j, \wh{w}_i}$.

            \paragraph{Proof of Claim 2.}
                Conditioned on $\evE$ it follows that 
                \begin{align*}
                    \abs{\rho_{j,V} - \rho_{j,W}}
                    = \abs{
                        \inangle{x,\wh{z}_i}\sinangle{v_j - w_j, \wh{z}_i}
                        + 
                        \inangle{x,\wh{w}_i}\sinangle{v_j - w_j, \wh{w}_i}
                    }\,.
                \end{align*}
                Therefore, conditioned on \cref{def:SS:event}, 
                \begin{align*}
                    \abs{\rho_{j,V} - \rho_{j,W}} 
                    &\leq 2\sqrt{\log{\nfrac{1}{\gamma}}}\cdot \abs{\sinangle{v_j - w_j, \wh{z}_i}}
                    + 
                    2R\sqrt{\log{\nfrac{1}{\gamma}}}\cdot \abs{\sinangle{v_j - w_j, \wh{w}_i}}\,.
                \end{align*}
                Next, Cauchy--Schwarz inequality implies that 
                \[
                    \abs{\rho_{j,V} - \rho_{j,W}}  \leq 2(1+R)\cdot \sqrt{\log{\nfrac{1}{\gamma}}}\cdot \norm{v_j - w_j}_2\,.
                    \yesnum\label{eq:SS:propertiesOfRHO:UB1}
                \]
                Finally, in the special case where $v_i$ is parallel to $w_i$, the result follows as $\rho_{j,V} - \rho_{j,W}
                    = {
                        \inangle{x,\wh{w}_i}\sinangle{v_j - w_j, \wh{w}_i}
                    }$ and one can verify that \cref{eq:SS:propertiesOfRHO:UB1} holds.

            \paragraph{Proof of Claim 3.}
                Observe that 
                \[
                    \abs{\rho_{i,V} - \rho_{i,W}}
                    = \abs{\proj_{{\rm span}(v_i,w_i)}(x)^\top (v_i -w_i)}
                    = 
                    \abs{\inangle{x,\wh{z}_i} \sinangle{v_i -w_i,\wh{z}_i}
                    +
                    \inangle{x,\wh{w}_i} \sinangle{v_i -w_i,\wh{w}_i}}\,.
                \]
                Since $\wh{z}_i$ is orthogonal to $\wh{w}_i$ (by definition) and $v_i-w_i$ lies in ${\rm span}(z_i,w_i)$, it follows that 
                \[
                    \abs{\rho_{i,V} - \rho_{i,W}}
                    \geq 
                    \abs{
                        \min\inbrace{
                        \inangle{x,\wh{z}_i},
                        \inangle{x,\wh{w}_i}
                        } 
                        - 
                        \min\inbrace{
                            \inangle{x,\wh{z}_i},
                            \inangle{x,\wh{w}_i}
                        }
                    }\cdot \norm{v_i - w_i}\,.
                \]
                Therefore, conditioned on $\evE$,
                \[
                    \abs{\rho_{i,V} - \rho_{i,W}}
                    \geq (R-2)\sqrt{\log{\nfrac{1}{\gamma}}}\cdot \norm{v_i - w_i}
                    ~~~\Stackrel{(R\geq 12)}{\geq}~~~ \frac{5R\sqrt{\log{\nfrac{1}{\gamma}}}}{6}\cdot \norm{v_i - w_i}\,.
                \]
                Finally, in the special case where $v_i$ is parallel to $w_i$, the result holds because 
                \[
                    \abs{\rho_{i,V} - \rho_{i,W}}
                    = \abs{\proj_{{\rm span}(v_i,w_i)}(x)^\top (v_i -w_i)}
                    = x^\top \wh{w}_i \norm{v_i - w_i}_2
                    ~~\quad\Stackrel{\rm\cref{def:SS:event}}{\geq}\quad~~ 2R\sqrt{\log{\nfrac{1}{\gamma}}}\cdot \norm{v_i - w_i}_2\,.
                \]
                

                
            \end{proof}
            
        \subsection{Proof of \cref{lem:SS:cdfratios} (Separating Ratios of CDFs)}
            \label{sec:proofof:lem:SS:cdfratios}
            In this section, we prove \cref{lem:SS:cdfratios}, which restate below.
            \lemSScdfratios*
            \noindent In the proof of \cref{lem:SS:cdfratios}, we use the following two technical lemmas, which are proved at the end of this section.
            Roughly speaking, the first lemma lower bounds the sensitivity of $\Phi(\mu;\sigma^2)$ to deviations of $\mu$ when $\abs{\mu}$ is close to 0.
            The second lemma upper bounds the sensitivity of $\Phi(\mu,\sigma^2)$ to small deviations in both $\mu$ and $\sigma^2$.
            \begin{lemma}[Lower Bound on Sensivity to Mean]\label{lem:SS:cdfRatiosLB}
                For each $\mu\neq 0$ and $\nu^2,\sigma^2 > 0$, it holds that:
                \[
                    \frac{
                        \Phi\inparen{\mu; \sigma^2}
                    }{
                        \Phi\inparen{0; \nu^2}
                    }\,,~ 
                    \frac{
                        \Phi\inparen{0; \nu^2}
                    }{
                        \Phi\inparen{\mu; \sigma^2}
                    }
                    ~~\notin~~ 
                    1\pm \frac{1}{\sqrt{2\pi e}}
                    \min\inbrace{1, \frac{\abs{\mu}}{\sigma}}
                    \,.  %
                \]
            \end{lemma}
            \vspace{-5mm}
            \begin{lemma}[Upper Bound on Sensivity to Mean and Variance]\label{lem:SS:cdfRatiosUB}
                Fix any $\alpha\geq 0$ and $\beta^2\in(0,1)$.
                For any $\lambda_1 \geq \lambda_2\geq 0$ and $\sigma_1^2,\sigma_2^2\geq 1$ such that 
                \[
                    \lambda_1 - \lambda_2 \leq \alpha
                    \qquadand
                    \frac{\sigma_1^2}{\sigma_2^2}\,,~\frac{\sigma_1^2}{\sigma_1^2} \in 1\pm \beta^2
                    \,,
                    \yesnum\label{eq:SS:cdfRatiosUB:assumption}
                \]
                it holds that 
                \[
                    \frac{
                        \Phi\inparen{\lambda_1; \sigma_1^2}
                    }{
                        \Phi\inparen{\lambda_2; \sigma_2^2}
                    }\,,~ 
                    \frac{
                        \Phi\inparen{\lambda_2; \sigma_2^2}
                    }{
                        \Phi\inparen{\lambda_1; \sigma_1^2}
                    }
                    ~~\in~~ 
                    1\pm \beta^2 
                    \pm
                        \sqrt{\frac{8}{\pi}}\frac{\alpha+\beta^2\lambda_1}{\min\inbrace{\sigma_1,\sigma_2}}
                        e^{-\frac{\lambda_2^2(1-\beta^2)}{2\max\sinbrace{\sigma_1^2,\sigma_2^2}}}
                        \,.
                \]
            \end{lemma}
            We divide the proof of \cref{lem:SS:cdfratios} into two parts corresponding to \cref{eq:SS:cdfratios:lb,eq:SS:cdfratios:ub}.
            
            \paragraph{Proof of \cref{eq:SS:cdfratios:lb}.}
            Recall that after the translation described earlier $\min\inbrace{\rho_{i,V}, \rho_{i,W}}=0$.
            Without loss of generality, let $\rho_{i, W}=0$.
            By \cref{lem:SS:propertiesOfRho} we have that 
            \[
                \abs{\rho_{i, V}} = \abs{\rho_{i, W} - \rho_{i, V}}\geq \frac{5R}{6} \cdot \sqrt{\log{\nfrac{1}{\gamma}}}\cdot \norm{v_i - w_i}_2\,.
            \]
            Substituting $\mu=-\rho_{i,V}$, $\sigma^2=\sigma_{i,V}^2$, and $\nu^2=\sigma_{i,W}^2$ in \cref{lem:SS:cdfRatiosLB} implies that 
            \[
                \frac{
                    \customcdf{-\rho_{i,V}}{i,V}
                }{
                    \customcdf{-\rho_{i,W}}{i,W}
                }
                \notin 1 \pm \frac{1}{\sqrt{2\pi e}}
                    \min\inbrace{
                        1\,, 
                        \frac{5R}{6\sigma_{i,V}} \cdot \sqrt{\log{\nfrac{1}{\gamma}}}\cdot \norm{v_i - w_i}_2
                    }
                \,.
            \]
            Substituting $\sigma_{i,V}^2=1+C^2+(16c^4/C^2)\leq 17+C^2\leq 18C^2$ (from \cref{eq:SS:distributionOfZeta}), it follows that 
            \[
                \frac{
                    \customcdf{-\rho_{i,V}}{i,V}
                }{
                    \customcdf{-\rho_{i,W}}{i,W}
                }
                \notin 1 \pm \frac{1}{6C\sqrt{\pi e}}\cdot 
                    \min\inbrace{
                        1\,,
                        \frac{5R}{6}\cdot \sqrt{\log{\nfrac{1}{\gamma}}}\cdot \norm{v_i-w_i}_2 
                    }
                \,.
            \]

            \paragraph{Proof of \cref{eq:SS:cdfratios:ub}.}
                \cref{lem:SS:propertiesOfRho} implies that, conditioned on $\evE$, 
                \[
                    -\rho_{j,V}\,,~-\rho_{j,W}
                    ~~\geq~~  c R \sqrt{\log{\nfrac{1}{\gamma}}}\,.
                    \qquadand
                    \abs{\rho_{i,V} - \rho_{i,W}} \leq 3R\cdot \sqrt{\log{\nfrac{1}{\gamma}}}\cdot \norm{v_j-w_j}_2\,.
                \]
                Further, we also have the following upper bound conditioned on $\evE$,
                \[
                    \abs{\rho_{j,V}}\,,~\abs{\rho_{j,W}}
                    \leq 2cR\sqrt{\log{\nfrac{1}{\gamma}}}\,.
                \]
                A useful property of $\sigma_{j,V}^2$ and $\sigma_{j,W}^2$ is that they are close to each other when $v_j$ and $w_j$ are close:
                \[
                    \abs{\sigma_{j,V}^2 - \sigma_{j,W}^2 }
                    \leq \norm{\proj_{{\rm span}(v_i,w_i)^\bot}(v_j) - \proj_{{\rm span}(v_i,w_i)^\bot}(w_j)}_2^2
                    \leq \norm{v_j - w_j}_2^2\,.
                    \yesnum\label{eq:SS:differenceOfVariance}
                \]
                Where we used the triangle inequality and the fact that projection to a linear space is contractive.
                This combined with  \cref{eq:SS:distributionOfZeta} shows that 
                \[
                    1\leq \sigma_{j,V}^2\,, \sigma_{j,W}^2\leq 18C^2
                    \qquadand
                    \abs{\sigma_{j,V}^2 -  \sigma_{j,W}^2}
                    \leq \norm{v_j-w_j}_2^2\,.
                \]
                Hence,
                \[
                    \frac{\sigma_{j,V}^2}{\sigma_{j,W}^2}\,,~
                    \frac{\sigma_{j,W}^2}{\sigma_{j,V}^2}
                    \in 1\pm \norm{v_j-w_j}_2^2\,.
                \]
                The above along with \cref{lem:SS:cdfRatiosUB}  implies that 
                \[
                    \frac{\customcdf{-\rho_{j,V}}{j,V}}{\customcdf{-\rho_{j,W}}{j,W}}
                        ~~\in~~
                        1\pm \norm{v_j-w_j}_2^2
                        \pm R\sqrt{\frac{8}{\pi} \log{\nfrac{1}{\gamma}}}\cdot 
                        \frac{3\norm{v_j-w_j}_2+2C\norm{v_j-w_j}_2^2}{1} \cdot e^{-\frac{c^2R^2{\log{\nfrac{1}{\gamma}}}(1-\snorm{v_j-w_j}_2^2)}{36C^2}}\,.
                \]
                On using that $\snorm{v_j-w_j}_2\leq \nfrac{1}{2}$, $c,\gamma\leq 1,$ $R,C\geq 1$, $\sqrt{\nfrac{8}{\pi}}\leq 2$, and simplifying, we obtain that 
                \[
                    \frac{\customcdf{-\rho_{j,V}}{j,V}}{\customcdf{-\rho_{j,W}}{j,W}}
                        ~~\in~~
                        1
                        \pm 5R\sqrt{\log{\nfrac{1}{\gamma}}}\cdot 
                        \inparen{\norm{v_j-w_j}_2+C\norm{v_j-w_j}_2^2}
                        \cdot \gamma^{-\frac{c^2R^2}{72C^2}}\,.
                 \]
            This completes the proof of \cref{lem:SS:cdfratios} up to proving \cref{lem:SS:cdfRatiosLB,lem:SS:cdfRatiosUB}.
            In the remainder of this section, we prove \cref{lem:SS:cdfRatiosLB,lem:SS:cdfRatiosUB}.
            
            \begin{proof}[Proof of \cref{lem:SS:cdfRatiosLB}]
                By definition,
                \begin{align*}
                    \Phi(\mu;\sigma^2) 
                    &= \frac{1}{2} + \frac{\sgn(\mu)}{\sqrt{2\pi}\sigma} \int_{0}^{\abs{\mu}} e^{-z^2/(2\sigma^2)}\d z
                    \geq \frac{1}{2} + \frac{
                        \max_{0\leq z\leq \abs{\mu}} ze^{-z^2/(2\sigma^2)}
                    }{\sqrt{2\pi}\sigma}\,.
                \end{align*}
                Since $\Phi(0;\nu^2)=\nfrac{1}{2}$, 
                \[
                    \frac{\Phi(\mu;\sigma^2)}{\Phi(0;\nu^2)}
                    \geq 1 + \sqrt{\frac{2}{\pi}} \frac{
                        \max_{0\leq z\leq \abs{\mu}} ze^{-z^2/(2\sigma^2)}
                    }{\sigma}
                    \,.
                \]
                Finally, as  
                \[
                    \sqrt{\frac{2}{\pi}}\frac{
                        \max_{0\leq z\leq \abs{\mu}} ze^{-z^2/(2\sigma^2)}
                    }{\sigma}
                    < \frac{1}{2}\,,
                \]
                and $\sfrac{1}{(1+z)} < 1-(\nfrac{z}{2})$ for all $0 < z < 1$, it follows that 
                \[
                    \frac{\Phi(0;\nu^2)}{\Phi(\mu;\sigma^2)}
                    > 1 - \frac{1}{\sqrt{2\pi}}\cdot \frac{\max_{0\leq z\leq \abs{\mu}} ze^{-z^2/(2\sigma^2)}}{\sigma}\,.
                \]
                The result follows from the following inequality (which can be checked by substitution)
                \[
                    \frac{
                        \max_{0\leq z\leq \abs{\mu}} ze^{-z^2/(2\sigma^2)}
                    }{
                        \sigma 
                    }
                    \geq 
                    \begin{cases}
                        \frac{1}{\sigma\sqrt{e}}\cdot\abs{\mu} & \text{if } {\mu^2}\leq \sigma^2\\
                        \frac{1}{\sqrt{e}} & \text{otherwise}
                    \end{cases}
                    \geq \frac{1}{\sqrt{e}} \min\inbrace{1, \frac{\abs{\mu}}{\sigma}}\,.
                \]
            \end{proof}
            
            \begin{proof}[Proof of \cref{lem:SS:cdfRatiosUB}]
                Observe that 
                \[
                    \frac{\Phi(\lambda_1;\sigma_1^2)}{\Phi(\lambda_2;\sigma_2^2)}
                    = \frac{
                        \inparen{\nfrac{1}{\sigma_1}}\cdot \int_{-\infty}^{\lambda_1} e^{-\sfrac{z^2}{2\sigma_1^2}} \d z
                    }{
                        \inparen{\nfrac{1}{\sigma_2}}\cdot \int_{-\infty}^{\lambda_2} e^{-\sfrac{z^2}{2\sigma_2^2}} \d z
                    }
                    = \frac{\sigma_2}{\sigma_1}\cdot \frac{
                        \inparen{\nfrac{1}{\sigma_2}}\cdot \int_{-\infty}^{\lambda_1\sigma_2/\sigma_1} e^{-\sfrac{z^2}{2\sigma_2^2}} \d z
                    }{
                        \inparen{\nfrac{1}{\sigma_2}}\cdot \int_{-\infty}^{\lambda_2} e^{-\sfrac{z^2}{2\sigma_2^2}} \d z
                    }
                    = \frac{\sigma_2}{\sigma_1}\cdot \frac{\Phi(\lambda_1\sigma_2/\sigma_1;\sigma_2^2)}{\Phi(\lambda_2;\sigma_2^2)}\,.
                    \yesnum\label{eq:SS:cdfRatiosUB:equality}
                \]
                Toward bounding this expression, we need to bound $\sfrac{\sigma_2}{\sigma_1}$ and $(\lambda_1\sigma_2/\sigma_1) - \lambda_2$:
                First, by our assumptions, $\sfrac{\sigma_2^2}{\sigma_1^2}\in 1\pm \beta^2$, which implies that 
                \[
                    \frac{\sigma_2}{\sigma_1}\in 1\pm \beta^2\,,
                    \yesnum\label{eq:SS:cdfRatiosUB:varRatio}
                \]
                since $\sqrt{1\pm z^2}\in 1\pm z^2$ for any $z$ with $\abs{z}<1$.
                Second 
                \[
                    \lambda_1\cdot \frac{\sigma_2}{\sigma_1}- \lambda_2
                    ~~\Stackrel{\eqref{eq:SS:cdfRatiosUB:varRatio}}{\in}~~ \lambda_1 - \lambda_2 \pm \beta^2 \lambda_1 
                    ~~\Stackrel{\eqref{eq:SS:cdfRatiosUB:assumption}}{\in}~~ [-\beta^2 \lambda_1 -\alpha,\alpha+\beta^2 \lambda_1 ]\,.
                \]
                Further, it holds that 
                \begin{align*}
                    \abs{\Phi(\lambda_1\sigma_2/\sigma_1;\sigma_2^2) 
                    - \Phi\inparen{\lambda_2; \sigma_2^2}}
                    &= \frac{1}{\sqrt{2\pi}\sigma_2}\int_{\lambda_2}^{\sfrac{\lambda_1\sigma_2}{\sigma_1}} e^{-z^2/(2\sigma_2^2)}\d z
                    \leq 
                    \frac{\alpha+\beta^2\lambda_1 }{\sqrt{2\pi}\sigma_2}
                     e^{-\lambda_2^2(1-\beta^2)/(2\sigma_2^2)}\,. %
                \end{align*}
                Since $\Phi(0;\sigma_2^2)=\nfrac{1}{2}$,
                $\lambda_2\geq 0$, and $\Phi(\cdot)$ is monotone and increasing, $\Phi(\lambda_2;\sigma_2^2)\geq \nfrac{1}{2}$.
                Therefore, dividing by $\Phi(\lambda_2;\sigma_2^2)$ implies that 
                \[
                    1 
                    - \sqrt{\frac{2}{\pi}}\frac{\alpha+\beta^2\lambda_1}{\sigma_2}
                     e^{-\lambda_2^2(1-\beta^2)/(2\sigma_2^2)}
                    \leq 
                    \frac{\Phi(\lambda_1\sigma_2/\sigma_1;\sigma_2^2) }{\Phi\inparen{\lambda_2; \sigma_2^2}}
                    \leq 1 
                    + \sqrt{\frac{2}{\pi}}\frac{\alpha+\beta^2\lambda_1}{\sigma_2}
                     e^{-\lambda_2^2(1-\beta^2)/(2\sigma_2^2)}\,.
                \]
                Combining this with \cref{eq:SS:cdfRatiosUB:equality,eq:SS:cdfRatiosUB:varRatio} and using that $1+\beta^2\leq 2$ implies that 
                \[
                    1 
                    -\beta^2 - \sqrt{\frac{2}{\pi}}\frac{\alpha+\beta^2\lambda_1}{\sigma_2}
                     e^{-\lambda_2^2(1-\beta^2)/(2\sigma_2^2)}
                    \leq 
                    \frac{\Phi(\lambda_1;\sigma_1^2)}{\Phi(\lambda_2;\sigma_2^2)}
                    \leq 1 
                    + \beta^2 + \sqrt{\frac{8}{\pi}}\frac{\alpha+\beta^2\lambda_1}{\sigma_2}
                     e^{-\lambda_2^2(1-\beta^2)/(2\sigma_2^2)}\,.
                \]
                This implies part of the result.
                It remains to bound $\sfrac{\Phi(\lambda_2;\sigma_2^2)}{\Phi(\lambda_1;\sigma_1^2)}$, which follows by replacing $\lambda_1,\lambda_2,\sigma_1,\sigma_2$ by $\lambda_2,\lambda_1,\sigma_2,\sigma_1$ respectively in the above proof.
            \end{proof}




    

    
    \subsection{Proof of \cref{fact:SS:LC:boundOnTruncatedMoments} (Upper Bounds on Moments of Truncated Gaussian with One-Sided Truncation)}
        \label{sec:proofof:fact:SS:LC:boundOnTruncatedMoments}
        In this section, we prove \cref{fact:SS:LC:boundOnTruncatedMoments}, which we restate below.
        \factboundOnTruncatedMoments*
        \begin{proof}[Proof of \cref{fact:SS:LC:boundOnTruncatedMoments}]
            \citet{orjebin2014recursive} provide the following formulas for the second and fourth moments of the normal distribution truncated to $(-\infty,b]$:
            \begin{align*}
                \Ex_{z\sim \normal{\mu}{1}}\insquare{z^2\mid z\leq b}
                    &= \mu^2 + 1 - \frac{(\mu+b)\,\phi(b-\mu)}{\Phi(b-\mu)}\,,
                    \yesnum\label{eq:SS:LC:momentBound2}\\
                \Ex_{z\sim \normal{\mu}{1}}\insquare{z^{\smash{4}}\mid z\leq b}
                    &=\mu^4 + 6\mu^2 + 3 - \frac{\bigl(b^3+b^2\mu+b\mu^2+3b+5\mu+\mu^3\bigr)\,\phi(b-\mu)}{\Phi(b-\mu)}\,.
                    \yesnum\label{eq:SS:LC:momentBound4}
            \end{align*} 
            Where $\phi(\cdot)$ and $\Phi(\cdot)$ are the probability density function and the cumulative density function of the standard normal distribution.
            
             
            Since $\nfrac{\phi(z)}{\Phi(z)}\leq \abs{z}+1$ for $z\in \R$ (see \cref{fact:BoundOnMillsLikeRatio}),  \cref{eq:SS:LC:momentBound2} implies,
            \[
                \Ex_{z\sim \normal{\mu}{1}}\insquare{z^{{2}}\mid z\leq b} 
                \le \mu^2+1 + \inparen{\abs{b}+\abs{\mu}}\,\inparen{1+\abs{b}+\abs{\mu}}
                \leq 2\inparen{1+\abs{b}+\abs{\mu}}^2 
                \,.
            \] 
            It remains to upper bound $\Ex_{z\sim \normal{\mu}{1}}\insquare{z^{{4}}\mid z\leq b}$.
            To simplify the notation, define,
            \[
                M \coloneqq |b|+|\mu|\,.
            \]
            Since $\left|b^3+b^2\mu+b\mu^2+\mu^3\right|\le M^3$, $|3b+5\mu|\le 5M$, and $\nfrac{\phi(z)}{\Phi(z)}\leq \abs{z}+1$ for $z\in \R$ (see \cref{fact:BoundOnMillsLikeRatio}),  \cref{eq:SS:LC:momentBound4} implies,
            \[
                \Ex_{z\sim \normal{\mu}{1}}\insquare{z^{{4}}\mid z\leq b}
                \le 
                    \mu^4+6\mu^2+3 
                    +
                    \inparen{
                        M^3
                        +
                        5 M
                    } 
                    \inparen{1 + M}
                \,.
            \]
            Using $\mu^4+6\mu^2+3 \leq M^4 + 9 M^2$ and $(z^3+5z)(z+1) = z^4 + z^3 + 5z^2 + 5z$ implies 
            \[
                \Ex_{z\sim \normal{\mu}{1}}\insquare{z^{{4}}\mid z\leq b}
                \leq 
                M^4 + 9M^2 + (M^3+5M)(M+1)
                \leq 2M^4 + M^3 + 14M^2 + 5M\,.
            \]
            This implies the result since $M\geq 0$, $2z^4 + z^3 + 14z^2 + 5z\leq 3\inparen{1+z}^4$ for any $z\geq 0$.
        \end{proof}

        \subsubsection*{Proof of Upper Bound on $\sfrac{\phi(\cdot)}{\Phi(\cdot)}$}
        \begin{fact}\label{fact:BoundOnMillsLikeRatio}
            For $z\in \R$,
            \[
                \frac{\phi(z)}{\Phi(z)} \le |z| + 1\,.
            \]
            Where $\phi(\cdot)$ and $\Phi(\cdot)$ are the probability density function and cumulative density function of the standard normal distribution.
        \end{fact}
        \begin{proof}
            We divide the proof into two cases.
            \begin{itemize}
                \item \textbf{Case A ($z\geq 0$):}~~
                    Since \(\Phi(z)\ge \Phi(0)=\nfrac{1}{2}\) for \(z\ge 0\), $\nfrac{\phi(z)}{\Phi(z)} \le 2\phi(z)$.
                    Further, $\phi(z)=\inparen{\nfrac{1}{\sqrt{2\pi}}}e^{-z^2/2}$ is maximized at \(z=0\), so that for \(z\ge 0\), $2\phi(z) \le 2\phi(0)=\sqrt{\sfrac{2}{\pi}} \le 0.8\,.$
                    But we also have \(|z|+1=z+1\ge 1\) and, hence, $\nfrac{\phi(z)}{\Phi(z)} \le 0.8 \le z+1 = |z|+1$.
                \item \textbf{Case B ($z < 0$):}~~
                    Set \(w=-z>0\). By the symmetry of the standard normal density, we have
                    \[
                        \frac{\phi(z)}{\Phi(z)} = \frac{\phi(-w)}{\Phi(-w)} = \frac{\phi(w)}{1-\Phi(w)}\,.
                    \]
                    Define the function
                    \[
                        g(w) = (w+1)(1-\Phi(w)) - \phi(w)\,,\quad w>0\,.
                    \]
                    We claim that \(g(w)\ge 0\) for all \(w>0\): this holds because of the standard lower bound $1-\Phi(x)\geq \phi(x) \cdot \frac{2}{\sqrt{4+x^2}+x}$ for $x\geq 0$ (see  Equation (3) in \citet{duembgen2010boundingstandardgaussiantail}) and the fact that $z+1\geq \frac{1}{2}\sinparen{\sqrt{4+z^2}+z}$ for $z\geq 0$.
                    Thus, we conclude $(w+1)(1-\Phi(w)) \ge \phi(w)$ or equivalently,
                    \[
                    \frac{\phi(w)}{1-\Phi(w)} \le w+1\,.
                    \]
                    Returning to the variable \(z\) (with \(w=-z\)), we deduce
                    \[
                    \frac{\phi(z)}{\Phi(z)} \le |z| + 1\,.
                    \]
            \end{itemize}
        \end{proof}
         

    
    
\section{Log-Concave Sampling over Convex Bodies}\label{apx:sampling-polytopes}
% --- begin logconcave_sampling.tex ---


In this section, we review well-known results about sampling from a log-concave density $\propto e^{-f}$ constrained to a convex body $K$, under mild assumptions. 
We use these results to implement the polytope sampling oracle (\Cref{asmp:samplingOracle}) in \Cref{sec:coarseGaussianEstimation}.
{Concretely,} we have a convex function $f: K\to \R$ that is bounded below, say there is some $x^\star\in K$ such that {$f(x) \geq f(x^\star)$} for all $x\in K$.
{Note that $x^\star$ exists and belongs in $K$ since $K$ is closed.}






    One of the most general tools for this purpose is the ``Hit-And-Run'' Markov Chain Monte Carlo (MCMC) algorithm \cite{lovasz2006fast}.
    \begin{theorem}[Mixing-Time of Hit-And-Run Markov Chains; \cite{lovasz2006simulated,lovasz2006fast}]\label{thm:convex body logconcave sampling}
        Consider a logconcave distribution $\pi_f\propto e^{-f}$ over a convex body $K$.
        Suppose we are provided the following.
        \begin{enumerate}[label=(S\arabic*)]
            \item Zero-order access to $f$.
            \label[assumption]{assum:evaluation-oracle}
            \item Membership access to $K$.
            \label[assumption]{assum:membership-oracle}
            \item A point $x^{(0)}$ {and constants $R\geq r > 0$} such that $B_2(x^{(0)}, r)\sset K\sset B_2(x^0, R)$.
            \label[assumption]{assum:interior-point-oracle}
            \item A bound $M\geq \max_{x\in K} f(x) - f(x^\star)$.
            \label[assumption]{assum:range-oracle}
        \end{enumerate}
        Then, there is an algorithm that makes
        \[
            O\left( d^{4.5}\cdot \polylog(d, M, \nicefrac{R}r, \nicefrac1\delta) \right)
        \]
        membership oracle calls to produce a random vector within $\delta$-TV distance of $\pi_f$.
    \end{theorem}
    To implement the desired Markov chain,
    we need \Crefrange{assum:evaluation-oracle}{assum:range-oracle} stated in \Cref{thm:convex body logconcave sampling}.
    We remark that \Cref{assum:membership-oracle} and \Cref{assum:interior-point-oracle}
    can both be implemented given a separation oracle to $K$.
    If we also know that $f$ is $L$-Lipschitz or $\beta$-smooth,
    we can take $M = LR$ and $M=LR^2$ respectively for \Cref{assum:range-oracle}.
    The smallest ratio $\nfrac{R}r$ of values $r, R$ from \Cref{assum:interior-point-oracle} depends on the structure of $K$.
    For the simple case of 1-dimensional intervals,
    we trivially have $\nfrac{R}r = 1$.
    For all $d$-dimensional $\ell_p$-balls,
    we have $\nfrac{R}r = \poly(d)$.

    \subsection{Removing \texorpdfstring{\Cref{asmp:samplingOracle}}{Assumption}} \label{apx:sampling-polytopes:removing-oracle-assumption}
    The sampling oracle assumption (\Cref{asmp:samplingOracle}) from our coarse Gaussian mean estimation algorithm (\Cref{sec:coarseGaussianEstimation}) can be relaxed to \Cref{asmp:wellDescribedPolyhedron} below.
    Before stating the exact assumption,
    we state a natural notion of the complexity of a polytope.
    \begin{definition}[Facet-Complexity]
        We say that a polytope $P\sset \R^d$ has \emph{facet-complexity} at most $\varphi_P$
        if there exists a system of inequalities with rational coefficients that has solution set $P$ and such that the bit-encoding length of each inequality of the system is at most $\varphi_P$.
        In case $P=\R^d$,
        we require $\varphi_P\geq d+1$.
    \end{definition}
    We are now ready to replace \Cref{asmp:samplingOracle}.
    \begin{assumption}[Well-Described Polyhedron]\label{asmp:wellDescribedPolyhedron}
        Each observation $P\in \hyP$ is a polyhedron and is provided in the form of a separation oracle and an upper bound $\varphi_P > 0$
        on the facet complexity of $P$.
    \end{assumption}
    The running time of our sampling algorithm then depends polynomially on the running time of the separation oracle
    as well as the facet-complexity $\varphi_P$ of the observed samples $P$.
    Indeed, let us check that \Cref{asmp:wellDescribedPolyhedron} allows us to apply \Cref{thm:convex body logconcave sampling}.
    It is clear that \Cref{asmp:wellDescribedPolyhedron} immediately satisfies \Cref{assum:evaluation-oracle,assum:membership-oracle}.
    We sketch how to address \Cref{assum:interior-point-oracle,assum:range-oracle}.

    {First,
    note that we wish to sample from the truncated standard Gaussian $\truncatedNormal{\mu}{I}{P\cap B_\infty(0, R)}$.
    In particular,
    it has density that is $O(1)$-smooth.
    Since $R = \poly(d, D, \log(\nfrac1\delta))$ is a polynomial of the dimension $d$,
    the warm-start radius $D$,
    and logarithmic in the inverse failure probability $\delta$.
    we can take $M = {O(R^2)}$ so that \Cref{assum:range-oracle} is satisfied.

    Next,
    we can again use the fact that we sample from polytopes contained in $B_\infty(0, R)$
    to deduce that $P\cap B_\infty(0, R)$ is contained in an $\ell_2$ ball of radius $R\sqrt{d}$.
    Moreover,
    the facet-complexity of $P\cap B_\infty(0, R)$ is at most $\varphi = \varphi_P + \log_2(R)$.
    On the other hand,
    to handle the inner ball,
    we draw on \cite[Lemma 6.2.5]{grotschel1988geometric} which states that a full-dimensional polytope with facet-complexity $\varphi$ must contain a ball of radius $2^{-7d^3\varphi}$.
    In the case that $P\cap B_\infty(0, R)$ is full-dimensional,
    this suffices to run the Markov chain from \Cref{thm:convex body logconcave sampling} in $\poly(d, \varphi)$ time.
    If $P$ is not full-dimensional,
    we can exactly compute the affine hull in polynomial time using the ellipsoid algorithm~\cite[(6.1.2)]{grotschel1988geometric}
    and then apply the full-dimensional argument on this affine subspace.}
    
    
    











% --- end logconcave_sampling.tex ---





\end{document}
